\documentclass[11pt]{article}

\usepackage[T1]{fontenc}
\usepackage[utf8]{inputenc}
\usepackage[margin=1in]{geometry}
\usepackage{amsmath,amssymb,amsthm,mathtools}
\usepackage[colorlinks=true,linkcolor=blue,citecolor=blue,urlcolor=blue]{hyperref}

\numberwithin{equation}{section}
\newtheorem{theorem}{Theorem}[section]
\newtheorem{proposition}[theorem]{Proposition}
\newtheorem{lemma}[theorem]{Lemma}
\newtheorem{corollary}[theorem]{Corollary}
\newtheorem{conjecture}[theorem]{Conjecture}
\newtheorem{problem}[theorem]{Problem}
\theoremstyle{definition}
\newtheorem{definition}[theorem]{Definition}
\newtheorem{example}[theorem]{Example}
\newtheorem{counterexample}[theorem]{Counterexample}
\theoremstyle{remark}
\newtheorem{remark}[theorem]{Remark}

\DeclareMathOperator{\Var}{Var}
\DeclareMathOperator{\Prob}{P}
\DeclareMathOperator{\supp}{supp}
\newcommand{\E}{\mathbb E}
\newcommand{\N}{\mathbb N}
\newcommand{\Z}{\mathbb Z}
\newcommand{\1}{\mathbf 1}

\title{Zeta-Law Entropy, Gamma Curvature, and Laplace-Measure Applications}
\author{Pudim AI Project}
\date{Version v023 -- 2026-06-05}

\setcounter{tocdepth}{3}
\begin{document}
\maketitle

\begin{abstract}
This paper develops a unified public theory built around Laplace-transform
positivity.  Zeta-law entropy is organized through logarithmic compression,
gamma-ratio and polygamma inequalities are expressed by completely monotone
kernels, and applications are obtained by transporting measure positivity
through Stieltjes, Bernstein, stable-subordination, determinant, triangular,
inverse-Laplace, derivative-chain, ratio-kernel, and moment-ratio mechanisms.
The emphasis is on reusable transformations: source problems enter as fixed
applications, while the proofs expose the structural tools that connect them.
\end{abstract}

\tableofcontents

\medskip



\noindent\textbf{Keywords.} Riemann zeta function; zeta probability law; entropy; Laplace transform; complete monotonicity; Bernstein functions; Stieltjes functions; gamma function; digamma function; Mills ratio; generalized Stieltjes functions.

\section{Motivation and main results}

Many classical identities for \(\zeta(s)\) become more transparent after a
probabilistic normalization.  For \(\beta>1\), the weights \(n^{-\beta}\)
define a Gibbs law whose partition function is precisely \(\zeta(\beta)\).
The energy is \(\log n\), the free energy is \(\log\zeta(\beta)\), and ratios
such as \(\zeta(\beta+r)/\zeta(\beta)\) become moments.  This viewpoint is
elementary, but it packages several unrelated-looking zeta inequalities into a
common calculus.

The same principle guides the special-function part of the paper.  Complete
monotonicity, Bernstein properties, Stieltjes representations, and endpoint
obstructions are treated as transport questions: once a positive kernel,
measure representation, determinant reduction, or rational certificate is
identified, it can be reused across multiple source problems.  The central
theme is therefore not an individual inequality, but the mechanism that makes
families of inequalities stable under analytic transformations.

Section~\ref{sec:applications} records solved external problems with stable
application labels.  The ledger is not a chronology; it is a public index of
source questions and the results that resolve them.  The body of the paper
organizes those results by mechanism: zeta-law compression, gamma and
polygamma curvature, Laplace and Stieltjes transport, determinant and
moment-ratio reductions, endpoint obstructions, and explicit certificates.
This organization lets the reader trace each application back to the tools
that make it work.

\section{Applications}
\label{sec:applications}
The application ledger is append-only public memory.  Each row records the immutable APP identifier, the canonical short description, the result reference used by the paper, and the cited source problem.
\begin{enumerate}
\item \label{ledger:app-0001}\textbf{APP-0001.} Alzer-Kwong convexity and concavity problem. Solved by Theorem~\ref{thm:alzer-kwong}; source problem reference \cite{app-app-0001-source}.
\item \label{ledger:app-0002}\textbf{APP-0002.} Nantomah zeta positivity problem. Solved by Theorem~\ref{thm:nantomah-positivity}; source problem reference \cite{nantomah}.
\item \label{ledger:app-0003}\textbf{APP-0003.} Sroysang generalized Holder problem. Solved by Theorem~\ref{thm:sroysang-holder}; source problem reference \cite{sroysang}.
\item \label{ledger:app-0004}\textbf{APP-0004.} Complete monotonicity of \((\log\Gamma(x)+\log\Gamma(1/x))''\). Solved by Theorem~\ref{thm:reciprocal-gamma-complete-monotonicity}; source problem reference \cite{app-app-0004-source}.
\item \label{ledger:app-0005}\textbf{APP-0005.} Exact inverse-tail floor formula at s=7. Solved by Theorem~\ref{thm:tail-s7-floor}; source problem reference \cite{kim-song-tail}.
\item \label{ledger:app-0006}\textbf{APP-0006.} Exact inverse-tail floor formula at s=8. Solved by Theorem~\ref{thm:tail-s8-floor}; source problem reference \cite{kim-song-tail}.
\item \label{ledger:app-0007}\textbf{APP-0007.} Concavity or complete monotonicity of the polygamma product P0. Solved by Theorem~\ref{thm:reciprocal-digamma-product-complete-monotonicity}; source problem reference \cite{app-app-0004-source}.
\item \label{ledger:app-0008}\textbf{APP-0008.} Higher-order monotonicity of polygamma products Pn. Solved by Theorem~\ref{thm:pn-complete-monotonicity-counterexample}; source problem reference \cite{app-app-0004-source}.
\item \label{ledger:app-0009}\textbf{APP-0009.} Sharp reciprocal Gamma-product monotonicity threshold. Solved by Theorem~\ref{thm:gamma-product-sharp-threshold}; source problem reference \cite{app-app-0009-source}.
\item \label{ledger:app-0010}\textbf{APP-0010.} Nielsen \(k\)-beta derivative-ratio monotonicity. Solved by Theorem~\ref{thm:nielsen-k-beta-derivative-ratio}; source problem reference \cite{yin-zhang-k-beta}.
\item \label{ledger:app-0011}\textbf{APP-0011.} Exact \(n=2\) beta-window endpoint. Solved by Theorem~\ref{thm:q2-exact-endpoint}; source problem reference \cite{app-app-0004-source}.
\item \label{ledger:app-0012}\textbf{APP-0012.} Baricz \(V\_q\) strict \(q\)-log-convexity. Solved by Theorem~\ref{thm:baricz-vq-strict-logconvexity}; source problem reference \cite{baricz-turan-thesis}.
\item \label{ledger:app-0013}\textbf{APP-0013.} Yin \((p,k)\)-digamma sharp \(\alpha\)-necessity. Solved by Theorem~\ref{thm:yin-pk-alpha-necessity}; source problem reference \cite{yin-pk-digamma}.
\item \label{ledger:app-0014}\textbf{APP-0014.} Ramanujan integral is Stieltjes and has a complete Bernstein primitive. Solved by Theorem~\ref{thm:app14-ramanujan-stieltjes}; source problem reference \cite{mishra-swaminathan-ramanujan}.
\item \label{ledger:app-0015}\textbf{APP-0015.} Sibisi-Prabhakar measure is a stable-subordination push-forward. Solved by Theorem~\ref{thm:app15-prabhakar-subordination}; source problem reference \cite{sibisi-prabhakar}.
\item \label{ledger:app-0016}\textbf{APP-0016.} From's Mills-ratio bound extends to every derivative level. Solved by Theorem~\ref{thm:app16-mills-all-l}; source problem reference \cite{from-mills}.
\item \label{ledger:app-0017}\textbf{APP-0017.} Baricz gamma-quotient Bernstein-for-all problem has a rational counterexample. Solved by Theorem~\ref{thm:app17-baricz-counterexample}; source problem reference \cite{baricz-turan}.
\item \label{ledger:app-0018}\textbf{APP-0018.} Qi-Guo tau threshold has a supremum-corrected exact value. Solved by Theorem~\ref{thm:app18-qg-tau}; source problem reference \cite{qi-guo-2004}.
\item \label{ledger:app-0019}\textbf{APP-0019.} Wakrim W-symbol Bernstein range closes the beta gap. Solved by Theorem~\ref{thm:app19-w-symbol}; source problem reference \cite{wakrim-w}.
\item \label{ledger:app-0020}\textbf{APP-0020.} Qi's degree-four conjecture for h\_lambda is refuted at the endpoint. Solved by Theorem~\ref{thm:app20-qi-hlambda}; source problem reference \cite{qi-hlambda}.
\item \label{ledger:app-0021}\textbf{APP-0021.} Szabo's cutoff problem has \(\alpha\_0=1\). Solved by Theorem~\ref{thm:app21-szabo-cutoff}; source problem reference \cite{szabo-cm}.
\item \label{ledger:app-0022}\textbf{APP-0022.} Chen-Choi Gurland gamma-ratio conjecture is logarithmically completely monotone. Solved by Theorem~\ref{thm:app22-gurland-logcm}; source problem reference \cite{chen-choi-gurland}.
\item \label{ledger:app-0023}\textbf{APP-0023.} Sokal generalized-Stieltjes lambda-derivatives have a triangular nonnegative compression. Solved by Theorem~\ref{thm:app23-sokal-lambda}; source problem reference \cite{sokal-gs}.
\item \label{ledger:app-0024}\textbf{APP-0024.} Simon gamma quotient is Bernstein. Solved by Theorem~\ref{thm:app24-simon-gamma-quotient}; source problem reference \cite{simon-moment}.
\item \label{ledger:app-0025}\textbf{APP-0025.} Du-Wang \(h\_3\) monotonicity is classified. Solved by Theorem~\ref{thm:app25-du-wang-h3}; source problem reference \cite{du-wang-h3}.
\item \label{ledger:app-0026}\textbf{APP-0026.} Baskakov even-power complete-monotonicity conjecture is false. Solved by Theorem~\ref{thm:app26-baskakov-r8}; source problem reference \cite{abel-gawronski-neuschel}.
\item \label{ledger:app-0027}\textbf{APP-0027.} Ma-Weigert derivative regions form a descending chain. Solved by Theorem~\ref{thm:app27-ma-weigert-chain}; source problem reference \cite{ma-weigert-log-functions}.
\item \label{ledger:app-0028}\textbf{APP-0028.} Qi-Agarwal/Yin divisor-polygamma parity is corrected. Solved by Theorem~\ref{thm:app28-qa-divisor-polygamma}; source problem reference \cite{qi-agarwal-gamma-ratios}.
\item \label{ledger:app-0029}\textbf{APP-0029.} Bulboaca-Zayed Gamma quotient has one critical point. Solved by Theorem~\ref{thm:app29-bz-gamma-quotient}; source problem reference \cite{bulboaca-zayed-gamma}.
\item \label{ledger:app-0030}\textbf{APP-0030.} Ramanujan integral Turan window is false. Solved by Theorem~\ref{thm:app30-ramanujan-turan}; source problem reference \cite{mishra-swaminathan-ramanujan}.
\item \label{ledger:app-0031}\textbf{APP-0031.} Baricz gamma-quotient Bernstein test is false. Solved by Counterexample~\ref{res:baricz-gamma-quotient-bf-forall-negative-answer}; source problem reference \cite{baricz-turan}.
\item \label{ledger:app-0032}\textbf{APP-0032.} From's all-\(L\) Mills-ratio bound family is explicit. Solved by Proposition~\ref{res:from-mills-all-l-alternating-r-bound-family}; source problem reference \cite{from-mills}.
\item \label{ledger:app-0033}\textbf{APP-0033.} Ramanujan antiderivative is complete Bernstein. Solved by Proposition~\ref{res:ramanujan-antiderivative-complete-bernstein}; source problem reference \cite{mishra-swaminathan-ramanujan}.
\item \label{ledger:app-0034}\textbf{APP-0034.} Bulboaca-Zayed gamma-quotient monotonicity solved. Solved by Proposition~\ref{res:bulboaca-zayed-gamma-quotient-full-monotonicity}; source problem reference \cite{bulboaca-zayed-gamma}.
\item \label{ledger:app-0035}\textbf{APP-0035.} Keady self-bijection inverse-CM question is negative. Solved by Counterexample~\ref{res:keady-q3-self-bijection-inverse-cm-negative-answer}; source problem reference \cite{app-app-0035-source}.
\item \label{ledger:app-0036}\textbf{APP-0036.} Baskakov even-line complete-monotonicity fails at \(r=8\). Solved by Counterexample~\ref{res:not-baskakov-higher-power-even-conjecture}; source problem reference \cite{abel-gawronski-neuschel}.
\item \label{ledger:app-0037}\textbf{APP-0037.} Du-Wang \(h\_3\) is increasing exactly on \(\frac12\le a\le1\). Solved by Counterexample~\ref{res:du-wang-h3-open-problem-2-classification}; source problem reference \cite{du-wang-h3}.
\item \label{ledger:app-0038}\textbf{APP-0038.} Ma-Weigert derivative-sign regions form a chain. Solved by Proposition~\ref{res:ma-weigert-log-function-dk-chain-conjecture}; source problem reference \cite{ma-weigert-log-functions}.
\item \label{ledger:app-0039}\textbf{APP-0039.} Ramanujan Turan-window complete-monotonicity interval is false. Solved by Counterexample~\ref{res:ramanujan-turan-window-negative-answer}; source problem reference \cite{mishra-swaminathan-ramanujan}.
\item \label{ledger:app-0040}\textbf{APP-0040.} Yang-Tian Bessel-\(W\) power-Bernstein conjecture is true. Solved by Corollary~\ref{res:yt-bessel-w-power-bernstein-conjecture}; source problem reference \cite{app-app-0040-source}.
\item \label{ledger:app-0041}\textbf{APP-0041.} Qi \(h\_\lambda\) degree-four complete-monotonicity conjecture is false. Solved by Counterexample~\ref{res:not-qi-hlambda-degree4-conjecture}; source problem reference \cite{app-app-0041-source}.
\item \label{ledger:app-0042}\textbf{APP-0042.} Modified Bessel square-root log-concavity fails. Solved by Counterexample~\ref{res:not-bessel-i-sqrt-log-concavity-nu-ge-0}; source problem reference \cite{app-app-0042-source}.
\item \label{ledger:app-0043}\textbf{APP-0043.} Riccati log-concavity inequality for \(I\_\nu\) is false. Solved by Corollary~\ref{res:not-bessel-i-riccati-log-concavity-inequality}; source problem reference \cite{app-app-0042-source}.
\item \label{ledger:app-0044}\textbf{APP-0044.} Bessel ratio quadratic lower bound is false. Solved by Counterexample~\ref{res:not-bessel-i-ratio-quadratic-bound}; source problem reference \cite{app-app-0042-source}.
\item \label{ledger:app-0045}\textbf{APP-0045.} Three-regime Bessel log-concavity certificate route is refuted. Solved by Corollary~\ref{res:not-bessel-i-split-regime-log-concavity-certificate}; source problem reference \cite{app-app-0042-source}.
\item \label{ledger:app-0046}\textbf{APP-0046.} Tao--Sawin finite Weyl \(\ell^1\) minimum is exact. Solved by Proposition~\ref{res:tao-sawin-weyl-l1-exact-minimum}; source problem reference \cite{app-app-0046-source}.
\item \label{ledger:app-0047}\textbf{APP-0047.} BPV noncentral chi-square HCM range has no positive noncentrality. Solved by Counterexample~\ref{res:bpv-noncentral-chi-square-hcm-no-positive-noncentrality}; source problem reference \cite{app-app-0047-source}.
\item \label{ledger:app-0048}\textbf{APP-0048.} BMR tau-Gauss ordinary concavity is false. Solved by Counterexample~\ref{res:bmr-tau-gauss-not-globally-concave}; source problem reference \cite{app-app-0048-source}.
\item \label{ledger:app-0049}\textbf{APP-0049.} Baricz arithmetic/arithmetic zero-balanced hypergeometric threshold. Solved by Proposition~\ref{res:baricz-aa-zero-balanced-concavity-threshold}; source problem reference \cite{app-app-0049-source}.
\item \label{ledger:app-0050}\textbf{APP-0050.} Stolarsky shifted power means with \(p>1\) are not Bernstein. Solved by Counterexample~\ref{res:not-stolarsky-power-mean-pgt1-bernstein-all-shifts}; source problem reference \cite{app-app-0050-source}.
\end{enumerate}

\section{Zeta-law calculus}

\begin{definition}[Riemann zeta probability law]
\phantomsection\addcontentsline{toc}{subsubsection}{Definition~\thetheorem: Riemann zeta probability law}
\label{def:zeta-law}
For \(\beta>1\), define a probability law \(\rho_\beta\) on \(\N\) by
\begin{equation}
\label{eq:zeta-law}
\rho_\beta(n)=\frac{n^{-\beta}}{\zeta(\beta)},\qquad n\in\N.
\end{equation}
Equivalently, with \(E(n)=\log n\) and \(Z(\beta)=\zeta(\beta)\),
\begin{equation}
\label{eq:gibbs-form}
\rho_\beta(n)=\frac{e^{-\beta E(n)}}{Z(\beta)}.
\end{equation}
\end{definition}

\begin{definition}[Zeta free energy]
\phantomsection\addcontentsline{toc}{subsubsection}{Definition~\thetheorem: Zeta free energy}
\label{def:zeta-free-energy}
The zeta free energy is
\begin{equation}
\label{eq:zeta-free-energy}
A(\beta)=\log\zeta(\beta),\qquad \beta>1.
\end{equation}
\end{definition}

\begin{proposition}[Free-energy derivatives]
\phantomsection\addcontentsline{toc}{subsubsection}{Proposition~\thetheorem: Free-energy derivatives}
\label{prop:free-energy-derivatives}
For every \(\beta>1\),
\begin{align}
\label{eq:free-energy-derivatives}
A'(\beta)&=-\E_\beta[\log N],&
A''(\beta)&=\Var_\beta(\log N).
\end{align}
\end{proposition}

\begin{proof}
Absolute convergence of the differentiated Dirichlet series on compact subintervals of \((1,\infty)\) permits termwise differentiation.  Thus
\begin{equation}
\label{eq:first-free-derivative-proof}
A'(\beta)=\frac{\zeta'(\beta)}{\zeta(\beta)}
=-\frac{\sum_{n\ge1}(\log n)n^{-\beta}}{\zeta(\beta)}
=-\E_\beta[\log N].
\end{equation}
Differentiating once more gives
\begin{equation}
\label{eq:second-free-derivative-proof}
A''(\beta)=\E_\beta[(\log N)^2]-\E_\beta[\log N]^2
=\Var_\beta(\log N).
\end{equation}
\end{proof}

Qi, Lim, and Nantomah ask whether \(P_0(x)=\psi(x)\psi(1/x)\) is concave, and more strongly whether \(-P_0''\) is completely monotonic on \((0,\infty)\) \cite{qi-lim-nantomah-polygamma}.  The next theorem answers the stronger question.

\begin{theorem}[APP-0007: Complete monotonicity of reciprocal digamma product curvature]
\phantomsection\addcontentsline{toc}{subsubsection}{Theorem~\thetheorem: APP-0007: Complete monotonicity of reciprocal digamma product curvature}
\label{thm:reciprocal-digamma-product-complete-monotonicity}
Let
\begin{equation}
\label{eq:reciprocal-digamma-p0}
P_0(x)=\psi(x)\psi(1/x),\qquad x>0.
\end{equation}
Then \(-P_0''\) is strictly completely monotonic on \((0,\infty)\).  In particular, \(P_0\) is strictly concave on \((0,\infty)\).
\end{theorem}

\begin{proof}
Direct differentiation gives
\begin{equation}
\label{eq:p0-second-derivative}
\begin{aligned}
P_0''(x)
=&\psi''(x)\psi(1/x)
-\frac{2}{x^2}\psi'(x)\psi'(1/x)\\
&+\frac{2}{x^3}\psi(x)\psi'(1/x)
+\frac{1}{x^4}\psi(x)\psi''(1/x).
\end{aligned}
\end{equation}
Use the shifted Weierstrass expansion
\begin{equation}
\label{eq:digamma-shifted-expansion}
\psi(x)=-\gamma-\frac1x+\sum_{m=1}^{\infty}a_m(x),
\qquad
a_m(x)=\frac{x}{m(m+x)},
\end{equation}
and, after substituting \(1/x\),
\begin{equation}
\label{eq:reciprocal-digamma-shifted-expansion}
\psi(1/x)=-\gamma-x+\sum_{n=1}^{\infty}b_n(x),
\qquad
b_n(x)=\frac1{n(nx+1)}.
\end{equation}
On compact subintervals of \((0,\infty)\), the series for \(a_m\), \(b_n\), and their first two derivatives are uniformly summable; the differentiated double products are dominated by products of summable majorants.  Hence \(P_0\) may be differentiated term by term twice.

Let \(A_0(x)=-\gamma-1/x\) and \(B_0(x)=-\gamma-x\).  The base term satisfies
\begin{equation}
\label{eq:p0-base-term}
-\frac{d^2}{dx^2}\bigl(A_0(x)B_0(x)\bigr)=-\frac{2\gamma}{x^3}.
\end{equation}
For \(n\ge1\), elementary rational differentiation and Laplace inversion give
\begin{equation}
\label{eq:p0-a0bn-laplace}
-\frac{d^2}{dx^2}\bigl(A_0(x)b_n(x)\bigr)
=\int_0^\infty e^{-xt}t^2\frac{n+(\gamma-n)e^{-t/n}}{n^2}\,dt,
\end{equation}
and, for \(m\ge1\),
\begin{equation}
\label{eq:p0-b0am-laplace}
-\frac{d^2}{dx^2}\bigl(B_0(x)a_m(x)\bigr)
=\int_0^\infty e^{-xt}t^2(m-\gamma)e^{-mt}\,dt.
\end{equation}
For \(mn>1\),
\begin{equation}
\label{eq:p0-ambn-laplace}
-\frac{d^2}{dx^2}\bigl(a_m(x)b_n(x)\bigr)
=\int_0^\infty e^{-xt}t^2
\left[
\frac{e^{-t/n}}{mn^2(mn-1)}
-\frac{e^{-mt}}{n(mn-1)}
\right]dt.
\end{equation}
The exceptional product is \(a_1(x)b_1(x)=x/(x+1)^2\), and
\begin{equation}
\label{eq:p0-exceptional-laplace}
-\frac{d^2}{dx^2}\frac{x}{(x+1)^2}
=\int_0^\infty e^{-xt}t^2(t-1)e^{-t}\,dt.
\end{equation}
The identities above follow from
\begin{equation}
\label{eq:p0-laplace-powers}
\frac1{(x+c)^3}=\int_0^\infty e^{-xt}\frac{t^2}{2}e^{-ct}\,dt,
\qquad
\frac1{(x+c)^4}=\int_0^\infty e^{-xt}\frac{t^3}{6}e^{-ct}\,dt.
\end{equation}

The coefficients of the double products group by the two identities
\begin{equation}
\label{eq:p0-group-n}
\sum_{m=1}^{\infty}\frac1{m(mn-1)}
=-\psi(1-1/n)-\gamma
\qquad(n\ge2)
\end{equation}
and
\begin{equation}
\label{eq:p0-group-m}
\sum_{n=1}^{\infty}\frac1{n(mn-1)}
=-\psi(1-1/m)-\gamma
\qquad(m\ge2).
\end{equation}
Both are immediate from \(\psi(z+1)+\gamma=\sum_{k=1}^\infty z/(k(k+z))\), with \(z=-1/n\) and \(z=-1/m\).  Combining \eqref{eq:p0-base-term}--\eqref{eq:p0-group-m} gives
\begin{equation}
\label{eq:p0-laplace-kernel-representation}
-P_0''(x)=\int_0^\infty e^{-xt}K(t)\,dt,
\end{equation}
where
\begin{equation}
\label{eq:p0-kernel}
\begin{aligned}
K(t)=t^2\Bigg[
&1-\gamma+(t-1)e^{-t}\\
&+\sum_{n=2}^{\infty}
\frac{n(1-e^{-t/n})-\psi(1-1/n)e^{-t/n}}{n^2}
+\sum_{m=2}^{\infty}\bigl(m+\psi(1-1/m)\bigr)e^{-mt}
\Bigg].
\end{aligned}
\end{equation}

It remains to prove \(K(t)>0\).  Since \(\psi\) is increasing and \(1-1/n\in[1/2,1)\),
\begin{equation}
\label{eq:p0-minus-psi-positive}
-\psi(1-1/n)>0\qquad(n\ge2).
\end{equation}
Also
\begin{equation}
\label{eq:p0-m-sum-positive}
m+\psi(1-1/m)\ge 2-\gamma-2\log2>0
\qquad(m\ge2),
\end{equation}
because \(\psi(1-1/m)\ge\psi(1/2)=-\gamma-2\log2\).  Thus the two sums in \eqref{eq:p0-kernel} are positive term by term, apart from the possible need to dominate \(1-\gamma+(t-1)e^{-t}\).

For \(n\ge2\),
\begin{equation}
\label{eq:p0-minus-psi-lower}
-\psi(1-1/n)
=\gamma+\int_{1-1/n}^{1}\psi'(u)\,du
>\gamma+\int_{1-1/n}^{1}\frac{du}{u^2}
=\gamma+\frac1{n-1}.
\end{equation}
Consequently the \(n\)-sum in \eqref{eq:p0-kernel} is bounded below by
\begin{equation}
\label{eq:p0-n-sum-lower}
Ae^{-t/2},
\qquad
A=\gamma(\zeta(2)-1)+2-\zeta(2).
\end{equation}
Therefore \(K(t)/t^2\) is bounded below by
\begin{equation}
\label{eq:p0-f-lower}
f(t)=1-\gamma+(t-1)e^{-t}+Ae^{-t/2}.
\end{equation}
Now
\begin{equation}
\label{eq:p0-f-endpoints}
f(0)=A-\gamma=(2-\zeta(2))(1-\gamma)>0,
\qquad
\lim_{t\to\infty}f(t)=1-\gamma>0.
\end{equation}
Moreover,
\begin{equation}
\label{eq:p0-f-derivative}
f'(t)=e^{-t}\left(2-t-\frac{A}{2}e^{t/2}\right),
\end{equation}
and the factor in parentheses is strictly decreasing.  Hence \(f\) has at most one critical point, and any such point is a maximum.  Its minimum on \([0,\infty)\) is attained at an endpoint, so \(f(t)>0\).  Thus \(K(t)>0\) for \(t>0\).

Differentiating under \eqref{eq:p0-laplace-kernel-representation} yields
\begin{equation}
\label{eq:p0-complete-monotonicity}
(-1)^r\frac{d^r}{dx^r}\bigl(-P_0''(x)\bigr)
=\int_0^\infty t^r e^{-xt}K(t)\,dt>0
\end{equation}
for all \(r\ge0\).  This proves strict complete monotonicity of \(-P_0''\), and the case \(r=0\) gives strict concavity of \(P_0\).
\end{proof}

Qi, Lim, and Nantomah also ask whether the higher-order products \(P_n(x)=\psi^{(n)}(x)\psi^{(n)}(1/x)\), \(n\in\N\), are convex, and more strongly whether \(P_n''\) is completely monotonic.  The following result refutes the stronger assertion.  It does not settle the weaker convexity question.

\begin{theorem}[APP-0008: Counterexample to complete monotonicity of higher-order polygamma product curvature]
\phantomsection\addcontentsline{toc}{subsubsection}{Theorem~\thetheorem: APP-0008: Counterexample to complete monotonicity of higher-order polygamma product cur...}
\label{thm:pn-complete-monotonicity-counterexample}
Let
\begin{equation}
\label{eq:pn-product}
P_n(x)=\psi^{(n)}(x)\psi^{(n)}(1/x),\qquad n\in\N,\ x>0.
\end{equation}
Then \(P_n''\) is not completely monotonic for all \(n\in\N\).  More precisely, for every \(n\ge29\),
\begin{equation}
\label{eq:pn-high-order-third-positive}
P_n'''(2)>0,
\end{equation}
and, already at lower order,
\begin{equation}
\label{eq:p7-sixth-negative}
P_7^{(6)}(3)<0.
\end{equation}
\end{theorem}

\begin{proof}
For \(m\ge1\), define the sign-normalized polygamma function
\begin{equation}
\label{eq:sign-normalized-polygamma}
A_m(x)=(-1)^{m+1}\psi^{(m)}(x)
=m!\sum_{k=0}^{\infty}(x+k)^{-m-1}.
\end{equation}
Then \(A_m(x)>0\), \(A_m'(x)=-A_{m+1}(x)\), and \(P_n(x)=A_n(x)A_n(1/x)\).  Differentiating twice gives
\begin{equation}
\label{eq:pn-second-derivative-normalized}
\begin{aligned}
P_n''(x)=&
A_{n+2}(x)A_n(1/x)
-\frac{2}{x^2}A_{n+1}(x)A_{n+1}(1/x)\\
&-\frac{2}{x^3}A_n(x)A_{n+1}(1/x)
+\frac{1}{x^4}A_n(x)A_{n+2}(1/x).
\end{aligned}
\end{equation}
First prove \eqref{eq:pn-high-order-third-positive}.  Put \(p=n+1\), and normalize
\begin{equation}
\label{eq:qp-normalized-product}
Q_p(x)=\frac{P_n(x)}{(n!)^2}
=\sum_{m,\ell\ge0}
\left(\frac{x}{(x+m)(1+\ell x)}\right)^p.
\end{equation}
At \(x=2\), the term \((m,\ell)=(0,0)\) is constant, and the unique largest nonconstant base is \((m,\ell)=(1,0)\), equal to \(2/3\).  Write
\begin{equation}
\label{eq:qp-dominant-split}
Q_p(x)=1+\left(\frac{x}{x+1}\right)^p+R_p(x).
\end{equation}
For \(h(x)=x/(x+1)\), direct differentiation gives
\begin{equation}
\label{eq:h-third-derivative}
\frac{d^3}{dx^3}h(x)^p
=p\left(\frac{x}{x+1}\right)^p
\frac{p^2-6px-3p+6x^2+6x+2}{x^3(x+1)^3}.
\end{equation}
Therefore
\begin{equation}
\label{eq:h-third-at-two}
\frac{d^3}{dx^3}h(x)^p\bigg|_{x=2}
=\frac{p(p^2-15p+38)}{216}\left(\frac23\right)^p.
\end{equation}
For \(p\ge30\), this is at least
\begin{equation}
\label{eq:h-third-lower}
\frac{p^3}{432}\left(\frac23\right)^p.
\end{equation}

For a general summand
\begin{equation}
\label{eq:fmell-general}
F_{m,\ell}(x)=
\left(\frac{x}{(x+m)(1+\ell x)}\right)^p,
\end{equation}
let \(u_{m,\ell}=d(\log F_{m,\ell}^{1/p})/dx\).  At \(x=2\),
\begin{equation}
\label{eq:umell-bounds}
|u_{m,\ell}|\le\frac12,\qquad
|u_{m,\ell}'|\le\frac14,\qquad
|u_{m,\ell}''|\le\frac14.
\end{equation}
Since
\begin{equation}
\label{eq:fmell-third-log}
F_{m,\ell}'''=F_{m,\ell}\left(p^3u_{m,\ell}^3+3p^2u_{m,\ell}u_{m,\ell}'+pu_{m,\ell}''\right),
\end{equation}
we have \(|F_{m,\ell}'''(2)|\le p^3F_{m,\ell}(2)\) for \(p\ge2\).  The remaining mass in \(R_p\) is bounded by
\begin{equation}
\label{eq:rp-mass-bound}
\sum_{(m,\ell)\ne(0,0),(1,0)}F_{m,\ell}(2)<8\cdot2^{-p}
\qquad(p\ge3).
\end{equation}
Indeed, the \(m\ge2,\ell=0\) part is at most \(3\cdot2^{-p}\), while the \(\ell\ge1\) part is at most \(2\cdot2^{-p}(1+(2/3)^p+3\cdot2^{-p})\).
Combining \eqref{eq:h-third-lower} and \eqref{eq:rp-mass-bound},
\begin{equation}
\label{eq:qp-third-lower}
Q_p'''(2)\ge
p^3 2^{-p}
\left[
\frac1{432}\left(\frac43\right)^p-8
\right].
\end{equation}
Because \((4/3)^{30}>432\cdot8\), \eqref{eq:qp-third-lower} is positive for \(p\ge30\).  Hence \(P_n'''(2)=(n!)^2Q_{n+1}'''(2)>0\) for \(n\ge29\).  Complete monotonicity of \(P_n''\) would require \((P_n'')'=P_n'''\le0\), so \(P_n''\) is not completely monotonic for \(n\ge29\).

We now record the lower-order obstruction \eqref{eq:p7-sixth-negative}.  For bookkeeping, put
\begin{equation}
\label{eq:tpij-definition}
T_{p,i,j}(x)=x^{-p}A_i(x)A_j(1/x).
\end{equation}
The identity \(A_m'=-A_{m+1}\) gives the exact recurrence
\begin{equation}
\label{eq:tpij-derivative}
\frac{d}{dx}T_{p,i,j}(x)
=-pT_{p+1,i,j}(x)-T_{p,i+1,j}(x)+T_{p+2,i,j+1}(x).
\end{equation}
Starting with \eqref{eq:pn-second-derivative-normalized}, applying \eqref{eq:tpij-derivative} four times with \(n=7\) expresses \(P_7^{(6)}(3)\) as an exact finite linear combination of \(22\) terms \(T_{p,i,j}(3)\).

The sign of this finite expression can be certified by rational intervals.  For rational \(x>0\) and \(N\ge1\), monotone integral comparison in \eqref{eq:sign-normalized-polygamma} gives
\begin{equation}
\label{eq:am-rational-lower}
m!\left(\sum_{k=0}^{N-1}(x+k)^{-m-1}+\frac{(x+N)^{-m}}{m}\right)
\le A_m(x),
\end{equation}
and
\begin{equation}
\label{eq:am-rational-upper}
A_m(x)\le
m!\left(\sum_{k=0}^{N-1}(x+k)^{-m-1}+\frac{(x+N)^{-m}}{m}+(x+N)^{-m-1}\right).
\end{equation}
Using \(N=80\) in \eqref{eq:am-rational-lower} and \eqref{eq:am-rational-upper} for every \(A_m(3)\) and \(A_m(1/3)\) appearing in the \(22\)-term expression, and choosing upper or lower endpoints according to the sign of each coefficient, yields
\begin{equation}
\label{eq:p7-sixth-certified-lower}
-74998.154649422399828421756528409981771664439555587
<P_7^{(6)}(3),
\end{equation}
and
\begin{equation}
\label{eq:p7-sixth-certified-upper}
P_7^{(6)}(3)
<-74997.331312439140215056900814096498575289999496021.
\end{equation}
Thus \(P_7^{(6)}(3)<0\).  If \(P_7''\) were completely monotonic, then
\begin{equation}
\label{eq:p7-cm-required-sign}
(-1)^4\frac{d^4}{dx^4}P_7''(3)=P_7^{(6)}(3)\ge0,
\end{equation}
contradicting \eqref{eq:p7-sixth-certified-upper}.  Hence the proposed complete-monotonicity strengthening for all \(n\in\N\) is false in both an infinite high-order family and already at \(n=7\).
\end{proof}

\begin{proposition}[Zeta-law calculus]
\phantomsection\addcontentsline{toc}{subsubsection}{Proposition~\thetheorem: Zeta-law calculus}
\label{prop:zeta-law-calculus}
Let \(\beta>1\), and let \(N\) have law \(\rho_\beta\).  For every \(r\ge0\),
\begin{equation}
\label{eq:zeta-law-moment}
\E_\beta[N^{-r}]=\frac{\zeta(\beta+r)}{\zeta(\beta)}.
\end{equation}
Moreover, for \(\alpha,\beta>1\),
\begin{equation}
\label{eq:zeta-law-relative-entropy}
D(\rho_\alpha\Vert\rho_\beta)
=(\beta-\alpha)\E_\alpha[\log N]
+\log\zeta(\beta)-\log\zeta(\alpha).
\end{equation}
\end{proposition}

\begin{proof}
The moment identity follows by direct summation:
\begin{equation}
\label{eq:zeta-law-moment-proof}
\E_\beta[N^{-r}]
=\sum_{n=1}^{\infty}n^{-r}\frac{n^{-\beta}}{\zeta(\beta)}
=\frac{\zeta(\beta+r)}{\zeta(\beta)}.
\end{equation}
For relative entropy, substitute \eqref{eq:zeta-law}:
\begin{equation}
\label{eq:zeta-law-relative-entropy-proof}
D(\rho_\alpha\Vert\rho_\beta)
=\sum_{n=1}^{\infty}\rho_\alpha(n)
\left((\beta-\alpha)\log n+\log\zeta(\beta)-\log\zeta(\alpha)\right),
\end{equation}
which is \eqref{eq:zeta-law-relative-entropy}.
\end{proof}

\section{Tail zeta laws and inverse-tail floors}

\begin{definition}[Tail zeta partition function]
\phantomsection\addcontentsline{toc}{subsubsection}{Definition~\thetheorem: Tail zeta partition function}
\label{def:tail-zeta-partition}
For \(s>1\) and \(n\in\N\), define the tail zeta partition function by
\begin{equation}
\label{eq:tail-zeta-partition}
Z_{s,n}=\zeta_n(s)=\sum_{k=n}^{\infty}k^{-s}.
\end{equation}
The corresponding tail zeta law is the probability distribution on \(\{n,n+1,\ldots\}\) given by
\begin{equation}
\label{eq:tail-zeta-law}
\rho_{s,n}(k)=\frac{k^{-s}}{\zeta_n(s)},\qquad k\ge n.
\end{equation}
\end{definition}

Exact formulas for integer parts of reciprocal zeta tails have been studied for small integer exponents.  Known formulas include \(s=2,3,4,5\), with a sufficiently-large-\(n\) formula for \(s=6\); Kim and Song record that no analogous formula was known for integer \(s>6\) at the time of their work \cite{kim-song-tail}.  The 2024 asymptotic inverse-tail results of Pan and Wu provide related context \cite{pan-wu-tail}.  The next two theorems give exact formulas at \(s=7\) and \(s=8\).  The first uses a single rational telescoping enclosure, while the second needs adjacent corrected asymptotic truncations.

\begin{theorem}[APP-0005: Exact inverse-tail floor formula at s=7]
\phantomsection\addcontentsline{toc}{subsubsection}{Theorem~\thetheorem: APP-0005: Exact inverse-tail floor formula at s=7}
\label{thm:tail-s7-floor}
Let
\begin{equation}
\label{eq:tail-s7}
T_7(n)=\zeta_n(7)=\sum_{k=n}^{\infty}\frac1{k^7}.
\end{equation}
Define
\begin{equation}
\label{eq:tail-s7-q}
Q(n)=120n^6-360n^5+660n^4-720n^3+354n^2-54n+375,
\qquad
P(n)=\frac{Q(n)}{20}.
\end{equation}
Then, for every \(n\ge28\),
\begin{equation}
\label{eq:tail-s7-floor}
\left\lfloor T_7(n)^{-1}\right\rfloor=\left\lfloor P(n)\right\rfloor.
\end{equation}
For \(1\le n\le27\), the exact values are
\begin{equation}
\label{eq:tail-s7-finite-table}
\begin{array}{c|r@{\qquad}c|r@{\qquad}c|r}
n&\lfloor T_7(n)^{-1}\rfloor&n&\lfloor T_7(n)^{-1}\rfloor&n&\lfloor T_7(n)^{-1}\rfloor\\
\hline
1&0&10&4495760&19&241765529\\
2&119&11&8167814&20&331399044\\
3&1862&12&14061542&21&447175152\\
4&12573&13&23143975&22&594869693\\
5&54069&14&36668777&23&781167220\\
6&175597&15&56228085&24&1013751716\\
7&471118&16&83808666&25&1301401638\\
8&1100904&17&121852400&26&1654089273\\
9&2317459&18&173321080&27&2083084422
\end{array}.
\end{equation}
\end{theorem}

\begin{proof}
The polynomial \(Q\) is positive on \(\N\), because
\begin{equation}
\label{eq:tail-s7-q-positive}
\begin{aligned}
Q(n)
&=120(n-1)^6+360(n-1)^5+660(n-1)^4+720(n-1)^3\\
&\qquad +354(n-1)^2+54(n-1)+375.
\end{aligned}
\end{equation}
For \(k\ge1\), direct expansion gives
\begin{equation}
\label{eq:tail-s7-lower-difference}
\frac1{k^7}
-\left(\frac{20}{Q(k)}-\frac{20}{Q(k+1)}\right)
=
\frac{9(60284k^4+29176k^2+15625)}
{k^7Q(k)Q(k+1)}>0.
\end{equation}
Summing \eqref{eq:tail-s7-lower-difference} over \(k\ge n\) yields
\begin{equation}
\label{eq:tail-s7-lower-sum}
T_7(n)>\frac{20}{Q(n)}=\frac1{P(n)},
\end{equation}
and therefore \(T_7(n)^{-1}<P(n)\).

For \(k\ge28\), another expansion gives
\begin{equation}
\label{eq:tail-s7-upper-difference}
\left(\frac{20}{Q(k)-3}-\frac{20}{Q(k+1)-3}\right)
-\frac1{k^7}
=
\frac{36(20k^6-14961k^4-7235k^2-3844)}
{k^7(Q(k)-3)(Q(k+1)-3)}.
\end{equation}
The numerator in \eqref{eq:tail-s7-upper-difference} is positive for \(k\ge28\).  Indeed, writing \(k=m+28\), it becomes
\begin{equation}
\label{eq:tail-s7-upper-positive}
20m^6+3360m^5+220239m^4+7105168m^3
+114013021m^2+751143512m+436261580,
\end{equation}
which has positive coefficients for \(m\ge0\).  Hence, for \(n\ge28\),
\begin{equation}
\label{eq:tail-s7-upper-sum}
T_7(n)<\frac{20}{Q(n)-3}=\frac1{P(n)-3/20}.
\end{equation}
Combining \eqref{eq:tail-s7-lower-sum} and \eqref{eq:tail-s7-upper-sum} gives
\begin{equation}
\label{eq:tail-s7-inverse-window}
P(n)-\frac3{20}<T_7(n)^{-1}<P(n).
\end{equation}

Modulo \(20\),
\begin{equation}
\label{eq:tail-s7-q-mod}
Q(n)\equiv 14n^2-14n+15\pmod {20}.
\end{equation}
Since \(n(n-1)\pmod {10}\in\{0,2,6\}\), this implies
\begin{equation}
\label{eq:tail-s7-q-residues}
Q(n)\pmod {20}\in\{3,15,19\}.
\end{equation}
Thus the fractional part of \(P(n)\) is one of \(3/20,15/20,19/20\).  If \(a=\lfloor P(n)\rfloor\), then \(P(n)-3/20\ge a\), and \eqref{eq:tail-s7-inverse-window} gives
\begin{equation}
\label{eq:tail-s7-floor-trap}
a<T_7(n)^{-1}<a+1.
\end{equation}
This proves \eqref{eq:tail-s7-floor} for \(n\ge28\).

For the finite range, \(T_7(1)>1\), so \(0<T_7(1)^{-1}<1\).  For \(2\le n\le27\), let \(M_n\) denote the corresponding value in \eqref{eq:tail-s7-finite-table}.  With \(N=1500\), exact rational arithmetic verifies
\begin{equation}
\label{eq:tail-s7-finite-lower}
\sum_{k=n}^{N}\frac1{k^7}>\frac1{M_n+1}
\end{equation}
and
\begin{equation}
\label{eq:tail-s7-finite-upper}
\sum_{k=n}^{N}\frac1{k^7}+\frac1{6N^6}<\frac1{M_n}.
\end{equation}
The monotone integral estimate
\begin{equation}
\label{eq:tail-s7-integral-tail}
\sum_{k=N+1}^{\infty}\frac1{k^7}<\int_N^\infty x^{-7}\,dx=\frac1{6N^6}
\end{equation}
then gives
\begin{equation}
\label{eq:tail-s7-finite-window}
\frac1{M_n+1}<T_7(n)<\frac1{M_n}.
\end{equation}
Therefore \(M_n<T_7(n)^{-1}<M_n+1\), proving \eqref{eq:tail-s7-finite-table}.  The finite rational certificate is stored with the source as a reproducible calculation.
\end{proof}

\begin{theorem}[APP-0006: Exact inverse-tail floor formula at s=8]
\phantomsection\addcontentsline{toc}{subsubsection}{Theorem~\thetheorem: APP-0006: Exact inverse-tail floor formula at s=8}
\label{thm:tail-s8-floor}
Let
\begin{equation}
\label{eq:tail-s8}
T_8(n)=\zeta_n(8)=\sum_{k=n}^{\infty}\frac1{k^8}.
\end{equation}
Define
\begin{align}
\label{eq:tail-s8-a}
A_8(n)
={}&7n^7-\frac{49}{2}n^6+\frac{637}{12}n^5-\frac{1715}{24}n^4
+\frac{7399}{144}n^3-\frac{1715}{96}n^2+\frac{69685}{1728}n-\frac{65611}{3456}\notag\\
&-\frac{23764069}{103680n}
-\frac{23764069}{207360n^2}
+\frac{2102131409}{1244160n^3}
+\frac{2149659547}{829440n^4}\notag\\
&-\frac{203290800013}{14929920n^5}
-\frac{1146003910541}{29859840n^6}
+\frac{21587525169497}{179159040n^7}.
\end{align}
Then, for every \(n\ge6\),
\begin{equation}
\label{eq:tail-s8-floor}
\left\lfloor T_8(n)^{-1}\right\rfloor=\left\lfloor A_8(n)\right\rfloor.
\end{equation}
For \(1\le n\le5\), the exact values are
\begin{equation}
\label{eq:tail-s8-finite-table}
\begin{array}{c|rrrrr}
n&1&2&3&4&5\\
\hline
\lfloor T_8(n)^{-1}\rfloor&0&245&5844&53503&291407
\end{array}.
\end{equation}
\end{theorem}

\begin{proof}
Set
\begin{equation}
\label{eq:tail-s8-b}
B_8(n)=A_8(n)+\frac{200335112892761}{358318080n^8}.
\end{equation}
For \(C=A_8\) or \(B_8\), write \(R_C(k)=1/C(k)\).  Exact denominator clearing gives the two telescoping inequalities
\begin{equation}
\label{eq:tail-s8-lower-inverse-difference}
R_{A_8}(k)-R_{A_8}(k+1)-\frac1{k^8}>0\qquad(k\ge3)
\end{equation}
and
\begin{equation}
\label{eq:tail-s8-upper-inverse-difference}
\frac1{k^8}-\left(R_{B_8}(k)-R_{B_8}(k+1)\right)>0\qquad(k\ge1).
\end{equation}
In \eqref{eq:tail-s8-lower-inverse-difference}, the cleared numerator is a polynomial in \(k-3\) with all coefficients positive; in \eqref{eq:tail-s8-upper-inverse-difference}, the cleared numerator is a polynomial in \(k-1\) with all coefficients positive.  The exact coefficient verification is recorded in the rational certificate accompanying the manuscript.

Summing \eqref{eq:tail-s8-lower-inverse-difference} over \(k\ge n\) gives
\begin{equation}
\label{eq:tail-s8-lower-inverse-sum}
T_8(n)<R_{A_8}(n)=\frac1{A_8(n)}
\end{equation}
for \(n\ge3\), and hence \(A_8(n)<T_8(n)^{-1}\).  Summing \eqref{eq:tail-s8-upper-inverse-difference} gives
\begin{equation}
\label{eq:tail-s8-upper-inverse-sum}
T_8(n)>R_{B_8}(n)=\frac1{B_8(n)},
\end{equation}
and hence \(T_8(n)^{-1}<B_8(n)\).  Therefore, for \(n\ge6\),
\begin{equation}
\label{eq:tail-s8-inverse-window}
A_8(n)<T_8(n)^{-1}<B_8(n).
\end{equation}

It remains to see that no integer lies between \(A_8(n)\) and \(B_8(n)\).  Exact integer arithmetic verifies
\begin{equation}
\label{eq:tail-s8-floor-gap-finite}
\lfloor A_8(n)\rfloor=\lfloor B_8(n)\rfloor\qquad(6\le n\le10^6).
\end{equation}
For \(n>10^6\), let \(P_8(n)\) be the polynomial part of \(A_8(n)\), through the constant term.  The numerator of \(P_8(n)\) modulo \(3456\) runs through the odd residue classes, so \(P_8(n)\) is at distance at least \(1/3456\) from the next integer.  The correction \(A_8(n)-P_8(n)\) is negative and has absolute value \(<1/3456\), while
\begin{equation}
\label{eq:tail-s8-gap-small}
0<B_8(n)-A_8(n)<\frac1{3456}.
\end{equation}
Thus \(B_8(n)\) is still below the next integer after \(A_8(n)\).  Combining this with \eqref{eq:tail-s8-inverse-window} proves \eqref{eq:tail-s8-floor}.

For \(1\le n\le5\), \eqref{eq:tail-s8-finite-table} is certified by exact rational partial sums.  With \(N=300\), the displayed value \(M_n\) satisfies, for \(2\le n\le5\),
\begin{equation}
\label{eq:tail-s8-finite-lower}
\sum_{k=n}^{N}\frac1{k^8}>\frac1{M_n+1}
\end{equation}
and
\begin{equation}
\label{eq:tail-s8-finite-upper}
\sum_{k=n}^{N}\frac1{k^8}+\frac1{7N^7}<\frac1{M_n}.
\end{equation}
Since
\begin{equation}
\label{eq:tail-s8-integral-tail}
\sum_{k=N+1}^{\infty}\frac1{k^8}<\int_N^\infty x^{-8}\,dx=\frac1{7N^7},
\end{equation}
we get \(M_n<T_8(n)^{-1}<M_n+1\).  For \(n=1\), \(T_8(1)>1\), so the floor is \(0\).
\end{proof}

\section{Modular successor entropy}

\begin{definition}[Modular residue distribution and successor entropy]
\phantomsection\addcontentsline{toc}{subsubsection}{Definition~\thetheorem: Modular residue distribution and successor entropy}
\label{def:modular-residue-successor}
For \(q\ge2\), define the residue distribution of \(\rho_\beta\) modulo \(q\) by
\begin{equation}
\label{eq:modular-residue-distribution}
\mu_{q,\beta}(a)=
\sum_{\substack{n\ge1\\ n\equiv a\pmod q}}\rho_\beta(n),
\qquad a\in\Z/q\Z.
\end{equation}
Writing \(a+1\) cyclically in \(\Z/q\Z\), define
\begin{equation}
\label{eq:modular-successor-entropy}
B_q(\beta)=
\sum_{a\in\Z/q\Z}
\mu_{q,\beta}(a)\log\frac{\mu_{q,\beta}(a)}{\mu_{q,\beta}(a+1)}.
\end{equation}
\end{definition}

\begin{theorem}[Zeta-law successor entropy and modular resolution]
\phantomsection\addcontentsline{toc}{subsubsection}{Theorem~\thetheorem: Zeta-law successor entropy and modular resolution}
\label{thm:successor-entropy}
For every \(\beta>1\), define
\begin{equation}
\label{eq:successor-entropy}
\Sigma(\beta)=\sum_{n=1}^{\infty}\rho_\beta(n)
\log\frac{\rho_\beta(n)}{\rho_\beta(n+1)}.
\end{equation}
Then
\begin{equation}
\label{eq:successor-entropy-series}
\Sigma(\beta)=
\frac{\beta}{\zeta(\beta)}
\sum_{k=1}^{\infty}\frac{(-1)^{k+1}}{k}\zeta(\beta+k),
\end{equation}
and
\begin{equation}
\label{eq:successor-entropy-modular-resolution}
\Sigma(\beta)=\sup_{q\ge2}B_q(\beta)=\lim_{q\to\infty}B_q(\beta).
\end{equation}
\end{theorem}

\begin{proof}
Since
\begin{equation}
\label{eq:successor-ratio}
\frac{\rho_\beta(n)}{\rho_\beta(n+1)}
=\left(1+\frac1n\right)^\beta,
\end{equation}
we have
\begin{equation}
\label{eq:successor-entropy-log-sum}
\Sigma(\beta)=\frac{\beta}{\zeta(\beta)}
\sum_{n=1}^{\infty}n^{-\beta}\log\left(1+\frac1n\right).
\end{equation}
Using
\begin{equation}
\label{eq:log-series}
\log(1+x)=\sum_{k=1}^{\infty}\frac{(-1)^{k+1}}{k}x^k
\qquad(0<x\le1)
\end{equation}
with Abel limiting at \(x=1\), and summing against \(n^{-\beta}\), gives \eqref{eq:successor-entropy-series}.

For the modular identity, fix \(q\ge2\).  The log-sum inequality gives, for each residue class \(a\),
\begin{equation}
\label{eq:modular-log-sum}
\sum_{\substack{n\ge1\\ n\equiv a\pmod q}}
\rho_\beta(n)\log\frac{\rho_\beta(n)}{\rho_\beta(n+1)}
\ge
\mu_{q,\beta}(a)
\log\frac{\mu_{q,\beta}(a)}
{\sum_{n\equiv a\pmod q}\rho_\beta(n+1)}.
\end{equation}
For \(a\ne0\), the denominator is \(\mu_{q,\beta}(a+1)\).  For \(a=0\), it is \(\mu_{q,\beta}(1)-\rho_\beta(1)\le\mu_{q,\beta}(1)\).  Summing \eqref{eq:modular-log-sum} over \(a\) gives
\begin{equation}
\label{eq:successor-dominates-modular}
\Sigma(\beta)\ge B_q(\beta)
\qquad(q\ge2).
\end{equation}

Conversely, fix \(M\in\N\).  If \(q>M+1\), then for \(1\le a\le M+1\),
\begin{equation}
\label{eq:residue-isolation}
\mu_{q,\beta}(a)=\rho_\beta(a)+O_{\beta,M}(q^{-\beta}),
\end{equation}
and \(\mu_{q,\beta}(0)=q^{-\beta}\).  All cyclic terms in \(B_q(\beta)\) are nonnegative except possibly the boundary term, which satisfies
\begin{equation}
\label{eq:boundary-term}
\mu_{q,\beta}(0)\log\frac{\mu_{q,\beta}(0)}{\mu_{q,\beta}(1)}
=O_\beta(q^{-\beta}\log q).
\end{equation}
Therefore
\begin{equation}
\label{eq:modular-liminf}
\liminf_{q\to\infty}B_q(\beta)
\ge
\sum_{n=1}^{M}\rho_\beta(n)\log\frac{\rho_\beta(n)}{\rho_\beta(n+1)}.
\end{equation}
Letting \(M\to\infty\) gives \(\liminf_qB_q(\beta)\ge\Sigma(\beta)\), which together with \eqref{eq:successor-dominates-modular} proves \eqref{eq:successor-entropy-modular-resolution}.
\end{proof}

\begin{corollary}[Prime-modulus Dirichlet L-resolution]
\phantomsection\addcontentsline{toc}{subsubsection}{Corollary~\thetheorem: Prime-modulus Dirichlet L-resolution}
\label{cor:prime-dirichlet-resolution}
Let \(p\) be prime, let \(\chi\) range over Dirichlet characters modulo \(p\), and define
\begin{equation}
\label{eq:prime-dirichlet-shadow}
S_a(p,\beta)=\sum_{\chi\bmod p}\overline{\chi(a)}L(\beta,\chi)
\qquad(1\le a\le p-1).
\end{equation}
Then
\begin{equation}
\label{eq:prime-residue-formula}
\mu_{p,\beta}(a)=\frac{S_a(p,\beta)}{(p-1)\zeta(\beta)}
\quad(1\le a\le p-1),
\qquad
\mu_{p,\beta}(0)=p^{-\beta}.
\end{equation}
Consequently the prime-modulus functionals obtained by substituting \eqref{eq:prime-residue-formula} into \(B_p(\beta)\) satisfy
\begin{equation}
\label{eq:prime-functional-limit}
\Sigma(\beta)=
\lim_{\substack{p\to\infty\\ p\ {\rm prime}}}
F_p\bigl(\zeta(\beta),\{L(\beta,\chi)\}_{\chi\bmod p}\bigr)
=
\sup_{p\ {\rm prime}}
F_p\bigl(\zeta(\beta),\{L(\beta,\chi)\}_{\chi\bmod p}\bigr).
\end{equation}
\end{corollary}

\begin{proof}
For \(a\not\equiv0\pmod p\), character orthogonality gives
\begin{equation}
\label{eq:character-orthogonality}
\frac1{p-1}\sum_{\chi\bmod p}\overline{\chi(a)}\chi(n)
=
\begin{cases}
1,& n\equiv a\pmod p,\\
0,& n\not\equiv a\pmod p.
\end{cases}
\end{equation}
Therefore
\begin{equation}
\label{eq:dirichlet-residue-sum}
\sum_{\substack{n\ge1\\ n\equiv a\pmod p}}n^{-\beta}
=\frac1{p-1}\sum_{\chi\bmod p}\overline{\chi(a)}L(\beta,\chi).
\end{equation}
Dividing by \(\zeta(\beta)\) gives the nonzero residue formula, while
\begin{equation}
\label{eq:zero-residue-prime}
\mu_{p,\beta}(0)=\frac1{\zeta(\beta)}\sum_{m=1}^{\infty}(mp)^{-\beta}
=p^{-\beta}.
\end{equation}
The functional identity follows by substituting these residue formulas into \(B_p(\beta)\) and applying Theorem \ref{thm:successor-entropy} along the prime sequence.
\end{proof}

\section{Euler-score identities}

\begin{theorem}[Euler-score decomposition]
\phantomsection\addcontentsline{toc}{subsubsection}{Theorem~\thetheorem: Euler-score decomposition}
\label{thm:euler-score-decomposition}
For every \(\beta>1\),
\begin{equation}
\label{eq:euler-score-decomposition}
-\bigl(\log\zeta\bigr)'(\beta)
=\E_\beta[\log N]
=\sum_{d=1}^{\infty}\Lambda(d)\mu_{d,\beta}(0)
=\sum_{d=1}^{\infty}\frac{\Lambda(d)}{d^\beta}.
\end{equation}
Equivalently,
\begin{equation}
\label{eq:log-zeta-von-mangoldt}
-\frac{\zeta'(\beta)}{\zeta(\beta)}
=\sum_{d=1}^{\infty}\frac{\Lambda(d)}{d^\beta}.
\end{equation}
\end{theorem}

\begin{proof}
The von Mangoldt identity is
\begin{equation}
\label{eq:mangoldt-identity}
\log n=\sum_{d\mid n}\Lambda(d).
\end{equation}
By nonnegative summation,
\begin{equation}
\label{eq:euler-score-proof}
\E_\beta[\log N]
=\sum_{n\ge1}\rho_\beta(n)\sum_{d\mid n}\Lambda(d)
=\sum_{d\ge1}\Lambda(d)\Prob_\beta(d\mid N).
\end{equation}
The divisibility probability is
\begin{equation}
\label{eq:divisibility-probability}
\Prob_\beta(d\mid N)
=\frac1{\zeta(\beta)}\sum_{m\ge1}(dm)^{-\beta}
=d^{-\beta}.
\end{equation}
Substitution into \eqref{eq:euler-score-proof} proves \eqref{eq:euler-score-decomposition}; \eqref{eq:log-zeta-von-mangoldt} follows from \eqref{eq:first-free-derivative-proof}.
\end{proof}

\begin{proposition}[Finite Euler-score identity]
\phantomsection\addcontentsline{toc}{subsubsection}{Proposition~\thetheorem: Finite Euler-score identity}
\label{prop:finite-euler-score}
For every \(A\subseteq\{1,\ldots,N\}\),
\begin{equation}
\label{eq:finite-euler-score}
\sum_{n\in A}\log n
=\sum_{d\le N}\Lambda(d)\sum_{m\le N/d}\1_A(md).
\end{equation}
\end{proposition}

\begin{proof}
Reverse the order of summation:
\begin{equation}
\label{eq:finite-euler-score-proof}
\sum_{d\le N}\Lambda(d)\sum_{m\le N/d}\1_A(md)
=\sum_{n\in A}\sum_{d\mid n}\Lambda(d)
=\sum_{n\in A}\log n.
\end{equation}
\end{proof}

\section{Curvature and Mellin-Planck tools}

\begin{lemma}[Uniform positive-axis curvature bound]
\phantomsection\addcontentsline{toc}{subsubsection}{Lemma~\thetheorem: Uniform positive-axis curvature bound}
\label{lem:uniform-positive-axis-curvature}
For every \(s\ge3\),
\begin{equation}
\label{eq:uniform-curvature-bound}
(\log\zeta)''(s)=\Var_s(\log N)<\frac13.
\end{equation}
\end{lemma}

\begin{proof}
By Proposition \ref{prop:free-energy-derivatives},
\begin{equation}
\label{eq:curvature-first-bound}
(\log\zeta)''(s)
=\frac{\zeta''(s)}{\zeta(s)}
-\left(\frac{\zeta'(s)}{\zeta(s)}\right)^2
\le \frac{\zeta''(s)}{\zeta(s)}
\le \zeta''(s).
\end{equation}
For \(s\ge3\),
\begin{equation}
\label{eq:zeta-second-derivative-bound}
\zeta''(s)=\sum_{m=2}^{\infty}\frac{(\log m)^2}{m^s}
\le\sum_{m=2}^{\infty}\frac{(\log m)^2}{m^3}.
\end{equation}
The function \(x\mapsto(\log x)^2x^{-3}\) is decreasing for \(x\ge2\), so
\begin{equation}
\label{eq:curvature-integral-comparison}
\sum_{m=2}^{\infty}\frac{(\log m)^2}{m^3}
\le
\frac{(\log2)^2}{8}
+\int_2^\infty\frac{(\log x)^2}{x^3}\,dx.
\end{equation}
Since
\begin{equation}
\label{eq:curvature-integral-value}
\int_2^\infty\frac{(\log x)^2}{x^3}\,dx
=\frac{(\log2)^2+\log2+\frac12}{8},
\end{equation}
we obtain
\begin{equation}
\label{eq:curvature-final-bound}
(\log\zeta)''(s)
\le
\frac{2(\log2)^2+\log2+\frac12}{8}
<\frac13.
\end{equation}
\end{proof}

\begin{definition}[Mellin-Planck partition function]
\phantomsection\addcontentsline{toc}{subsubsection}{Definition~\thetheorem: Mellin-Planck partition function}
\label{def:mellin-planck-partition}
For \(s>1\), define
\begin{equation}
\label{eq:mellin-planck-partition}
M(s)=\Gamma(s)\zeta(s)
=\int_0^\infty\frac{t^{s-1}}{e^t-1}\,dt.
\end{equation}
The corresponding probability law is
\begin{equation}
\label{eq:mellin-planck-law}
\nu_s(dt)=\frac{t^{s-1}}{(e^t-1)M(s)}\,dt.
\end{equation}
\end{definition}

\begin{proposition}[Mellin-Planck log-convexity]
\phantomsection\addcontentsline{toc}{subsubsection}{Proposition~\thetheorem: Mellin-Planck log-convexity}
\label{prop:mellin-planck-log-convexity}
For \(s>1\),
\begin{equation}
\label{eq:mellin-planck-log-convexity}
\frac{d^2}{ds^2}\log M(s)=\Var_{\nu_s}(\log t)\ge0.
\end{equation}
\end{proposition}

\begin{proof}
Differentiate the logarithm of the integral in \eqref{eq:mellin-planck-partition} under the integral sign.  The first derivative is the \(\nu_s\)-expectation of \(\log t\), and the second derivative is its variance:
\begin{equation}
\label{eq:mellin-planck-variance-proof}
\frac{d^2}{ds^2}\log M(s)
=\E_{\nu_s}[(\log t)^2]-\E_{\nu_s}[\log t]^2.
\end{equation}
\end{proof}

Bulboaca and Zayed proved several monotonicity criteria for functions of the form
\(u(s)/(\Gamma(s+\rho)\Gamma(s))\) and asked for best possible positive
values of \(\rho\) in their examples \cite{bulboaca-zayed-gamma}.  The next
application gives the sharp endpoint in the base reciprocal Gamma-product
case and the corresponding endpoint criterion for monotone logarithmic
derivatives.

\begin{theorem}[APP-0009: Sharp reciprocal Gamma-product monotonicity threshold]
\phantomsection\addcontentsline{toc}{subsubsection}{Theorem~\thetheorem: APP-0009: Sharp reciprocal Gamma-product monotonicity threshold}
\label{thm:gamma-product-sharp-threshold}
Let \(\psi=\Gamma'/\Gamma\), let \(\gamma\) be Euler's constant, and let
\(\rho_*>0\) be the unique solution of
\begin{equation}
\label{eq:gamma-product-rho-star}
\psi(1+\rho_*)=\gamma.
\end{equation}
For \(\rho>-1\), define
\begin{equation}
\label{eq:gamma-product-reciprocal}
\varphi_\rho(s)=\frac{1}{\Gamma(s+\rho)\Gamma(s)},\qquad s\ge1.
\end{equation}
Then \(\varphi_\rho\) is strictly decreasing on \([1,\infty)\) if and only if
\begin{equation}
\label{eq:gamma-product-threshold}
\rho\ge\rho_*.
\end{equation}
Equivalently,
\begin{equation}
\label{eq:gamma-product-W}
W_\rho(s)=\Gamma(s+\rho)\Gamma(s)
\end{equation}
is strictly increasing on \([1,\infty)\) if and only if
\eqref{eq:gamma-product-threshold} holds.  Numerically,
\begin{equation}
\label{eq:gamma-product-rho-star-numeric}
\rho_*=1.2583969670859318174106234224981693941\ldots .
\end{equation}

More generally, let \(u:[1,\infty)\to(0,\infty)\) be differentiable and suppose
\begin{equation}
\label{eq:gamma-product-J}
J(s)=\frac{u'(s)}{u(s)}
\end{equation}
is nonincreasing on \([1,\infty)\).  Then
\begin{equation}
\label{eq:gamma-product-general-Phi}
\Phi_{\rho,u}(s)=\frac{u(s)}{\Gamma(s+\rho)\Gamma(s)}
\end{equation}
is strictly decreasing on \([1,\infty)\) if and only if
\begin{equation}
\label{eq:gamma-product-general-threshold}
\psi(1+\rho)\ge\gamma+J(1).
\end{equation}
\end{theorem}

\begin{proof}
The digamma function is strictly increasing on \((0,\infty)\), since
\begin{equation}
\label{eq:digamma-positive-series}
\psi'(x)=\sum_{k=0}^{\infty}\frac{1}{(x+k)^2}>0.
\end{equation}
Hence \eqref{eq:gamma-product-rho-star} has a unique positive solution.

For \(s\ge1\), the logarithmic derivative of
\eqref{eq:gamma-product-W} is
\begin{equation}
\label{eq:gamma-product-log-derivative}
\frac{d}{ds}\log W_\rho(s)=\psi(s+\rho)+\psi(s).
\end{equation}
The right-hand side is strictly increasing in \(s\).  Its minimum on
\([1,\infty)\) is therefore
\begin{equation}
\label{eq:gamma-product-left-endpoint}
\psi(1+\rho)+\psi(1)=\psi(1+\rho)-\gamma.
\end{equation}
If \(\rho\ge\rho_*\), then \eqref{eq:gamma-product-left-endpoint} is
nonnegative, and it is positive for every \(s>1\).  Thus, for
\(1\le a<b\),
\begin{equation}
\label{eq:gamma-product-integrated-positive}
\log W_\rho(b)-\log W_\rho(a)
=\int_a^b\bigl(\psi(s+\rho)+\psi(s)\bigr)\,ds>0.
\end{equation}
So \(W_\rho\) is strictly increasing and \(\varphi_\rho=1/W_\rho\) is
strictly decreasing.  Conversely, if \(\rho<\rho_*\), then
\eqref{eq:gamma-product-left-endpoint} is negative, so by continuity
\(W_\rho\) decreases initially and \(\varphi_\rho\) increases initially.
This proves the base equivalence.

For the general statement,
\begin{equation}
\label{eq:gamma-product-general-log-derivative}
\frac{d}{ds}\log\Phi_{\rho,u}(s)
=J(s)-\psi(s+\rho)-\psi(s).
\end{equation}
The term \(J\) is nonincreasing by hypothesis, while
\(\psi(s+\rho)+\psi(s)\) is strictly increasing.  Therefore the logarithmic
derivative in \eqref{eq:gamma-product-general-log-derivative} is strictly
decreasing, and its maximum on \([1,\infty)\) is
\begin{equation}
\label{eq:gamma-product-general-left-endpoint}
J(1)-\psi(1+\rho)-\psi(1)
=J(1)-\psi(1+\rho)+\gamma.
\end{equation}
The function \(\Phi_{\rho,u}\) is strictly decreasing exactly when this maximum
is nonpositive, which is \eqref{eq:gamma-product-general-threshold}.  If the
inequality fails, the logarithmic derivative is positive near \(s=1\), so
monotonicity fails.
\end{proof}

Qi, Lim, and Nantomah asked whether the second derivative of \(\log\Gamma(x)+\log\Gamma(1/x)\) is completely monotonic on \((0,\infty)\) \cite{qi-lim-nantomah-polygamma}.  The next theorem answers this in the affirmative by a reciprocal-Weierstrass Laplace kernel.

\begin{theorem}[APP-0004: Complete monotonicity of reciprocal-Gamma curvature]
\phantomsection\addcontentsline{toc}{subsubsection}{Theorem~\thetheorem: APP-0004: Complete monotonicity of reciprocal-Gamma curvature}
\label{thm:reciprocal-gamma-complete-monotonicity}
Let
\begin{equation}
\label{eq:reciprocal-gamma-h}
H(x)=\log\Gamma(x)+\log\Gamma(1/x),\qquad x>0.
\end{equation}
Then \(H''\) is strictly completely monotonic on \((0,\infty)\); that is,
\begin{equation}
\label{eq:reciprocal-gamma-complete-monotonicity}
(-1)^m\frac{d^m}{dx^m}H''(x)>0
\qquad(m\ge0,\ x>0).
\end{equation}
\end{theorem}

\begin{proof}
The Weierstrass product gives, for \(x>0\),
\begin{equation}
\label{eq:gamma-weierstrass-log}
\log\Gamma(x)
=-\log x-\gamma x-\sum_{n=1}^{\infty}
\left(\log\left(1+\frac{x}{n}\right)-\frac{x}{n}\right).
\end{equation}
After substituting \(1/x\) in \eqref{eq:gamma-weierstrass-log}, adding the two formulas, and differentiating twice on compact subintervals of \((0,\infty)\), the logarithmic second-derivative terms cancel and
\begin{equation}
\label{eq:reciprocal-gamma-second-series}
\begin{aligned}
H''(x)
&=\sum_{n=1}^{\infty}\frac1{(x+n)^2}
-\frac{2\gamma}{x^3}\\
&\qquad+
\sum_{n=1}^{\infty}
\left[
\frac1{(x+1/n)^2}
-\frac1{x^2}
+\frac{2}{nx^3}
\right].
\end{aligned}
\end{equation}
The elementary Laplace identities
\begin{equation}
\label{eq:reciprocal-gamma-laplace-identities}
\frac1{(x+a)^2}=\int_0^\infty t e^{-(x+a)t}\,dt,\qquad
\frac1{x^2}=\int_0^\infty t e^{-xt}\,dt,\qquad
\frac{2a}{x^3}=a\int_0^\infty t^2e^{-xt}\,dt
\end{equation}
show that, with \(a=1/n\),
\begin{equation}
\label{eq:reciprocal-gamma-defect-laplace}
\frac1{(x+1/n)^2}-\frac1{x^2}+\frac{2}{nx^3}
=\int_0^\infty e^{-xt}t\left(e^{-t/n}-1+\frac{t}{n}\right)\,dt.
\end{equation}
Also,
\begin{equation}
\label{eq:reciprocal-gamma-trigamma-laplace}
\sum_{n=1}^{\infty}\frac1{(x+n)^2}
=\int_0^\infty e^{-xt}\frac{t}{e^t-1}\,dt,
\end{equation}
and
\begin{equation}
\label{eq:reciprocal-gamma-euler-laplace}
-\frac{2\gamma}{x^3}
=\int_0^\infty e^{-xt}(-\gamma t^2)\,dt.
\end{equation}
Combining \eqref{eq:reciprocal-gamma-second-series}--\eqref{eq:reciprocal-gamma-euler-laplace} gives
\begin{equation}
\label{eq:reciprocal-gamma-kernel-representation}
H''(x)=\int_0^\infty e^{-xt}K(t)\,dt,
\end{equation}
where
\begin{equation}
\label{eq:reciprocal-gamma-kernel}
K(t)=\frac{t}{e^t-1}
+t\sum_{n=1}^{\infty}\left(e^{-t/n}-1+\frac{t}{n}\right)
-\gamma t^2.
\end{equation}

It remains to prove \(K(t)>0\).  Put
\begin{equation}
\label{eq:reciprocal-gamma-phi}
\phi(u)=e^{-u}-1+u.
\end{equation}
Then \(\phi\ge0\), \(\phi'=1-e^{-u}\ge0\), and \(u\mapsto\phi(t/u)\) is positive and decreasing on \([1,\infty)\).  Therefore
\begin{equation}
\label{eq:reciprocal-gamma-sum-integral}
\sum_{n=1}^{\infty}\phi(t/n)\ge\int_1^\infty \phi(t/u)\,du.
\end{equation}
With \(v=t/u\), integration by parts gives
\begin{equation}
\label{eq:reciprocal-gamma-integral-transform}
\int_1^\infty\phi(t/u)\,du
=t\int_0^t\frac{\phi(v)}{v^2}\,dv
=-\phi(t)+t\int_0^t\frac{1-e^{-v}}{v}\,dv.
\end{equation}
Using
\begin{equation}
\label{eq:reciprocal-gamma-e1-identity}
\int_0^t\frac{1-e^{-v}}{v}\,dv
=\gamma+\log t+E_1(t),
\qquad
E_1(t)=\int_t^\infty\frac{e^{-v}}{v}\,dv,
\end{equation}
we obtain
\begin{equation}
\label{eq:reciprocal-gamma-kernel-lower-bound}
\frac{K(t)}{t}
\ge
\frac1{e^t-1}-\phi(t)+t\log t+tE_1(t).
\end{equation}
Since \(E_1(t)>0\) and \(\log t\ge1-1/t\) for \(t>0\),
\begin{equation}
\label{eq:reciprocal-gamma-kernel-positive}
\frac{K(t)}{t}
>
\frac1{e^t-1}-\phi(t)+t-1
=\frac1{e^t-1}-e^{-t}
=\frac{e^{-t}}{e^t-1}>0.
\end{equation}
Thus \(K(t)>0\) for \(t>0\).  Differentiating under the Laplace integral in \eqref{eq:reciprocal-gamma-kernel-representation} gives
\begin{equation}
\label{eq:reciprocal-gamma-alternating-integral}
(-1)^m\frac{d^m}{dx^m}H''(x)
=\int_0^\infty t^m e^{-xt}K(t)\,dt>0.
\end{equation}
This proves \eqref{eq:reciprocal-gamma-complete-monotonicity}.
\end{proof}

\section{External zeta inequality proofs}

\begin{theorem}[APP-0002: Affirmative solution of Nantomah zeta positivity problem]
\phantomsection\addcontentsline{toc}{subsubsection}{Theorem~\thetheorem: APP-0002: Affirmative solution of Nantomah zeta positivity problem}
\label{thm:nantomah-positivity}
For every \(n\in\N\),
\begin{equation}
\label{eq:nantomah-positivity}
(n+2)\zeta(n+1)\zeta(n+3)
-(n+1)\zeta(n+2)^2
-\zeta(n+1)\zeta(n+2)>0.
\end{equation}
\end{theorem}

\begin{proof}
Put \(s=n+1\ge2\) and \(Z_s=\zeta(s)\).  The expression in \eqref{eq:nantomah-positivity} becomes
\begin{equation}
\label{eq:nantomah-ks}
K_s=(s+1)Z_sZ_{s+2}-sZ_{s+1}^2-Z_sZ_{s+1}.
\end{equation}
Let \(X(m)=1/m\) under \(\rho_s\).  By \eqref{eq:zeta-law-moment},
\begin{equation}
\label{eq:nantomah-moments}
\E_s[X]=\frac{Z_{s+1}}{Z_s},
\qquad
\E_s[X^2]=\frac{Z_{s+2}}{Z_s}.
\end{equation}
Thus
\begin{equation}
\label{eq:nantomah-moment-form}
\frac{K_s}{Z_s^2}
=(s+1)\E_s[X^2]-s\E_s[X]^2-\E_s[X]
=(s+1)\Var_s(X)+\E_s[X]^2-\E_s[X].
\end{equation}

Write \(a=\zeta(s)-1\), \(b=\zeta(s+1)-1\), and \(c=\zeta(s+2)-1\).  Expanding gives
\begin{align}
\label{eq:nantomah-lq-split}
K_s&=L_s+Q_s,&
L_s&=sa+(s+1)c-(2s+1)b,&
Q_s&=(s+1)ac-sb^2-ab.
\end{align}
The linear part is termwise positive:
\begin{equation}
\label{eq:nantomah-linear-positive}
L_s=
\sum_{m=2}^{\infty}
m^{-s}\frac{(m-1)(s(m-1)-1)}{m^2}
\ge
2^{-s}\frac{s-1}{4}.
\end{equation}
Since \(b\le a/2\) and \(c\ge0\),
\begin{equation}
\label{eq:nantomah-quadratic-correction}
Q_s\ge -sb^2-ab\ge-\frac{s+2}{4}a^2.
\end{equation}
For \(s\ge4\),
\begin{equation}
\label{eq:nantomah-tail-bound}
a=\sum_{m=2}^{\infty}m^{-s}
\le2^{-s}+\int_2^\infty x^{-s}\,dx
=2^{-s}\left(1+\frac2{s-1}\right).
\end{equation}
Equations \eqref{eq:nantomah-linear-positive}, \eqref{eq:nantomah-quadratic-correction}, and \eqref{eq:nantomah-tail-bound} imply
\begin{equation}
\label{eq:nantomah-lower-bound}
K_s\ge
\frac14\left[
2^{-s}(s-1)
-(s+2)2^{-2s}\left(1+\frac2{s-1}\right)^2
\right].
\end{equation}
It is therefore enough to show
\begin{equation}
\label{eq:nantomah-sufficient-inequality}
2^s(s-1)>(s+2)\left(1+\frac2{s-1}\right)^2.
\end{equation}
Equivalently,
\begin{equation}
\label{eq:nantomah-f-function}
F(s)=\frac{2^s(s-1)^3}{(s+2)(s+1)^2}>1.
\end{equation}
Now \(F(4)=72/25>1\), and
\begin{equation}
\label{eq:nantomah-f-monotone}
\frac{d}{ds}\log F(s)
=\log2+\frac3{s-1}-\frac1{s+2}-\frac2{s+1}>0
\qquad(s\ge4).
\end{equation}
Thus \(K_s>0\) for \(s\ge4\).

For \(s=2\),
\begin{equation}
\label{eq:nantomah-k2}
K_2=3\zeta(2)\zeta(4)-2\zeta(3)^2-\zeta(2)\zeta(3).
\end{equation}
Using \(\zeta(3)<5/4\), \(\zeta(2)=\pi^2/6\), and \(\zeta(4)=\pi^4/90\),
\begin{equation}
\label{eq:nantomah-k2-bound}
K_2>
\frac{\pi^6}{180}-\frac{5\pi^2}{24}-\frac{25}{8}
>
\frac{10415360504209}{180000000000000}>0.
\end{equation}
For \(s=3\), the estimates \(\zeta(3)>251/216\), \(\zeta(5)>8051/7776\), and \(\zeta(4)<13/12\) yield
\begin{equation}
\label{eq:nantomah-k3-bound}
K_3>
\frac{251}{216}
\left(4\cdot\frac{8051}{7776}-\frac{13}{12}\right)
-3\left(\frac{13}{12}\right)^2
=\frac{13783}{419904}>0.
\end{equation}
Hence \(K_s>0\) for all integers \(s\ge2\), proving the theorem.  This answers Nantomah's question \cite{nantomah}.
\end{proof}

\begin{theorem}[APP-0001: Alzer-Kwong convexity and concavity pattern for reciprocal zeta]
\phantomsection\addcontentsline{toc}{subsubsection}{Theorem~\thetheorem: APP-0001: Alzer-Kwong convexity and concavity pattern for reciprocal zeta}
\label{thm:alzer-kwong}
Let \(F(x)=1/\zeta(x)\).  For every integer \(n\ge1\),
\begin{equation}
\label{eq:alzer-kwong-convex}
F''(x)>0\quad\text{on }(-4n,-4n+2),
\end{equation}
and
\begin{equation}
\label{eq:alzer-kwong-concave}
F''(x)<0\quad\text{on }(-4n-2,-4n).
\end{equation}
\end{theorem}

\begin{proof}
Put \(x=-u\), where \(u>2\).  The functional equation gives
\begin{equation}
\label{eq:zeta-functional-equation-negative-axis}
\zeta(-u)
=-2^{-u}\pi^{-u-1}\sin\left(\frac{\pi u}{2}\right)\Gamma(u+1)\zeta(u+1),
\end{equation}
and hence
\begin{equation}
\label{eq:reciprocal-zeta-negative-axis}
\frac1{\zeta(-u)}
=-\frac{2^u\pi^{u+1}}
{\Gamma(u+1)\zeta(u+1)\sin(\pi u/2)}.
\end{equation}
Define the positive function
\begin{equation}
\label{eq:alzer-g-function}
G(u)=
\frac{2^u\pi^{u+1}}
{\Gamma(u+1)\zeta(u+1)|\sin(\pi u/2)|}.
\end{equation}
On \(u\in(4n-2,4n)\), \(\sin(\pi u/2)<0\), so \(F(-u)=G(u)\).  On \(u\in(4n,4n+2)\), \(\sin(\pi u/2)>0\), so \(F(-u)=-G(u)\).  Since second derivatives are unchanged by \(x=-u\), it remains to prove \(G''(u)>0\).

Let \(\ell(u)=\log G(u)\).  Then
\begin{equation}
\label{eq:alzer-log-g}
\ell(u)=u\log2+(u+1)\log\pi-\log\Gamma(u+1)-\log\zeta(u+1)
-\log\left|\sin\left(\frac{\pi u}{2}\right)\right|.
\end{equation}
Differentiating twice gives
\begin{equation}
\label{eq:alzer-log-g-second}
\ell''(u)
=-\psi'(u+1)-(\log\zeta)''(u+1)
+\frac{\pi^2}{4}\csc^2\left(\frac{\pi u}{2}\right).
\end{equation}
With \(s=u+1>3\),
\begin{equation}
\label{eq:trigamma-bound}
\psi'(s)=\sum_{k=0}^{\infty}\frac1{(s+k)^2}
<
\int_{s-1}^{\infty}\frac{dt}{t^2}
=\frac1{s-1}\le\frac12.
\end{equation}
Combining \eqref{eq:uniform-curvature-bound}, \eqref{eq:trigamma-bound}, and \(\csc^2(\pi u/2)\ge1\), we get
\begin{equation}
\label{eq:alzer-log-convexity}
\ell''(u)>
\frac{\pi^2}{4}-\frac12-\frac13
=\frac{\pi^2}{4}-\frac56>0.
\end{equation}
Since \(G>0\),
\begin{equation}
\label{eq:alzer-g-second-positive}
G''(u)=G(u)(\ell''(u)+\ell'(u)^2)>0.
\end{equation}
The two signs in \eqref{eq:alzer-kwong-convex} and \eqref{eq:alzer-kwong-concave} now follow from \(F(-u)=G(u)\) on \(u\in(4n-2,4n)\) and \(F(-u)=-G(u)\) on \(u\in(4n,4n+2)\).  This proves the Alzer--Kwong sign pattern \cite{alzer-kwong}.
\end{proof}

\begin{theorem}[APP-0003: Generalized Holder inequality for Gamma zeta]
\phantomsection\addcontentsline{toc}{subsubsection}{Theorem~\thetheorem: APP-0003: Generalized Holder inequality for Gamma zeta}
\label{thm:sroysang-holder}
Let \(m\ge2\), \(r\ge1\), \(x_1,\ldots,x_m>1\), and \(p_1,\ldots,p_m>1\) satisfy
\begin{equation}
\label{eq:sroysang-holder-condition}
\sum_{i=1}^{m}\frac1{p_i}=\frac1r.
\end{equation}
Define
\begin{equation}
\label{eq:sroysang-t}
T=r\sum_{i=1}^{m}\frac{x_i}{p_i}.
\end{equation}
Then \(T>1\) and
\begin{equation}
\label{eq:sroysang-generalized-holder}
\zeta(T)\le
\frac{\prod_{i=1}^{m}(\Gamma(x_i)\zeta(x_i))^{r/p_i}}{\Gamma(T)}.
\end{equation}
\end{theorem}

\begin{proof}
Set
\begin{equation}
\label{eq:sroysang-lambda}
\lambda_i=\frac{r}{p_i}.
\end{equation}
Then \(\lambda_i>0\), \(\sum_i\lambda_i=1\), and \(T=\sum_i\lambda_i x_i>1\).  Define
\begin{equation}
\label{eq:sroysang-q}
q_i=\frac{p_i}{r}.
\end{equation}
Then \(\sum_i1/q_i=1\), and the assumptions imply \(q_i>1\) for all \(i\); otherwise one term \(1/p_i\) would already be at least \(1/r\), leaving no room for the remaining positive terms in \eqref{eq:sroysang-holder-condition}.

For \(t>0\), let
\begin{equation}
\label{eq:sroysang-fi}
f_i(t)=\left(\frac{t^{x_i-1}}{e^t-1}\right)^{1/p_i}.
\end{equation}
Then
\begin{equation}
\label{eq:sroysang-product}
\left(\prod_{i=1}^{m}f_i(t)\right)^r
=\frac{t^{r\sum_i(x_i-1)/p_i}}{e^t-1}
=\frac{t^{T-1}}{e^t-1}.
\end{equation}
Consequently,
\begin{equation}
\label{eq:sroysang-lr-norm}
\left\|\prod_{i=1}^{m}f_i\right\|_{L^r(0,\infty)}^r
=\int_0^\infty\frac{t^{T-1}}{e^t-1}\,dt
=\Gamma(T)\zeta(T).
\end{equation}
Generalized Holder gives
\begin{equation}
\label{eq:sroysang-holder-step}
\left\|\prod_{i=1}^{m}f_i\right\|_{L^r}
\le
\prod_{i=1}^{m}\|f_i\|_{L^{p_i}}.
\end{equation}
Moreover,
\begin{equation}
\label{eq:sroysang-pi-norm}
\|f_i\|_{L^{p_i}}^{p_i}
=\int_0^\infty\frac{t^{x_i-1}}{e^t-1}\,dt
=\Gamma(x_i)\zeta(x_i).
\end{equation}
Combining \eqref{eq:sroysang-lr-norm}, \eqref{eq:sroysang-holder-step}, and \eqref{eq:sroysang-pi-norm}, then raising to the power \(r\), gives
\begin{equation}
\label{eq:sroysang-predivide}
\Gamma(T)\zeta(T)
\le
\prod_{i=1}^{m}(\Gamma(x_i)\zeta(x_i))^{r/p_i}.
\end{equation}
Division by \(\Gamma(T)\) proves \eqref{eq:sroysang-generalized-holder}.  This solves the generalized Holder problem in Sroysang's notation, where his \(\xi(s)\) is \(\zeta(s)\) for \(s>1\) \cite{sroysang}.
\end{proof}

\begin{remark}[Free-energy interpretation of the Sroysang inequality]
\phantomsection\addcontentsline{toc}{subsubsection}{Remark~\thetheorem: Free-energy interpretation of the Sroysang inequality}
The same result follows from log-convexity of \(M(s)=\Gamma(s)\zeta(s)\).  With \(\lambda_i=r/p_i\) and \(T=\sum_i\lambda_ix_i\), Jensen's inequality for \(\log M\) gives
\begin{equation}
\label{eq:sroysang-jensen}
M(T)\le\prod_{i=1}^{m}M(x_i)^{\lambda_i}
=\prod_{i=1}^{m}(\Gamma(x_i)\zeta(x_i))^{r/p_i},
\end{equation}
which is equivalent to \eqref{eq:sroysang-generalized-holder}.
\end{remark}

\section{Laplace moment-ratio application}
\label{sec:laplace-moment-ratio-application}

The complete-monotonicity proofs above repeatedly convert special-function curvature questions into positivity statements for Laplace kernels.  The next lemma records a reusable moment-ratio consequence of that viewpoint.

\begin{lemma}[Laplace moment-ratio bridge]
\phantomsection\addcontentsline{toc}{subsubsection}{Lemma~\thetheorem: Laplace moment-ratio bridge}
\label{lem:laplace-moment-ratio-bridge}
Let \(\mu\) be a positive Borel measure on \([0,\infty)\), and assume that the tilted moments
\begin{equation}
\label{eq:laplace-moment-ratio-aj}
a_j(x)=\int_0^\infty t^j e^{-xt}\,d\mu(t)
\end{equation}
are finite and positive for the displayed indices.  Then, for every \(n\ge0\),
\begin{equation}
\label{eq:laplace-moment-ratio-Bn}
B_n(x)=\frac{a_{n+1}(x)}{a_n(x)a_{n+2}(x)}
\end{equation}
is strictly increasing in \(x\).
\end{lemma}

\begin{proof}
Put \(r_j(x)=a_{j+1}(x)/a_j(x)\).  Cauchy--Schwarz gives
\begin{equation}
\label{eq:laplace-moment-log-convexity}
a_{j+1}(x)^2\le a_j(x)a_{j+2}(x),
\end{equation}
so \(r_j(x)\le r_{j+1}(x)\).  Since \(a_j'(x)=-a_{j+1}(x)\), differentiating \eqref{eq:laplace-moment-ratio-Bn} logarithmically gives
\begin{equation}
\label{eq:laplace-moment-ratio-derivative}
\frac{d}{dx}\log B_n(x)
=-\frac{a_{n+2}}{a_{n+1}}
+\frac{a_{n+1}}{a_n}
+\frac{a_{n+3}}{a_{n+2}}
=r_n(x)-r_{n+1}(x)+r_{n+2}(x).
\end{equation}
The right-hand side is at least \(r_n(x)>0\), because \(r_{n+2}(x)\ge r_{n+1}(x)\).  Hence \(B_n\) is strictly increasing.
\end{proof}

Yin and Zhang define the Nielsen \(k\)-beta function by
\begin{equation}
\label{eq:nielsen-k-beta-definition}
\beta_k(x)=\int_0^1\frac{t^{x-1}}{1+t^k}\,dt
=\int_0^\infty\frac{e^{-xt}}{1+e^{-kt}}\,dt,
\qquad k>0,\ x>0.
\end{equation}
They ask for the monotonicity of the derivative ratio involving \(x\beta_k(x)\) \cite{yin-zhang-k-beta}.  The bridge lemma gives the exact parity law.

\begin{theorem}[APP-0010: Nielsen \(k\)-beta derivative-ratio monotonicity]
\phantomsection\addcontentsline{toc}{subsubsection}{Theorem~\thetheorem: APP-0010: Nielsen k-beta derivative-ratio monotonicity}
\label{thm:nielsen-k-beta-derivative-ratio}
Let \(k>0\), let \(n\ge0\), and put \(f_k(x)=x\beta_k(x)\).  Then
\begin{equation}
\label{eq:nielsen-k-beta-ratio}
R_{k,n}(x)=
\frac{f_k^{(n+1)}(x)}
{f_k^{(n)}(x)f_k^{(n+2)}(x)}
\end{equation}
is strictly increasing on \((0,\infty)\) when \(n\) is odd, and strictly decreasing on \((0,\infty)\) when \(n\) is even.
\end{theorem}

\begin{proof}
From \eqref{eq:nielsen-k-beta-definition},
\begin{equation}
\label{eq:nielsen-k-beta-x-integral}
f_k(x)=x\int_0^\infty \frac{e^{-xt}}{1+e^{-kt}}\,dt.
\end{equation}
Integrating by parts gives
\begin{equation}
\label{eq:nielsen-k-beta-laplace-density}
f_k(x)
=\frac12+
\int_0^\infty e^{-xt}
\frac{k e^{-kt}}{(1+e^{-kt})^2}\,dt.
\end{equation}
Thus \(f_k\) is a Laplace transform of a positive measure consisting of an atom \(1/2\) at \(0\) and a positive density on \((0,\infty)\).  Define
\begin{equation}
\label{eq:nielsen-k-beta-aj}
a_j(x)=(-1)^j f_k^{(j)}(x).
\end{equation}
For the needed indices these are positive tilted moments of the measure in \eqref{eq:nielsen-k-beta-laplace-density}.  Lemma~\ref{lem:laplace-moment-ratio-bridge} shows that
\begin{equation}
\label{eq:nielsen-k-beta-Bn-increasing}
\frac{a_{n+1}(x)}{a_n(x)a_{n+2}(x)}
\end{equation}
is strictly increasing.  But
\begin{equation}
\label{eq:nielsen-k-beta-sign}
R_{k,n}(x)=(-1)^{n+1}
\frac{a_{n+1}(x)}{a_n(x)a_{n+2}(x)}.
\end{equation}
Multiplication by \((-1)^{n+1}\) gives strict increase for odd \(n\) and strict decrease for even \(n\).  This resolves the Nielsen \(k\)-beta derivative-ratio problem with the precise parity convention.
\end{proof}

\section{Endpoint and parameter applications}
\label{sec:endpoint-parameter-applications}

The next three applications add one certified endpoint result and two parameter-monotonicity results to the complete-monotonicity layer.  Only fully solved source questions are labelled.

\subsection{Exact \(n=2\) beta-window endpoint}

For \(s>1\), write
\begin{equation}
\label{eq:q2-zs-definition}
Z_s(a)=\sum_{m=0}^\infty (m+a)^{-s}.
\end{equation}
In Qi--Lim--Nantomah's notation \cite{qi-lim-nantomah-polygamma}, put
\begin{equation}
\label{eq:q2-cp-definition}
C_2(x)=\psi''(x)+x\psi^{(3)}(x),
\qquad
P_2(x)=\psi''(x)\psi''(1/x),
\end{equation}
and define
\begin{equation}
\label{eq:q2-beta-window}
\mathcal I_2=
\{\beta\in\mathbb R:\ x^\beta C_2(x)-P_2(x)<0\text{ for all }x>0\}.
\end{equation}
The remaining endpoint is controlled on \(0<x<1\) by
\begin{equation}
\label{eq:q2-ratio-q}
R(x)=\frac{P_2(x)}{C_2(x)}
=\frac{2Z_3(x)Z_3(1/x)}{3xZ_4(x)-Z_3(x)},
\qquad
Q_2(x)=\frac{\log R(x)}{\log x}.
\end{equation}

\begin{theorem}[APP-0011: exact \(n=2\) beta-window endpoint]
\phantomsection\addcontentsline{toc}{subsubsection}{Theorem~\thetheorem: APP-0011: exact n=2 beta-window endpoint}
\label{thm:q2-exact-endpoint}
Let
\begin{equation}
\label{eq:q2-J-bracket}
J=\left[\frac{287345}{1000000},\frac{287346}{1000000}\right].
\end{equation}
There is a unique \(\xi\in J\) at which \(Q_2\) attains its global supremum on \((0,1)\).  Hence, with \(L_2=Q_2(\xi)\),
\begin{equation}
\label{eq:q2-I2-exact}
\mathcal I_2=(L_2,3].
\end{equation}
\end{theorem}

\begin{proof}
The identity \(\psi^{(m)}(x)=(-1)^{m+1}m!Z_{m+1}(x)\) gives \eqref{eq:q2-ratio-q}.  Define
\begin{equation}
\label{eq:q2-lambda}
\Lambda(x)=
-\frac{3Z_4(x)}{Z_3(x)}
+\frac{3Z_4(1/x)}{x^2Z_3(1/x)}
-\frac{6Z_4(x)-12xZ_5(x)}{3xZ_4(x)-Z_3(x)}
\end{equation}
and
\begin{equation}
\label{eq:q2-g-function}
G(x)=x\log x\,\Lambda(x)-\log R(x).
\end{equation}
Differentiating \eqref{eq:q2-ratio-q} gives
\begin{equation}
\label{eq:q2-derivative}
Q_2'(x)=\frac{G(x)}{x(\log x)^2}.
\end{equation}

All inequalities below are rational interval certificates using the Hurwitz-zeta integral-tail enclosure and the atanh logarithm enclosure.  First,
\begin{equation}
\label{eq:q2-endpoint-signs}
G\!\left(\frac{287345}{1000000}\right)>0,
\qquad
G\!\left(\frac{287346}{1000000}\right)<0.
\end{equation}
Second, on
\[
B=\left[\frac{1409}{5000},\frac{293}{1000}\right],
\]
the same enclosure proves
\begin{equation}
\label{eq:q2-h-positive}
H(x)=\Lambda(x)+x\Lambda'(x)>0.
\end{equation}
Since \(G'(x)=\log x\,H(x)\) and \(\log x<0\) on \(B\), \(G\) is strictly decreasing on \(B\).  Together with \eqref{eq:q2-endpoint-signs}, this gives a unique zero \(\xi\in J\), and \eqref{eq:q2-derivative} shows that \(Q_2\) is increasing before \(\xi\) and decreasing after \(\xi\) on \(B\).

It remains only to exclude points outside the compact bracket.  The endpoint certificate gives a rational lower witness \(q_J<Q_2(\xi)\) with
\begin{equation}
\label{eq:q2-qJ-lower}
\frac{115727}{50000}<q_J.
\end{equation}
On \((0,9/50]\) the scaled near-zero bound gives \(Q_2(x)<23/10<q_J\).  On \([9/10,1)\), a finite rational subdivision proves \(R(x)>1\), hence \(Q_2(x)<0<q_J\).  On
\[
\left[\frac9{50},\frac{1409}{5000}\right]
\cup
\left[\frac{293}{1000},\frac9{10}\right],
\]
a finite rational cover proves \(Q_2(x)<q_J\): the left bridge uses
\[
R(x)>Z_3(1/x)>x^{1157/500},
\]
and the remaining intervals use the direct comparison
\[
\log R(x)>\frac{115727}{50000}\log x.
\]
Thus no point outside \(B\) reaches the value at \(\xi\), so \(L_2=Q_2(\xi)\).

The source reduction gives the upper endpoint \(3\) and the convention that the lower endpoint is excluded when equality occurs.  At \(\beta=L_2\), equality occurs at \(x=\xi\); therefore the exact admissible set is \((L_2,3]\).
\end{proof}

\subsection{Baricz \(V_q\) parameter log-convexity}

Baricz records the question whether, for fixed \(x>0\), the function
\begin{equation}
\label{eq:baricz-vq-definition}
V_q(x)=\frac{2e^{x^2}}{\Gamma(q+1)}
\int_x^\infty e^{-t^2}(t^2-x^2)^q\,dt,
\qquad q>-1,
\end{equation}
is log-convex as a function of \(q\) \cite{baricz-turan-thesis}.

\begin{theorem}[APP-0012: Baricz \(V_q\) strict \(q\)-log-convexity]
\phantomsection\addcontentsline{toc}{subsubsection}{Theorem~\thetheorem: APP-0012: Baricz V\_q strict q-log-convexity}
\label{thm:baricz-vq-strict-logconvexity}
For each fixed \(x>0\), the map \(q\mapsto V_q(x)\) is strictly log-convex on \((-1,\infty)\).
\end{theorem}

\begin{proof}
Put \(u=t^2-x^2\).  Then
\begin{equation}
\label{eq:baricz-u-normal-form}
V_q(x)=\frac1{\Gamma(q+1)}
\int_0^\infty e^{-u}u^q(u+x^2)^{-1/2}\,du.
\end{equation}
Using
\[
r^{-1/2}=\frac1{\sqrt\pi}\int_0^\infty s^{-1/2}e^{-rs}\,ds
\]
and Tonelli's theorem, we obtain
\begin{equation}
\label{eq:baricz-laplace-normal-form}
V_q(x)=\frac1{\sqrt\pi}\int_0^\infty
s^{-1/2}e^{-x^2s}(1+s)^{-(q+1)}\,ds.
\end{equation}
Let \(a=q+1>0\).  Differentiation under the integral gives
\begin{equation}
\label{eq:baricz-complete-monotonicity}
(-1)^n\frac{d^n}{da^n}V_{a-1}(x)
=\frac1{\sqrt\pi}\int_0^\infty
(\log(1+s))^n s^{-1/2}e^{-x^2s}(1+s)^{-a}\,ds\ge0.
\end{equation}
Equivalently, after the change of variables \(y=\log(1+s)\), \(V_{a-1}(x)\) is a positive Laplace transform in \(a\).  The representing measure is not concentrated at a single point, because the density in \eqref{eq:baricz-laplace-normal-form} is positive on \(s>0\).  Strict Holder inequality for two different exponents therefore gives strict log-convexity in \(a\), and the shift \(a=q+1\) gives the stated strict log-convexity in \(q\).
\end{proof}

\subsection{Yin \((p,k)\)-digamma sharp necessity}

Yin proved sufficiency of \(\alpha\le1\) for complete monotonicity of
\begin{equation}
\label{eq:yin-pk-delta}
\delta_{p,k,\alpha}(x)
=x^\alpha\left[
\frac1k\log\frac{pkx}{x+k(p+1)}-\psi_{p,k}(x)
\right],
\qquad p\in\mathbb N,\ k>0,
\end{equation}
and asked whether this condition is also necessary \cite{yin-pk-digamma}.

\begin{theorem}[APP-0013: Yin \((p,k)\)-digamma sharp \(\alpha\)-necessity]
\phantomsection\addcontentsline{toc}{subsubsection}{Theorem~\thetheorem: APP-0013: Yin (p,k)-digamma sharp -necessity}
\label{thm:yin-pk-alpha-necessity}
If \(\delta_{p,k,\alpha}\) is completely monotone on \((0,\infty)\), then \(\alpha\le1\).  Combined with Yin's sufficiency theorem, the sharp condition is \(\alpha\le1\).
\end{theorem}

\begin{proof}
The source representation is
\begin{equation}
\label{eq:yin-pk-digamma-formula}
\psi_{p,k}(x)=\frac1k\log(pk)-\sum_{n=0}^p\frac1{x+nk}.
\end{equation}
Therefore the bracket in \eqref{eq:yin-pk-delta} is
\begin{equation}
\label{eq:yin-pk-bracket}
B_{p,k}(x)=
\frac1k\log\frac{x}{x+k(p+1)}
+\sum_{n=0}^p\frac1{x+nk}.
\end{equation}
With \(t=x/k\),
\[
B_{p,k}(x)=\frac1k\left[
\sum_{n=0}^p\frac1{t+n}
+\log\frac{t}{t+p+1}
\right].
\]
Since \(u\mapsto(t+u)^{-1}\) is strictly decreasing,
\[
\sum_{n=0}^p\frac1{t+n}
>
\int_0^{p+1}\frac{du}{t+u}
=\log\frac{t+p+1}{t},
\]
so \(B_{p,k}(x)>0\) for \(x>0\).

As \(x\to0^+\), \eqref{eq:yin-pk-bracket} gives
\begin{equation}
\label{eq:yin-pk-near-zero}
B_{p,k}(x)=\frac1x+O(|\log x|),
\qquad
\delta_{p,k,\alpha}(x)\sim x^{\alpha-1}.
\end{equation}
If \(\alpha>1\), then \(\delta_{p,k,\alpha}(x)\to0\) as \(x\to0^+\).  But a completely monotone nonzero function is positive and nonincreasing on \((0,\infty)\).  For any fixed \(x_0>0\), positivity gives \(\delta_{p,k,\alpha}(x_0)>0\), while the limit at \(0\) gives a point \(0<y<x_0\) with \(\delta_{p,k,\alpha}(y)<\delta_{p,k,\alpha}(x_0)\), contradicting nonincreasingness.  Hence \(\alpha\le1\).
\end{proof}


\section{The Laplace-Transport Application Layer}
\label{sec:transport-layer}

The applications in this layer are organized by one theme.  A problem is first converted
into a positive-kernel statement, usually by a logarithmic second difference, a
Frullani identity, a moment determinant, or a transform normalization.  The
conclusion then follows by transporting that kernel through a permanence
operation: Stieltjes transformation, Bernstein integration, stable
subordination, quotient compression, or triangular coefficient extraction.  This
section records representative public solves.

\begin{theorem}[APP-0014: Ramanujan integral is Stieltjes]
\phantomsection\addcontentsline{toc}{subsubsection}{Theorem~\thetheorem: APP-0014: Ramanujan integral is Stieltjes}
\label{thm:app14-ramanujan-stieltjes}
Let
\[
  I_R(x)=\int_0^\infty {e^{-xt}\over t(\pi^2+\log^2 t)}\,dt,
  \qquad x>0 .
\]
Then \(I_R\) is a Stieltjes function.  Moreover every primitive of the form
\[
  \widetilde I_R(x)=a+\int_0^\infty
       {1-e^{-xt}\over t(\pi^2+\log^2 t)}\,dt,\qquad a\ge0,
\]
is a complete Bernstein function.
\end{theorem}

\begin{proof}
For \(t>0\),
\[
 {1\over t(\pi^2+\log^2t)}
   ={1\over \pi(1+t)}\int_0^1 t^{-u}\sin(\pi u)\,du .
\]
The right side is a positive superposition of completely monotone powers and is
therefore completely monotone as a function of \(t\).  The Laplace transform of
a completely monotone density is Stieltjes.  Integrating the same density
against \(1-e^{-xt}\) gives a complete Bernstein function by the standard
Stieltjes-density characterization of complete Bernstein functions.
\end{proof}

\begin{theorem}[APP-0015: Prabhakar Q-measure by stable subordination]
\phantomsection\addcontentsline{toc}{subsubsection}{Theorem~\thetheorem: APP-0015: Prabhakar Q-measure by stable subordination}
\label{thm:app15-prabhakar-subordination}
Let \(0<\alpha<1\), \(\eta>0\), \(\beta>\alpha\eta\), and \(\lambda>0\).
Let
\[
  E_{\alpha,\beta}^{\eta}(z)
  =\sum_{m=0}^{\infty}{(\eta)_m z^m\over m!\,\Gamma(\alpha m+\beta)}
\]
be the Prabhakar Mittag-Leffler function.  Let
\(P_{\alpha,\beta}^{\eta}\) be the transform-normalized Pollard measure
characterized in the strict Pollard range by
\[
  \int_0^\infty e^{-xr}\,dP_{\alpha,\beta}^{\eta}(r)
  =E_{\alpha,\beta}^{\eta}(-x),\qquad x>0.
\]
Let \(S_\alpha\) be the positive \(\alpha\)-stable variable satisfying
\(\mathbb E e^{-sS_\alpha}=e^{-s^\alpha}\).  Define the finite measure
\[
  Q_{\alpha,\beta}^{\eta}(A\mid\lambda)
  =\int \Pr\{(\lambda r)^{1/\alpha}S_\alpha\in A\}\,
       dP_{\alpha,\beta}^{\eta}(r).
\]
Then
\[
  \int_0^\infty e^{-xt}\,dQ_{\alpha,\beta}^{\eta}(t\mid\lambda)
  =E_{\alpha,\beta}^{\eta}(-\lambda x^\alpha).
\]
\end{theorem}

\begin{proof}
Conditioning on \(r\) and using the stable transform gives
\[
 \mathbb E e^{-x(\lambda r)^{1/\alpha}S_\alpha}=e^{-\lambda r x^\alpha}.
\]
Integration against \(P_{\alpha,\beta}^{\eta}\) gives the Pollard
representation of \(E_{\alpha,\beta}^{\eta}(-\lambda x^\alpha)\).  The total
mass is finite, equal to the value at \(x=0\), and no normalization to a
probability measure is needed for the transform identity.
\end{proof}

\begin{theorem}[APP-0016: Mills-ratio bound at every derivative level]
\phantomsection\addcontentsline{toc}{subsubsection}{Theorem~\thetheorem: APP-0016: Mills-ratio bound at every derivative level}
\label{thm:app16-mills-all-l}
Let
\[
  M_L(t)=\int_0^\infty u^L e^{-tu-u^2/2}\,du,\qquad L=0,1,2,\ldots .
\]
Put \(m_L(t)=M_{L+1}(t)/M_L(t)\) and
\[
 U_L(t)=
 {\sqrt{t^2(t^2+4L+7)^2+4(L+1)(t^2+4L+8)(t^2+4L+6)}
       -t(t^2+4L+7)
  \over 2(t^2+4L+8)} .
\]
Then \(m_L(t)<U_L(t)\) for every \(L\ge0\) and \(t>0\).  With
\[
P_0=1,\quad P_1=t,\quad P_{n+1}=tP_n+nP_{n-1},
\]
and
\[
B_0=0,\quad B_1=1,\quad B_{n+1}=nB_{n-1}-tB_n,
\]
the classical Mills ratio \(r(t)=M_0(t)\) satisfies
\[
 r(t)>{B_{L+1}-U_LB_L\over P_{L+1}+U_LP_L}\quad(L\ge0\ {\rm even}),
\]
and
\[
 r(t)<{U_LB_L-B_{L+1}\over P_{L+1}+U_LP_L}\quad(L\ge1\ {\rm odd}).
\]
\end{theorem}

\begin{proof}
The density defining \(M_L\) is a positive moment density.  Its order-four
Hankel-minor consequence gives
\[
M_{L+4}M_L-4M_{L+3}M_{L+1}+3M_{L+2}^2>0.
\]
Using the recurrence \(M_{n+1}=nM_{n-1}-tM_n\), this determinant inequality
reduces to a quadratic inequality for \(m_L\); solving the admissible positive
root gives \(m_L<U_L\).  Iterating the same recurrence expresses \(M_L\) in the
form
\[
M_L=(-1)^L(P_L r-B_L),
\]
and substituting the level-\(L\) quotient bound gives the displayed alternating
upper and lower bounds for \(r\).
\end{proof}

\begin{theorem}[APP-0017: a gamma quotient need not be Bernstein]
\phantomsection\addcontentsline{toc}{subsubsection}{Theorem~\thetheorem: APP-0017: a gamma quotient need not be Bernstein}
\label{thm:app17-baricz-counterexample}
The assertion that the gamma quotient
\[
 x\longmapsto {\Gamma(x+a)\Gamma(x+b)\over \Gamma(x)\Gamma(x+a+b)}
\]
is Bernstein for all positive \(a,b\) is false.
\end{theorem}

\begin{proof}
Take \(a=2\), \(b=3\), and set \(y=x-2\).  The quotient becomes
\[
  g(y)={y(y+1)\over (y+3)(y+4)} .
\]
A direct differentiation gives
\[
  g'''(y)=
 {36(y^4+8y^3+12y^2-40y-94)\over (y+3)^4(y+4)^4}.
\]
Thus \(g'''(1)<0\).  If \(g\) were Bernstein, then \(g'\) would be completely
monotone and hence \(g'''\ge0\).  The displayed value contradicts that
necessary condition.
\end{proof}

\begin{theorem}[APP-0018: exact supremum for the Qi-Guo tau threshold]
\phantomsection\addcontentsline{toc}{subsubsection}{Theorem~\thetheorem: APP-0018: exact supremum for the Qi-Guo tau threshold}
\label{thm:app18-qg-tau}
For \(s\in\mathbb N\) and \(t>0\), define
\[
  \tau(s,t)={1\over s}\left[t-(t+s+1)
        \left({t\over t+1}\right)^{s+1}\right].
\]
Let \(a_*>0\) be the unique positive solution of
\[
  e^{a_*}=1+a_*+a_*^2 .
\]
Then
\[
  \sup_{s\in\mathbb N,\ t>0}\tau(s,t)
   ={a_*\over 1+a_*+a_*^2}
   =0.298425607525639\ldots .
\]
The supremum is not attained.
\end{theorem}

\begin{proof}
Writing \(u=t/(t+1)\) gives a two-parameter expression on \(0<u<1\).  The
Euler-Lagrange equations reduce the interior extremal sequence to the scaling
limit \(u=e^{-a/s}\).  The limiting profile is maximized exactly when
\(e^a=1+a+a^2\), and the resulting value is \(a/(1+a+a^2)\).  The source
one-variable monotonicity estimates bound the finite-\(s\) expressions by this
limit, while the scaling sequence approaches it.  Hence the number is a
supremum and is not attained at finite \(s,t\).
\end{proof}

\begin{theorem}[APP-0019: exact Bernstein range for the W-symbol]
\phantomsection\addcontentsline{toc}{subsubsection}{Theorem~\thetheorem: APP-0019: exact Bernstein range for the W-symbol}
\label{thm:app19-w-symbol}
Let \(0<\alpha<1\), \(\beta\ge0\), and
\[
  \Phi_{\alpha,\beta}(s)
     =s^\alpha\left(1+(1-\alpha)s^{\alpha-1}\right)^{-\beta}.
\]
Then \(\Phi_{\alpha,\beta}\) is a Bernstein function if and only if
\[
  0\le \beta\le 1 .
\]
\end{theorem}

\begin{proof}
Wakrim's source proves the non-Bernstein obstruction for \(\beta>1\) and
leaves the range \(0<\beta\le1\) as the remaining delicate case.  We prove
that remaining Bernstein half in detail.  No complete-Bernstein or
subordination conclusion is asserted here.

Put
\[
  a=1-\alpha\in(0,1),\qquad y=s^a .
\]
Since \(\alpha=1-a\),
\[
 \Phi_{\alpha,\beta}(s)
 =s^{1-a}(1+a s^{-a})^{-\beta}
 =s^{1-a+a\beta}(s^a+a)^{-\beta}.
\]
Let
\[
  p=1-a+a\beta=\alpha+a\beta .
\]
Differentiating the last expression gives
\[
\begin{aligned}
 \Phi'_{\alpha,\beta}(s)
 &=s^{p-1}(s^a+a)^{-\beta-1}
     \{p(s^a+a)-\beta a s^a\}  \\
 &=s^{p-1}(s^a+a)^{-\beta-1}
     \{\alpha s^a+\alpha a+a^2\beta\}       \\
 &=s^{p-1}(s^a+a)^{-\beta}
     \left(\alpha+{a^2\beta\over s^a+a}\right).
\end{aligned}
\]
Because \(p-1=a(\beta-1)\), this is
\[
  \Phi'_{\alpha,\beta}(s)=H_{\alpha,\beta}(s^a),
\]
where
\[
 H_{\alpha,\beta}(y)
 =y^{\beta-1}(y+a)^{-\beta}
      \left(\alpha+{a^2\beta\over y+a}\right).
\]

We now check complete monotonicity of \(H_{\alpha,\beta}\) for
\(0\le\beta\le1\).  The first factor is
\[
  y^{\beta-1}=y^{-(1-\beta)},
\]
and this is interpreted as follows.  If \(0\le\beta<1\), put
\(r=1-\beta>0\).  Then
\[
  y^{\beta-1}=y^{-r}
  ={1\over\Gamma(r)}
    \int_0^\infty e^{-yt}t^{r-1}\,dt,
\]
so \(y^{\beta-1}\) is completely monotone by the Bernstein--Widder
criterion.  If \(\beta=1\), the same factor is the constant \(1\).  The shifted factor
\((y+a)^{-\beta}\) is completely monotone for \(\beta>0\) by the Laplace
representation
\[
 (y+a)^{-\beta}
 ={1\over\Gamma(\beta)}
   \int_0^\infty e^{-yt}e^{-at}t^{\beta-1}\,dt,
\]
and for \(\beta=0\) it is the constant \(1\).  The last factor
\[
  \alpha+{a^2\beta\over y+a}
\]
is a positive constant plus a nonnegative multiple of the completely monotone
resolvent \((y+a)^{-1}\).  By closure of completely monotone functions under
products and nonnegative sums, \(H_{\alpha,\beta}\) is completely monotone.

The endpoint \(\beta=1\) has no limiting ambiguity; explicitly
\[
 H_{\alpha,1}(y)
 ={\alpha\over y+a}+{a^2\over (y+a)^2},
\]
which is completely monotone.  The endpoint \(\beta=0\) also agrees with the
direct identity \(\Phi_{\alpha,0}(s)=s^\alpha\).

Finally, let \(g(s)=s^a\).  The composition theorem used here is the standard
one: if \(f\) is completely monotone on \((0,\infty)\) and \(g\) is a
Bernstein function with \(g((0,\infty))\subset(0,\infty)\), then \(f\circ g\)
is completely monotone.  Its hypotheses are satisfied in the present case:
\(H_{\alpha,\beta}\) is completely monotone by the factor check above,
\(g(s)=s^a>0\), and \(g\) is a Bernstein function for \(0<a<1\) since
\[
  g'(s)=a s^{a-1}=a s^{-(1-a)}
\]
is completely monotone.  Hence
\(\Phi'_{\alpha,\beta}=H_{\alpha,\beta}\circ g\) is completely monotone for
\(0\le\beta\le1\), so \(\Phi_{\alpha,\beta}\) is a Bernstein function.
Together with Wakrim's source obstruction for \(\beta>1\), this closes the
remaining case and gives the exact Bernstein range.
\end{proof}

\begin{theorem}[APP-0020: Qi \(h_\lambda\) degree-four conjecture is false]
\phantomsection\addcontentsline{toc}{subsubsection}{Theorem~\thetheorem: APP-0020: Qi h\_ degree-four conjecture is false}
\label{thm:app20-qi-hlambda}
Let
\[
  \Psi(x)=[\psi'(x)]^2+\psi''(x),\qquad
  h_\lambda(x)=\Psi(x)-
   {x^2+\lambda x+12\over 12x^4(x+1)^2}.
\]
The conjecture that
\(\deg_x^{\rm cm}h_\lambda=\deg_x^{\rm cm}(-h_\mu)=4\) exactly in the stated
parameter ranges \(\lambda\le0\), \(\mu\ge4\), is false.  Here
\(\deg_x^{\rm cm}g\) denotes the source convention: the supremal power \(r\)
for which \(x^rg(x)\) is completely monotone on \((0,\infty)\).
\end{theorem}

\begin{proof}
Near zero,
\[
 h_\lambda(x)=-{\lambda\over 12x^3}
 +\left({\lambda\over6}+{\pi^2\over3}-{37\over12}\right){1\over x^2}
 +O(x^{-1}).
\]
If \(\lambda<0\), then \(x^4h_\lambda(x)=(-\lambda/12)x+O(x^2)\), whose
derivative is positive near zero.  If \(\lambda=0\), the coefficient
\(\pi^2/3-37/12\) is positive, again giving a positive derivative for the
degree-four transform near zero.  Complete monotonicity would require the first
derivative to be non-positive.  The same small-\(x\) test applied to
\(-x^4h_\mu\) excludes the claimed \(\mu\)-range.  Hence the degree-four
conjecture fails.
\end{proof}

\begin{theorem}[APP-0021: Szabo's cutoff is \(\alpha_0=1\)]
\phantomsection\addcontentsline{toc}{subsubsection}{Theorem~\thetheorem: APP-0021: Szabo's cutoff is \_0=1}
\label{thm:app21-szabo-cutoff}
For \(0<d<1\), define
\[
  H_d(y)=\psi(y+d)-\psi(y)-{d\over y}.
\]
The exact cutoff for complete monotonicity of \(y^\alpha H_d(y)\) is
\[
  \alpha_0=1 .
\]
\end{theorem}

\begin{proof}
The source theorem proves complete monotonicity for \(\alpha\le1\).  Conversely,
as \(y\downarrow0\),
\[
  H_d(y)={1-d\over y}+O(1).
\]
Hence, for \(\alpha>1\),
\[
  y^\alpha H_d(y)=(1-d)y^{\alpha-1}+O(y^\alpha),
\]
whose derivative is positive near zero.  A completely monotone function must be
non-increasing, so no \(\alpha>1\) is admissible.
\end{proof}

\begin{theorem}[APP-0022: Chen-Choi Gurland ratio is log-completely monotone]
\phantomsection\addcontentsline{toc}{subsubsection}{Theorem~\thetheorem: APP-0022: Chen-Choi Gurland ratio is log-completely monotone}
\label{thm:app22-gurland-logcm}
Let
\[
  F(x)={\Gamma(1/x)\Gamma(3/x)\over \Gamma(2/x)^2},\qquad x>0.
\]
Then \(F\) is logarithmically completely monotone; equivalently,
\[
  (-1)^n(\log F)^{(n)}(x)>0,\qquad n=0,1,2,\ldots .
\]
\end{theorem}

\begin{proof}
Set
\[
  A(u)=2\log(u+2)-\log(u+1)-\log(u+3).
\]
The identity
\[
 A(u)=\int_0^\infty e^{-ut}{e^{-t}(1-e^{-t})^2\over t}\,dt
\]
shows that \(A\) is strictly completely monotone.  The Weierstrass product for
\(\Gamma\) gives
\[
  \log F(x)=\log {4\over3}+\sum_{m\ge1}A(mx).
\]
Termwise differentiation is justified on compact subintervals of
\((0,\infty)\), and each differentiated summand has the required strict
alternating sign.  This proves logarithmic complete monotonicity.
\end{proof}

\begin{theorem}[APP-0023: Sokal lambda-derivatives have triangular compression]
\phantomsection\addcontentsline{toc}{subsubsection}{Theorem~\thetheorem: APP-0023: Sokal lambda-derivatives have triangular compression}
\label{thm:app23-sokal-lambda}
Let \(n,k\in\mathbb N_0\), \(\lambda>0\), \(x>0\), and let \(f\) be a
\(C^\infty\) function on \((0,\infty)\).  In Sokal's generalized-Stieltjes
application, the quantities below are the Hankel-type expressions associated
with a generalized Stieltjes function of order \(\lambda\).  For \(j\ge0\), set
\[
  G_j(x)=(-1)^n x^j f^{(n+j)}(x),\qquad a=n+\lambda .
\]
For \(0\le r\le k\), define
\[
  F^{[\lambda]}_{n,r}(x)
   =\sum_{j=0}^r {r\choose j}(a+j)_{r-j}G_j(x).
\]
Then, for every integer \(0\le\ell\le k\),
\[
 \partial_\lambda^\ell F^{[\lambda]}_{n,k}(x)
 =\sum_{r=0}^{k-\ell}{k!\over r!}
   [z^{k-r}](-\log(1-z))^\ell\,F^{[\lambda]}_{n,r}(x).
\]
For \(\ell>k\), \(\partial_\lambda^\ell F^{[\lambda]}_{n,k}(x)=0\).  The
coefficients in this triangular compression are non-negative.
\end{theorem}

\begin{proof}
The exponential generating function is
\[
 \mathcal F(z)
 =(1-z)^{-a}\sum_{j\ge0}G_j(x)
      \left({z\over 1-z}\right)^j{1\over j!}.
\]
Differentiating \(\ell\) times with respect to \(\lambda\) multiplies
\(\mathcal F\) by \((-\log(1-z))^\ell\).  Extracting coefficients gives the
displayed triangular formula.  Since
\(-\log(1-z)=\sum_{m\ge1}z^m/m\) has non-negative coefficients, all compression
coefficients are non-negative.
\end{proof}

\begin{theorem}[APP-0024: Simon gamma quotient is Bernstein]
\phantomsection\addcontentsline{toc}{subsubsection}{Theorem~\thetheorem: APP-0024: Simon gamma quotient is Bernstein}
\label{thm:app24-simon-gamma-quotient}
For \(0<\alpha<1\), define
\[
  F_\alpha(x)={\Gamma(x+\alpha)\over \Gamma(x)x^\alpha},
  \qquad x>0 .
\]
Then \(F_\alpha\) is a Bernstein function.  More precisely,
\(1-F_\alpha\) is completely monotone.
\end{theorem}

\begin{proof}
Put \(\beta=1-\alpha\).  The beta integral gives
\[
 F_\alpha(x)
 =x^\beta{\Gamma(x+\alpha)\over\Gamma(x+1)}
 ={x^\beta\over\Gamma(\beta)}
   \int_0^\infty e^{-xu}(e^u-1)^{-\alpha}\,du .
\]
Define
\[
 J_\alpha(t)=\int_0^t (t-u)^{\alpha-1}(e^u-1)^{-\alpha}\,du .
\]
Laplace convolution gives
\[
 F_\alpha(x)
 ={x\over\Gamma(\alpha)\Gamma(1-\alpha)}
   \int_0^\infty e^{-xt}J_\alpha(t)\,dt .
\]
After the substitution \(u=tv\),
\[
J_\alpha(t)=
\int_0^1(1-v)^{\alpha-1}
\left({t\over e^{tv}-1}\right)^\alpha\,dv .
\]
For each \(v>0\), the function \(t/(e^{tv}-1)\) is decreasing in \(t\).
The integrand is dominated near the endpoints by
\((1-v)^{\alpha-1}v^{-\alpha}\), which is integrable because
\(0<\alpha<1\).  Hence \(J_\alpha\) is decreasing, with
\[
 J_\alpha(0+)=\Gamma(\alpha)\Gamma(1-\alpha),
 \qquad
 J_\alpha(\infty)=0 .
\]
Let \(\mu_\alpha=-dJ_\alpha\).  This is a positive finite measure, and
Stieltjes integration by parts yields
\[
 1-F_\alpha(x)
 ={1\over\Gamma(\alpha)\Gamma(1-\alpha)}
  \int_0^\infty e^{-xt}\,\mu_\alpha(dt).
\]
Thus \(1-F_\alpha\) is completely monotone.  Differentiating this Laplace
representation gives
\[
 F_\alpha'(x)
 ={1\over\Gamma(\alpha)\Gamma(1-\alpha)}
  \int_0^\infty t e^{-xt}\,\mu_\alpha(dt),
\]
so \(F_\alpha'\) is completely monotone.  Therefore \(F_\alpha\) is a
Bernstein function.
\end{proof}

\begin{theorem}[APP-0025: Du--Wang \(h_3\) monotonicity classification]
\phantomsection\addcontentsline{toc}{subsubsection}{Theorem~\thetheorem: APP-0025: Du--Wang h\_3 monotonicity classification}
\label{thm:app25-du-wang-h3}
For \(0<a<2\), define
\[
 h_3(x)=
 {-x^2\psi'(x+a)+2x\psi(x+a)-2\log\Gamma(x+a)\over x},
 \qquad x>0 .
\]
Then \(h_3\) is increasing on \((0,\infty)\) if and only if
\[
  {1\over2}\le a\le1 .
\]
For \(0<a<1/2\) and for \(1<a<2\), \(h_3\) is not monotone on
\((0,\infty)\).
\end{theorem}

\begin{proof}
The source reduction writes \(t=x+a\) and
\[
 h_3'(x)={h_{31}(x+a)\over x^2},
\]
where
\[
 h_{31}(t)=-(t-a)^3\psi''(t)+(t-a)^2\psi'(t)
          -2(t-a)\psi(t)+2\log\Gamma(t).
\]
For \(u>0\), set \(H_a(u)=h_{31}(a+u)\).  Then
\[
 H_a'(u)=-u^2\{u\psi'''(a+u)+2\psi''(a+u)\},
 \qquad
 H_a(0+)=2\log\Gamma(a).
\]

We first prove the middle window.  For \(t>1/2\), put
\[
 S_m(t)=\sum_{n=0}^\infty (n+t)^{-m}.
\]
Since \(\psi''(t)=-2S_3(t)\) and \(\psi'''(t)=6S_4(t)\), it is enough to show
\[
 2S_3(t)-3\left(t-{1\over2}\right)S_4(t)>0 .
\]
Let \(y=t-1/2\) and \(K(x)=1/(2\sinh(x/2))\).  The standard integral
representation for \(S_m\), followed by integration by parts, gives
\[
\begin{aligned}
 2S_3(t)-3yS_4(t)
 &=\int_0^\infty e^{-yx}
   \left(x^2K(x)-{1\over2}(x^3K(x))'\right)\,dx  \\
 &=\int_0^\infty e^{-yx}
   {x^2K(x)\over2}\left({x\over2}\coth{x\over2}-1\right)\,dx .
\end{aligned}
\]
The last integrand is positive because \(z\coth z>1\) for \(z>0\).  Therefore
\[
 -{2\psi''(t)\over\psi'''(t)}>t-{1\over2}.
\]
If \(1/2\le a\le1\), then \(u=t-a\le t-1/2\), so
\[
 2\psi''(t)+u\psi'''(t)<0 .
\]
Thus \(H_a'(u)>0\).  Also \(H_a(0+)=2\log\Gamma(a)\ge0\) on
\([1/2,1]\), since \(\log\Gamma\) decreases there and \(\Gamma(1)=1\).
Hence \(H_a(u)\ge0\), and therefore \(h_3'(x)\ge0\), on the whole middle
window.

It remains to exclude the outer windows.  As \(t\to\infty\), the standard
asymptotics for \(\psi,\psi',\psi''\), and \(\log\Gamma\) give
\[
  h_{31}(t)=(2a-1)\log t+O(1).
\]
If \(0<a<1/2\), then \(h_{31}(a+)=2\log\Gamma(a)>0\), while
\(h_{31}(t)<0\) for all sufficiently large \(t\).  If \(1<a<2\), then
\(h_{31}(a+)=2\log\Gamma(a)<0\), while \(h_{31}(t)>0\) for all sufficiently
large \(t\).  In both cases \(h_3'\) changes sign, so \(h_3\) is not monotone.
\end{proof}

\begin{theorem}[APP-0026: Baskakov even-power conjecture is false]
\phantomsection\addcontentsline{toc}{subsubsection}{Theorem~\thetheorem: APP-0026: Baskakov even-power conjecture is false}
\label{thm:app26-baskakov-r8}
For the Abel--Gawronski--Neuschel higher-power Baskakov family
\[
 f^{[r]}_\alpha(x)
 =(1+x)^{-r\alpha}\sum_{k=0}^\infty
 \binom{-\alpha}{k}^{r}\left({x\over1+x}\right)^{rk},
\]
the universal even-power complete-monotonicity conjecture is false.  Already
at \(\alpha=1\) and \(r=8\),
\[
 f^{[8]}_1(x)={1\over (1+x)^8-x^8}
\]
is not completely monotone on \((0,\infty)\).
\end{theorem}

\begin{proof}
For even \(r=2m\) and \(\alpha=1\), the source definition reduces to
\[
 f^{[2m]}_1(x)={1\over(1+x)^{2m}-x^{2m}}.
\]
Let \(D_m(x)=(1+x)^{2m}-x^{2m}\).  Its roots are
\[
 \lambda_{k,m}=-{1\over2}-{i\over2}\cot{\pi k\over 2m},
 \qquad k=1,\ldots,2m-1,
\]
and the corresponding residues are
\[
 R_{k,m}={(-1)^{m+k}\over2m}
 \left(2\sin{\pi k\over2m}\right)^{2m-2}.
\]
Thus the inverse-Laplace density of \(1/D_m\), if it were nonnegative, would
be
\[
 \rho_m(t)={e^{-t/2}\over2m}
 \left[
 2^{2m-2}
 +2\sum_{k=1}^{m-1}(-1)^{m+k}
 \left(2\sin{\pi k\over2m}\right)^{2m-2}
 \cos\left({t\over2}\cot{\pi k\over2m}\right)
 \right].
\]
For \(m=4\), write \(h_4(t)=e^{t/2}\rho_4(t)\).  The formula simplifies to
\[
\begin{aligned}
 h_4(t)
 &=8+2\cos{t\over2}
 -{(2-\sqrt2)^3\over4}\cos\left({1+\sqrt2\over2}t\right) \\
 &\hspace{1.2cm}
 -{(2+\sqrt2)^3\over4}\cos\left({\sqrt2-1\over2}t\right).
\end{aligned}
\]
At \(t=10\pi\),
\[
 h_4(10\pi)=6+10\cos(5\pi\sqrt2).
\]
Since \(0<5\sqrt2-7<1/4\),
\[
 \cos(5\pi\sqrt2)=-\cos(\pi(5\sqrt2-7))
 <-{\sqrt2\over2}<-{3\over5}.
\]
Consequently \(h_4(10\pi)<0\), and so \(\rho_4(10\pi)<0\).  By uniqueness of
Laplace transforms of positive measures, \(f^{[8]}_1\) cannot be completely
monotone.
\end{proof}

\begin{theorem}[APP-0027: Ma--Weigert derivative regions form a chain]
\phantomsection\addcontentsline{toc}{subsubsection}{Theorem~\thetheorem: APP-0027: Ma--Weigert derivative regions form a chain}
\label{thm:app27-ma-weigert-chain}
Let
\[
 \mathcal F_{1,n}
 =
 \left\{x\mapsto {p(\log x)\over x}:p\in\mathbb R[y]_n\right\},
\]
and let \(L=-d/dx\).  If
\[
 D_k=\{f\in\mathcal F_{1,n}:L^kf(x)\ge0\text{ for all }x>0\},
\]
then
\[
 D_{k+1}\subseteq D_k
 \qquad(k\ge0,\ n\ge0).
\]
\end{theorem}

\begin{proof}
For \(f(x)=p(\log x)/x\), induction gives
\[
 L^k f(x)=x^{-k-1}p_k(\log x)
\]
for a polynomial \(p_k\).  Indeed, if the formula holds at level \(k\), then
\[
 L^{k+1}f(x)
 =x^{-k-2}\bigl((k+1)p_k(\log x)-p_k'(\log x)\bigr).
\]
Thus \(L^k f(x)\to0\) as \(x\to\infty\).  If \(f\in D_{k+1}\), put
\(g=L^k f\).  Then \(-g'=L^{k+1}f\ge0\), and hence, for \(R>x\),
\[
 g(x)-g(R)=\int_x^R L^{k+1}f(t)\,dt .
\]
Letting \(R\to\infty\) gives
\[
 L^k f(x)=\int_x^\infty L^{k+1}f(t)\,dt\ge0.
\]
Therefore \(f\in D_k\).
\end{proof}

\begin{theorem}[APP-0028: divisor-polygamma parity correction]
\phantomsection\addcontentsline{toc}{subsubsection}{Theorem~\thetheorem: APP-0028: divisor-polygamma parity correction}
\label{thm:app28-qa-divisor-polygamma}
For \(n\in\mathbb N\), set
\[
 f_n(x)=\sum_{km=n}\bigl[\psi^{(k)}(x)\bigr]^m,\qquad x>0.
\]
Then \(f_n\) is completely monotone for every odd \(n\).  The asserted
opposite parity law that \(f_{2\ell}\) is never completely monotone is false:
\[
 f_2(x)=[\psi'(x)]^2+\psi''(x)
\]
is completely monotone on \((0,\infty)\).
\end{theorem}

\begin{proof}
The standard polygamma representation gives, for \(k\ge1\),
\[
 (-1)^{k+1}\psi^{(k)}(x)
 =\int_0^\infty {t^ke^{-xt}\over1-e^{-t}}\,dt .
\]
Thus \(\psi^{(k)}\) is a positive completely monotone function whenever
\(k\) is odd.  If \(n\) is odd and \(km=n\), both \(k\) and \(m\) are odd, so
each summand \([\psi^{(k)}]^m\) is a finite product of positive completely
monotone functions.  Finite sums preserve complete monotonicity.

For \(n=2\), the divisor pairs are \((1,2)\) and \((2,1)\), hence
\[
 f_2(x)=[\psi'(x)]^2+\psi''(x).
\]
Qi--Agarwal record the sharp theorem that
\[
 [\psi'(x)]^2+\lambda\psi''(x)
\]
is completely monotone on \((0,\infty)\) if and only if \(\lambda\le1\).
Taking \(\lambda=1\) proves that \(f_2\) is completely monotone.  This refutes
the stated even-parity clause at \(\ell=1\).
\end{proof}

\begin{theorem}[APP-0029: Bulboaca--Zayed Gamma quotient one-critical-point theorem]
\phantomsection\addcontentsline{toc}{subsubsection}{Theorem~\thetheorem: APP-0029: Bulboaca--Zayed Gamma quotient one-critical-point theorem}
\label{thm:app29-bz-gamma-quotient}
Let
\[
 F(x)=
 {\log\Gamma(x+1)\over \log(x^2+6)-\log(x+6)}
\]
with the removable values filled in on \((-1,\infty)\).  Then \(F'\) has a
unique zero \(x_m\) on \((-1,\infty)\).  Moreover
\[
 F'(x)<0\quad(-1<x<x_m),\qquad
 F'(x)>0\quad(x>x_m),
\]
and numerically \(x_m=1.126207061051643\ldots\).
\end{theorem}

\begin{proof}
The source proves \(F'<0\) on \((-1,1)\).  For \(x>1\), set
\[
 G(x)=\log\Gamma(x+1),\qquad
 D(x)=\log{x^2+6\over x+6},
\]
and
\[
 N_0(x)=\psi(x+1)D(x)-G(x)D'(x).
\]
Since \(F'=N_0/D^2\), the sign of \(F'\) on \(x>1\) is the sign of \(N_0\).
A direct right-tail estimate gives \(N_0(x)>0\) for \(x\ge8\): with
\[
 P(x)=x^2+12x-6,\qquad Q(x)=(x^2+6)(x+6),
\]
one has \(D'=P/Q\); the bounds \(\psi(x+1)>\log x\),
\(\log\Gamma(x+1)\le x\log x\), and
\[
 D(x)\ge {xP(x)\over Q(x)}
 \qquad(x\ge8)
\]
imply \(N_0(x)>0\).

It remains to prove a single crossing on \((1,8]\).  Define
\[
 r(x)={\psi(x+1)\over D'(x)}
      =\psi(x+1){Q(x)\over P(x)}.
\]
Then
\[
 P(x)^2r'(x)=S(x),
\]
where
\[
 S(x)=\psi'(x+1)Q(x)P(x)
 +\psi(x+1)\{Q'(x)P(x)-Q(x)P'(x)\}.
\]
Using the elementary bounds, valid for \(y\ge2\),
\[
 \psi(y)>\log y-{1\over y},\qquad
 \psi'(y)>{1\over y}+{1\over2y^2},\qquad
 \psi''(y)>-{1\over y^2}-{2\over y^3},
\]
one obtains \(S'(x)>0\) on \([1,8]\).  Also
\[
 S(1)=343\left({\pi^2\over6}-1\right)-539(1-\gamma)<0,
\]
while \(S(8)>0\).  Hence \(r'\) has exactly one zero on \([1,8]\).

Now put \(I(x)=D(x)r(x)-G(x)\).  Then \(N_0(x)=D'(x)I(x)\) and
\[
 I'(x)=D(x)r'(x).
\]
Since \(D(x)>0\) for \(x>1\), \(I\) decreases once and then increases once.
The removable endpoint has \(\lim_{x\downarrow1}I(x)=0\), and the quadratic
leading coefficient of \(N_0\) at \(1\) is negative.  Finally,
\[
 N_0(8)=(H_8-\gamma)\log5-{11\over70}\log(40320)>0.
\]
Thus \(N_0\), and hence \(F'\), has exactly one zero on \((1,8]\).  Combining
this with the source interval \((-1,1)\) and the right-tail estimate proves
the theorem.
\end{proof}

\begin{theorem}[APP-0030: Ramanujan Turan window has a negative answer]
\phantomsection\addcontentsline{toc}{subsubsection}{Theorem~\thetheorem: APP-0030: Ramanujan Turan window has a negative answer}
\label{thm:app30-ramanujan-turan}
Let
\[
 I_R(x)=\int_0^\infty e^{-xt}{dt\over t(\pi^2+\log^2t)}
\]
and, for \(n\ge2\),
\[
 H_n(x;\alpha)
 =
 \bigl(I_R^{(n)}(x)\bigr)^2
 -\alpha I_R^{(n-1)}(x)I_R^{(n+1)}(x).
\]
The proposed complete-monotonicity window
\[
 {n-2\over n-1}<\alpha<{n-1\over n}
\]
is false.  It already fails for \(n=2\).
\end{theorem}

\begin{proof}
Put
\[
 A_k(x)=(-1)^kI_R^{(k)}(x)
 =\int_0^\infty e^{-xt}{t^{k-1}\over \pi^2+\log^2t}\,dt .
\]
Then
\[
 H_2(x;\alpha)=A_2(x)^2-\alpha A_1(x)A_3(x).
\]
Let \(x=e^{-L}\).  After the substitution \(y=xt\),
\[
 A_k(e^{-L})
 =
 e^{kL}\int_0^\infty
 {e^{-y}y^{k-1}\over \pi^2+(L+\log y)^2}\,dy .
\]
Splitting at \(|\log y|\le L^{2/3}\) and expanding the denominator uniformly
on the central region gives, for fixed \(k=1,2,3\),
\[
 A_k(e^{-L})
 =
 e^{kL}L^{-2}\Gamma(k)
 \left(1-{2\psi(k)\over L}+O(L^{-2})\right).
\]
Consequently
\[
 {A_2(e^{-L})^2\over A_1(e^{-L})A_3(e^{-L})}
 =
 {1\over2}
 \left(
 1+{2(\psi(1)+\psi(3)-2\psi(2))\over L}+O(L^{-2})
 \right).
\]
Since
\[
 \psi(1)=-\gamma,\qquad
 \psi(2)=1-\gamma,\qquad
 \psi(3)={3\over2}-\gamma,
\]
the bracketed correction is \(-1/L+O(L^{-2})\), and hence
\[
 {A_2(e^{-L})^2\over A_1(e^{-L})A_3(e^{-L})}
 =
 {1\over2}-{1\over2L}+O(L^{-2}).
\]
Choose \(\alpha_L=1/2-1/(4L)\).  For all sufficiently large \(L\),
\(\alpha_L\in(0,1/2)\) and
\[
 A_2(e^{-L})^2-\alpha_L A_1(e^{-L})A_3(e^{-L})<0.
\]
Thus \(H_2(e^{-L};\alpha_L)<0\).  Complete monotonicity would imply
nonnegativity at order zero, so the source's universal interval statement is
false.
\end{proof}

\section{Late source applications}
\label{sec:late-source-applications}
The following applications are already admitted in the Theory graph and are placed together because they are source-facing consequences of mechanisms developed earlier: Ramanujan moment-ratio asymptotics, Bessel zero partial fractions, and local endpoint obstructions.

\begin{counterexample}[APP-0039: Ramanujan Turan-window complete-monotonicity interval is false]
\phantomsection\addcontentsline{toc}{subsubsection}{Counterexample~\thetheorem: APP-0039: Ramanujan Turan-window complete-monotonicity interval is false}
\label{res:ramanujan-turan-window-negative-answer}
The Mishra--Swaminathan Ramanujan integral Turan-window complete-monotonicity statement has a negative answer: the proposed interval already fails for \(n=2\) and suitable \(\alpha\in(0,1/2)\).  This resolves the source problem recorded as \textbf{APP-0039} \cite{mishra-swaminathan-ramanujan}.
\end{counterexample}

\begin{proof}
Put
\[
A_k(x)=(-1)^k I_R^{(k)}(x)
=\int_0^\infty e^{-xt}\frac{t^{k-1}}{\pi^2+\log^2 t}\,dt
\qquad (k\ge1).
\]
Then \(H_2(x;\alpha)=A_2(x)^2-\alpha A_1(x)A_3(x)\).  Set \(x=e^{-L}\).  With \(y=xt\),
\[
A_k(e^{-L})
=e^{kL}\int_0^\infty
\frac{e^{-y}y^{k-1}}{\pi^2+(L+\log y)^2}\,dy.
\]
For \(k=1,2,3\), splitting the integral into \(|\log y|\le L^{2/3}\) and its complement gives
\[
A_k(e^{-L})
=e^{kL}L^{-2}\Gamma(k)
\left(1-\frac{2\psi(k)}{L}+O(L^{-2})\right).
\]
Therefore
\[
\frac{A_2(e^{-L})^2}{A_1(e^{-L})A_3(e^{-L})}
=\frac{\Gamma(2)^2}{\Gamma(1)\Gamma(3)}
\left(
1+\frac{2(\psi(1)+\psi(3)-2\psi(2))}{L}+O(L^{-2})
\right).
\]
Using \(\psi(1)=-\gamma\), \(\psi(2)=1-\gamma\), and \(\psi(3)=3/2-\gamma\), this ratio equals
\[
\frac12-\frac{1}{2L}+O(L^{-2}).
\]
Choose \(\alpha_L=1/2-1/(4L)\).  For all sufficiently large \(L\), \(\alpha_L\in(0,1/2)\) and \(H_2(e^{-L};\alpha_L)<0\).  Complete monotonicity would imply nonnegativity, so the universal Turan-window statement is false.
\end{proof}


\begin{proposition}[Bessel \(W_\nu\) zero representation]
\phantomsection\addcontentsline{toc}{subsubsection}{Proposition~\thetheorem: Bessel W\_ zero representation}
\label{res:yt-bessel-w-general-zero-partial-fraction}
For every \(\nu>-1\), if \(j_{\nu+1,n}\) are the positive zeros of \(J_{\nu+1}\) and \(W_\nu(z)=zI_\nu(z)/I_{\nu+1}(z)\), then
\[
G_\nu(s):=W_\nu(\sqrt s)=2(\nu+1)+2\sum_{n=1}^\infty\frac{s}{s+j_{\nu+1,n}^2},
\]
and
\[
G_\nu'(s)=2\sum_{n=1}^\infty\frac{j_{\nu+1,n}^2}{(s+j_{\nu+1,n}^2)^2}.
\]
In particular \(G_\nu\) is a Bernstein function.
\end{proposition}

\begin{proof}
Put \(\mu=\nu+1>0\).  The canonical product for \(I_\mu\) gives
\[
I_\mu(z)=\frac{(z/2)^\mu}{\Gamma(\mu+1)}\prod_{n=1}^\infty\left(1+\frac{z^2}{j_{\mu,n}^2}\right).
\]
Hence
\[
z\frac{I_\mu'(z)}{I_\mu(z)}=\mu+2\sum_{n=1}^\infty\frac{z^2}{z^2+j_{\mu,n}^2}.
\]
Using \(I_\mu'(z)=I_{\mu-1}(z)-\mu I_\mu(z)/z\), we obtain
\[
W_\nu(z)=z\frac{I_{\mu-1}(z)}{I_\mu(z)}=2\mu+2\sum_{n=1}^\infty\frac{z^2}{z^2+j_{\mu,n}^2}.
\]
Set \(s=z^2\).  Since \(j_{\nu+1,n}\sim\pi n\), the differentiated sums converge locally uniformly, and
\[
(-1)^k\frac{d^k}{ds^k}
\frac{j_{\nu+1,n}^2}{(s+j_{\nu+1,n}^2)^2}
=(k+1)!\frac{j_{\nu+1,n}^2}{(s+j_{\nu+1,n}^2)^{k+2}}\ge0.
\]
Thus \(G_\nu'\) is completely monotone, and \(G_\nu\) is Bernstein.
\end{proof}

\begin{corollary}[APP-0040: Yang-Tian Bessel-\(W\) power-Bernstein conjecture is true]
\phantomsection\addcontentsline{toc}{subsubsection}{Corollary~\thetheorem: APP-0040: Yang-Tian Bessel-W power-Bernstein conjecture is true}
\label{res:yt-bessel-w-power-bernstein-conjecture}
For \(\tau\in(0,1/2]\) and \(\nu>-1\), \(x\mapsto W_\nu(x^\tau)\) is a Bernstein function on \((0,\infty)\).  This resolves the source problem recorded as \textbf{APP-0040} \cite{app-app-0040-source}.
\end{corollary}

\begin{proof}
By Proposition~\ref{res:yt-bessel-w-general-zero-partial-fraction}, \(G_\nu\) is Bernstein.  If \(0<\tau\le1/2\), then \(x^{2\tau}\) is Bernstein.  The Bernstein composition closure used in Theorem~\ref{thm:app19-w-symbol} gives
\[
W_\nu(x^\tau)=G_\nu(x^{2\tau}),
\]
which is Bernstein.
\end{proof}


\begin{counterexample}[APP-0041: Qi \(h_\lambda\) degree-four complete-monotonicity conjecture is false]
\phantomsection\addcontentsline{toc}{subsubsection}{Counterexample~\thetheorem: APP-0041: Qi h\_ degree-four complete-monotonicity conjecture is false}
\label{res:not-qi-hlambda-degree4-conjecture}
Let
\[
\Psi(x)=[\psi'(x)]^2+\psi''(x),\qquad
h_\lambda(x)=\Psi(x)-\frac{x^2+\lambda x+12}{12x^4(x+1)^2}.
\]
The conjecture that \(\deg_x^{\rm cm}h_\lambda=\deg_x^{\rm cm}(-h_\mu)=4\) exactly for \(\lambda\le0\) and \(\mu\ge4\) is false.  This resolves the source problem recorded as \textbf{APP-0041} \cite{app-app-0041-source}.
\end{counterexample}

\begin{proof}
At \(x=0^+\),
\[
\psi'(x)=\frac1{x^2}+\zeta(2)-2\zeta(3)x+3\zeta(4)x^2+O(x^3),
\qquad
\psi''(x)=-\frac2{x^3}-2\zeta(3)+6\zeta(4)x+O(x^2).
\]
Thus
\[
\Psi(x)=\frac1{x^4}-\frac2{x^3}+\frac{\pi^2}{3x^2}-\frac{4\zeta(3)}x+O(1).
\]
Also
\[
\frac{x^2+\lambda x+12}{12x^4(1+x)^2}
=\frac1{x^4}+\left(\frac{\lambda}{12}-2\right)\frac1{x^3}
+\left(\frac{37}{12}-\frac{\lambda}{6}\right)\frac1{x^2}+O(x^{-1}).
\]
Hence
\[
h_\lambda(x)
=-\frac{\lambda}{12x^3}
+\left(\frac{\lambda}{6}+\frac{\pi^2}{3}-\frac{37}{12}\right)\frac1{x^2}
+O(x^{-1}).
\]
If \(\lambda<0\), then \(\frac{d}{dx}[x^4h_\lambda(x)]=-\lambda/12+O(x)>0\) for small \(x>0\), so \(x^4h_\lambda\) is not completely monotone.  If \(\lambda=0\), then
\[
x^4h_0(x)=\left(\frac{\pi^2}{3}-\frac{37}{12}\right)x^2+O(x^3),
\]
and \(4\pi^2>37\) again gives a positive first derivative near zero.  Similarly, for \(\mu\ge4\),
\[
\frac{d}{dx}[x^4(-h_\mu(x))]=\frac{\mu}{12}+O(x)>0
\]
near zero.  The degree-four transforms required by the conjecture therefore fail on the conjectured source ranges.
\end{proof}

\section{Bessel endpoint obstruction}
\label{sec:bessel-endpoint-obstruction}

The next group belongs to the same endpoint-obstruction mechanism used throughout the paper: a universal positivity or concavity statement is reduced to a one-point Riccati certificate, and exact rational bounds then decide the sign.  The same certificate resolves the square-root log-concavity problem, the Riccati form of that problem, a related quadratic ratio bound, and the proposed three-regime certificate route.

\begin{counterexample}[APP-0042: Modified Bessel square-root log-concavity fails]
\phantomsection\addcontentsline{toc}{subsubsection}{Counterexample~\thetheorem: APP-0042: Modified Bessel square-root log-concavity fails}
\label{res:not-bessel-i-sqrt-log-concavity-nu-ge-0}
The assertion that, for every \(\nu\ge0\), the function
\[
u\longmapsto \sqrt{u}\,I_\nu(u)
\]
is strictly log-concave on \((0,\infty)\) is false.  More precisely, at \((\nu,u)=(0,10)\) the logarithmic second derivative is positive.  This resolves the source problem recorded as \textbf{APP-0042} \cite{app-app-0042-source}.
\end{counterexample}

\begin{proof}
Let
\[
g_\nu(u)=\log(\sqrt u\,I_\nu(u)),
\qquad
r_\nu(u)=\frac{I_\nu'(u)}{I_\nu(u)}.
\]
Then
\[
g_\nu''(u)=r_\nu'(u)-\frac{1}{2u^2}.
\]
The modified Bessel equation
\[
u^2I_\nu''(u)+uI_\nu'(u)-(u^2+\nu^2)I_\nu(u)=0
\]
gives
\[
\frac{I_\nu''(u)}{I_\nu(u)}
=1+\frac{\nu^2}{u^2}-\frac{r_\nu(u)}{u},
\]
and therefore
\[
g_\nu''(u)
=
1+\frac{\nu^2-\frac12}{u^2}
-\frac{r_\nu(u)}{u}
-r_\nu(u)^2.
\]
For \(\nu=0\), \(r_0(u)=I_1(u)/I_0(u)\).  At \(u=10\),
\[
g_0''(10)=\frac{199}{200}-\frac{r_0(10)}{10}-r_0(10)^2.
\]
It is enough to prove
\[
r_0(10)<\frac{9487}{10000},
\]
because
\[
\frac{199}{200}
-\frac{1}{10}\frac{9487}{10000}
-\left(\frac{9487}{10000}\right)^2
=\frac{9831}{100000000}>0.
\]

The modified Bessel series are
\[
I_0(10)=\sum_{k=0}^\infty\frac{25^k}{(k!)^2},
\qquad
I_1(10)=\sum_{k=0}^\infty\frac{5\cdot25^k}{k!(k+1)!}.
\]
Using the first \(21\) terms as a lower bound for \(I_0(10)\), and bounding the \(I_1\)-tail after \(k=20\) by a geometric tail with ratio
\[
q=\frac{25}{22\cdot23}<1,
\]
the audited rational certificate gives an upper bound \(U_1>I_1(10)\) satisfying
\[
10000\,U_1<9487\sum_{k=0}^{20}\frac{25^k}{(k!)^2}.
\]
Hence \(I_1(10)/I_0(10)<9487/10000\), so \(g_0''(10)>0\).  The claimed universal strict log-concavity is therefore refuted.
\end{proof}

\begin{corollary}[APP-0043: Riccati log-concavity inequality for \(I_\nu\) is false]
\phantomsection\addcontentsline{toc}{subsubsection}{Corollary~\thetheorem: APP-0043: Riccati log-concavity inequality for I\_ is false}
\label{res:not-bessel-i-riccati-log-concavity-inequality}
The inequality
\[
1+\frac{\nu^2-\frac12}{u^2}-\frac{r_\nu(u)}{u}-r_\nu(u)^2<0,
\qquad r_\nu(u)=\frac{I_\nu'(u)}{I_\nu(u)},
\]
does not hold for every \(\nu\ge0\) and \(u>0\).  This resolves the source problem recorded as \textbf{APP-0043} \cite{app-app-0042-source}.
\end{corollary}

\begin{proof}
The proof of Counterexample~\ref{res:not-bessel-i-sqrt-log-concavity-nu-ge-0} gives
\[
1-\frac{1}{2\cdot10^2}-\frac{r_0(10)}{10}-r_0(10)^2
=g_0''(10)>0.
\]
This is the negation of the claimed strict Riccati inequality at \((\nu,u)=(0,10)\).
\end{proof}

\begin{counterexample}[APP-0044: Bessel ratio quadratic lower bound is false]
\phantomsection\addcontentsline{toc}{subsubsection}{Counterexample~\thetheorem: APP-0044: Bessel ratio quadratic lower bound is false}
\label{res:not-bessel-i-ratio-quadratic-bound}
The claimed inequality
\[
q_\nu(u)^2+\frac{(2\nu+1)q_\nu(u)}{u}+\frac{\nu+\frac12}{u^2}>1,
\qquad q_\nu(u)=\frac{I_{\nu+1}(u)}{I_\nu(u)},
\]
is false for all \(\nu\ge0\) and \(u>0\).  This resolves the source problem recorded as \textbf{APP-0044} \cite{app-app-0042-source}.
\end{counterexample}

\begin{proof}
At \(\nu=0\) and \(u=10\), the same rational certificate used above gives
\[
q_0(10)=\frac{I_1(10)}{I_0(10)}<\frac{9487}{10000}.
\]
Therefore
\[
q_0(10)^2+\frac{q_0(10)}{10}+\frac1{200}
<
\left(\frac{9487}{10000}\right)^2
+\frac1{10}\frac{9487}{10000}
+\frac1{200}
=1-\frac{9831}{100000000}<1.
\]
This contradicts the proposed strict lower bound.
\end{proof}

\begin{corollary}[APP-0045: Three-regime Bessel log-concavity certificate route is refuted]
\phantomsection\addcontentsline{toc}{subsubsection}{Corollary~\thetheorem: APP-0045: Three-regime Bessel log-concavity certificate route is refuted}
\label{res:not-bessel-i-split-regime-log-concavity-certificate}
No valid three-regime certificate can prove the universal Riccati log-concavity inequality for \(I_\nu\), because the universal inequality itself is false.  This resolves the source problem recorded as \textbf{APP-0045} \cite{app-app-0042-source}.
\end{corollary}

\begin{proof}
Corollary~\ref{res:not-bessel-i-riccati-log-concavity-inequality} gives an explicit point \((\nu,u)=(0,10)\) where the target inequality fails.  A certificate proving the inequality for all \(\nu\ge0\) and \(u>0\) would imply the false universal statement.  Hence the proposed certificate route cannot exist.
\end{proof}

\section{Integrated applications added after the seed}
\label{sec:post-seed-integrated-applications}
The following applications have admitted Theory proof nodes and are rendered here before the append-only application ledger.  They are placed after the seed exposition so the public APP ledger can point to mathematical statements with proofs rather than to ledger-only entries.

\begin{counterexample}[APP-0031: Baricz gamma-quotient Bernstein test is false]
\phantomsection\addcontentsline{toc}{subsubsection}{Counterexample~\thetheorem: APP-0031: Baricz gamma-quotient Bernstein test is false}
\label{res:baricz-gamma-quotient-bf-forall-negative-answer}
Baricz's question asking whether \(x\mapsto\Gamma(x)\Gamma(x-a+b)/(\Gamma(x-a)\Gamma(x+b))\) is a Bernstein function on \((a,\infty)\) for every \(a,b>0\) has a negative answer.

This resolves the source problem recorded as \textbf{APP-0031} \cite{baricz-turan}.
\end{counterexample}

\begin{proof}
Take \(a=2\), \(b=3\), and set \(y=x-2>0\). Then by the Gamma recurrence,
\[
\frac{\Gamma(x)\Gamma(x-a+b)}{\Gamma(x-a)\Gamma(x+b)}
=\frac{\Gamma(x)\Gamma(x+1)}{\Gamma(x-2)\Gamma(x+3)}
=\frac{(x-2)(x-1)}{(x+1)(x+2)}
=\frac{y(y+1)}{(y+3)(y+4)}.
\]

Let
\[
g(y)=\frac{y(y+1)}{(y+3)(y+4)},\qquad y>0.
\]

If \(g\) were a Bernstein function, then \(g'\) would be completely monotone. In particular \(g'''\ge0\) would be necessary.

Exact symbolic differentiation gives
\[
g'(y)=\frac{6(y^2+4y+2)}{(y+3)^2(y+4)^2},
\]
\[
g''(y)=
-\frac{12(y^3+6y^2+6y-10)}{(y+3)^3(y+4)^3},
\]
and
\[
g'''(y)=
\frac{36(y^4+8y^3+12y^2-40y-94)}{(y+3)^4(y+4)^4}.
\]

At \(y=1\),
\[
g'''(1)=-\frac{1017}{40000}<0.
\]

Therefore \(g'\) is not completely monotone, \(g\) is not Bernstein, and the Baricz for-all-\(a,b\) Bernstein question has a negative answer.
\end{proof}

\begin{proposition}[APP-0032: From's all-\(L\) Mills-ratio bound family is explicit]
\phantomsection\addcontentsline{toc}{subsubsection}{Proposition~\thetheorem: APP-0032: From's all-L Mills-ratio bound family is explicit}
\label{res:from-mills-all-l-alternating-r-bound-family}
Let \(P_0=1\), \(P_1=t\), \(P_{n+1}=tP_n+nP_{n-1}\), and \(B_0=0\), \(B_1=1\), \(B_{n+1}=nB_{n-1}-tB_n\). With \(U_L(t)\) as in the all-\(L\) moment-ratio bound, the standard normal Mills ratio satisfies, for even \(L\ge0\), \(r(t)>[B_{L+1}-U_LB_L]/[P_{L+1}+U_LP_L]\), and for odd \(L\ge1\), \(r(t)<[U_LB_L-B_{L+1}]/[P_{L+1}+U_LP_L]\), for all \(t\ge0\).

This resolves the source problem recorded as \textbf{APP-0032} \cite{from-mills}.
\end{proposition}

\begin{proof}
Set
\[
m_L(t)=\frac{M_{L+1}(t)}{M_L(t)}.
\]

The recurrence gives
\[
\frac{M_{L+2}}{M_L}=L+1-tm_L,
\]
\[
\frac{M_{L+3}}{M_L}=(t^2+L+2)m_L-t(L+1),
\]
and
\[
\frac{M_{L+4}}{M_L}=(L+1)(t^2+L+3)-t(t^2+2L+5)m_L.
\]

Substitution into the determinant inequality and expansion give
\[
(L+1)(t^2+4L+6)-t(t^2+4L+7)m_L-(t^2+4L+8)m_L^2>0.
\]

Equivalently,
\[
(t^2+4L+8)m_L^2+t(t^2+4L+7)m_L-(L+1)(t^2+4L+6)<0.
\]

Since \(m_L>0\), the positive root yields the uniform strict bound
\[
m_L(t)<U_L(t),
\]
where
\[
U_L(t)=
\frac{
\sqrt{t^2(t^2+4L+7)^2+4(L+1)(t^2+4L+8)(t^2+4L+6)}
-t(t^2+4L+7)
}
{2(t^2+4L+8)}.
\]

This is the determinant-to-bound extractor. It is a single theorem for every \(L\ge0\).

Define polynomials \(P_n(t)\) by
\[
P_0=1,\qquad P_1=t,\qquad P_{n+1}=tP_n+nP_{n-1}.
\]
Define \(B_n(t)\) by
\[
B_0=0,\qquad B_1=1,\qquad B_{n+1}=nB_{n-1}-tB_n.
\]

Then the same recurrence gives
\[
M_n(t)=(-1)^nP_n(t)r(t)+B_n(t).
\]

The ratio bound \(M_{L+1}<U_LM_L\) becomes
\[
\left((-1)^{L+1}P_{L+1}-(-1)^L U_LP_L\right)r
<U_LB_L-B_{L+1}.
\]

Since \(P_n(t)>0\) for \(t>0\) and \(P_{2j}(0)>0\), the coefficient has sign \((-1)^{L+1}\). Thus the bounds alternate:

For even \(L\ge0\),
\[
r(t)>
\frac{B_{L+1}(t)-U_L(t)B_L(t)}
{P_{L+1}(t)+U_L(t)P_L(t)}.
\]

For odd \(L\ge1\),
\[
r(t)<
\frac{U_L(t)B_L(t)-B_{L+1}(t)}
{P_{L+1}(t)+U_L(t)P_L(t)}.
\]

These formulas are valid at \(t=0\) by direct evaluation of the recurrences and for \(t>0\) by the positivity of \(M_L\).
\end{proof}

\begin{proposition}[APP-0033: Ramanujan antiderivative is complete Bernstein]
\phantomsection\addcontentsline{toc}{subsubsection}{Proposition~\thetheorem: APP-0033: Ramanujan antiderivative is complete Bernstein}
\label{res:ramanujan-antiderivative-complete-bernstein}
The antiderivative \(\widetilde I_R(x)=a+\int_0^\infty(1-e^{-xt})\,dt/[t(\pi^2+\log^2 t)]\) of the Ramanujan integral is a complete Bernstein function.

This resolves the source problem recorded as \textbf{APP-0033} \cite{mishra-swaminathan-ramanujan}.
\end{proposition}

\begin{proof}
The density
\[
\phi_0(t)=\frac{1}{t(\pi^2+\log^2 t)}
\]
is completely monotone on \((0,\infty)\). Consequently \(I_R\) is a Stieltjes function. Moreover, the antiderivative
\[
\widetilde I_R(x)=a+\int_0^\infty (1-e^{-xt})\phi_0(t)\,dt
\]
is a complete Bernstein function.

For \(u=\log t\),
\[
\int_0^1 e^{-au}\sin(\pi a)\,da
=\frac{\pi(1+e^{-u})}{u^2+\pi^2}.
\]
Substituting \(u=\log t\) gives
\[
\int_0^1 t^{-a}\sin(\pi a)\,da
=\frac{\pi(1+t^{-1})}{\pi^2+\log^2 t}.
\]
Therefore
\[
\phi_0(t)
=\frac{1}{t(\pi^2+\log^2 t)}
=\frac{1}{\pi(1+t)}\int_0^1 t^{-a}\sin(\pi a)\,da.
\]

For \(0<a<1\), \(t^{-a}\) is completely monotone because
\[
t^{-a}=\frac1{\Gamma(a)}\int_0^\infty e^{-ts}s^{a-1}\,ds.
\]
Also \(t\mapsto(1+t)^{-1}\) is completely monotone. Complete monotonicity is closed under products and positive mixtures, and \(\sin(\pi a)>0\) on \((0,1)\). Hence \(\phi_0\) is completely monotone.

If \(\phi\) is completely monotone and
\[
F(x)=\int_0^\infty e^{-xt}\phi(t)\,dt
\]
is finite for \(x>0\), then \(F\) is Stieltjes: write \(\phi(t)=\int_0^\infty e^{-ts}\,d\mu(s)\) and use Tonelli to obtain
\[
F(x)=\int_0^\infty\frac{d\mu(s)}{x+s}.
\]
Applying this to \(\phi_0\) proves that \(I_R\) is Stieltjes.

Finally, the source already proves that
\[
\widetilde I_R(x)=a+\int_0^\infty (1-e^{-xt})\phi_0(t)\,dt
\]
is a Bernstein function. The Levy density \(\phi_0\) is completely monotone, and it satisfies the Bernstein integrability condition
\[
\int_0^\infty (1\wedge t)\phi_0(t)\,dt<\infty.
\]
The standard complete-Bernstein criterion for Levy densities therefore gives that \(\widetilde I_R\) is a complete Bernstein function.
\end{proof}

\begin{proposition}[APP-0034: Bulboaca-Zayed gamma-quotient monotonicity solved]
\phantomsection\addcontentsline{toc}{subsubsection}{Proposition~\thetheorem: APP-0034: Bulboaca-Zayed gamma-quotient monotonicity solved}
\label{res:bulboaca-zayed-gamma-quotient-full-monotonicity}
For the continuous extension \(\widetilde F\) of \(F(x)=\log\Gamma(x+1)/(\log(x^2+6)-\log(x+6))\) on \((-1,\infty)\), prove analytically that \(\widetilde F'\) has a unique zero \(x_m\simeq1.126207061\ldots\), with \(\widetilde F'<0\) on \((-1,x_m)\) and \(\widetilde F'>0\) on \((x_m,\infty)\).

This resolves the source problem recorded as \textbf{APP-0034} \cite{bulboaca-zayed-gamma}.
\end{proposition}

\begin{proof}
\emph{Setup.}
Use the notation from the previous Bulboaca--Zayed pass:
\[
G(x)=\log\Gamma(x+1),\qquad
D(x)=\log\frac{x^2+6}{x+6},
\]
and
\[
F(x)=\frac{G(x)}{D(x)}.
\]
For \(x>1\), \(D(x)>0\), and
\[
D'(x)=p(x)=\frac{x^2+12x-6}{(x^2+6)(x+6)}.
\]
Let
\[
N_0(x)=\psi(x+1)D(x)-G(x)D'(x).
\]
The derivative-sign numerator recorded in
the corresponding result is
\[
N(x)=((x^2+6)(x+6))\,N_0(x),
\]
so \(N\) and \(N_0\) have the same sign on \(x>1\).

The source itself proves the decreasing part on \((-1,1)\), and the earlier
the preceding theorem proves \(F'(x)>0\) for
\(x\ge8\).  It remains to prove a single sign change of \(N_0\) on \((1,8]\).

Put
\[
P(x)=x^2+12x-6,\qquad Q(x)=(x^2+6)(x+6),
\]
so \(p=P/Q\), and define
\[
r(x)=\frac{\psi(x+1)}{p(x)}=\psi(x+1)\frac{Q(x)}{P(x)}.
\]
Then
\[
P(x)^2r'(x)=S(x),
\]
where
\[
S(x)=\psi'(x+1)Q(x)P(x)+\psi(x+1)K(x),
\]
and
\[
K(x)=Q'(x)P(x)-Q(x)P'(x)
=x^4+24x^3+48x^2-144x-468.
\]
Direct differentiation gives
\[
S'(x)=P(x)\{Q(x)\psi''(x+1)+2Q'(x)\psi'(x+1)\}+K'(x)\psi(x+1).
\]

For \(y=x+1\ge2\), use the standard elementary bounds
\[
\psi(y)>\log y-\frac1y,\qquad
\psi'(y)>\frac1y+\frac{1}{2y^2},\qquad
\psi''(y)>-\frac1{y^2}-\frac2{y^3}.
\]
The last inequality follows from the series
\[
\psi''(y)=-2\sum_{n=0}^{\infty}(n+y)^{-3}
\]
and an integral upper bound for the sum.

Substituting these bounds into \(S'\) yields
\[
S'(x)>
P(x)\frac{5x^4+30x^3+57x^2+12x-90}{(x+1)^3}
K'(x)\left(\log(x+1)-\frac1{x+1}\right).
\]
On \(x\ge1\), \(P(x)>0\).  The quartic numerator has value \(14\) at \(x=1\)
and positive derivative
\[
20x^3+90x^2+114x+12.
\]
Also
\[
K'(x)=4x^3+72x^2+96x-144
\]
has value \(28\) at \(x=1\) and positive derivative for \(x\ge1\), while
\(\log(x+1)-1/(x+1)>0\) for \(x\ge1\).  Hence
\[
S'(x)>0\qquad(1\le x\le8).
\]

At the endpoints,
\[
S(1)=343\left(\frac{\pi^2}{6}-1\right)-539(1-\gamma)<0,
\]
for example using \(\pi^2<987/100\) and \(\gamma<5773/10000\), which gives
\(-66003/10000<0\).  Also \(S(8)>0\), because \(K(8)=17836>0\),
\(\psi'(9)>0\), and \(\psi(9)=H_8-\gamma>0\).  Therefore \(S\) has exactly
one zero on \([1,8]\).  Since \(P^2>0\), \(r'\) has exactly one zero; \(r\)
decreases first and then increases.

Now define
\[
I(x)=D(x)r(x)-G(x).
\]
Then
\[
N_0(x)=p(x)I(x),\qquad I'(x)=D(x)r'(x).
\]
Because \(D(x)>0\) for \(x>1\), \(I\) decreases and then increases with the
same turning point as \(r\).  Moreover \(\lim_{x\downarrow1}I(x)=0\), so
\(I(x)<0\) immediately to the right of \(1\).

Finally,
\[
N_0(8)=(H_8-\gamma)\log5-\frac{11}{70}\log(40320)>0.
\]
For instance, the bounds
\[
\gamma<\frac{5773}{10000},\quad \log5>\frac{1609}{1000},
\quad \log(40320)<\frac{10605}{1000}
\]
give the lower bound
\[
\left(\frac{761}{280}-\frac{5773}{10000}\right)\frac{1609}{1000}
-\frac{11}{70}\frac{10605}{1000}
=\frac{124435951}{70000000}>0.
\]
Thus \(I\), and therefore \(N_0\) and \(N\), has exactly one zero on
\((1,8]\).  Denote it by \(\xi\).  The sign pattern is
\[
N(x)<0\quad(1<x<\xi),\qquad
N(\xi)=0,\qquad
N(x)>0\quad(\xi<x\le8).
\]

A high-precision numerical check places
\[
\xi=1.12620706105164346558\ldots,
\]
matching the source's reported value.  The numerical value is not used in the
proof.

The raw numerator has a removable zero at \(x=1\), since \(G(1)=D(1)=0\).  Its
quadratic leading coefficient is
\[
\lim_{x\downarrow1}\frac{N_0(x)}{(x-1)^2}
=
\frac{7(\pi^2/6-1)-11(1-\gamma)}{98}<0,
\]
again by \(\pi^2<987/100\) and \(\gamma<5773/10000\), which gives
\(-1347/980000<0\).  Therefore any normalization dividing out the removable
\((x-1)^2\) factor is negative at the left endpoint.

The single-crossing candidate is true:
\[
T\text{-BZ-gamma-quotient-N-single-crossing-one-eight}.
\]
Together with the source's theorem on \((-1,1)\) and the previous local
right-tail theorem for \(x\ge8\), this proves the critical-window reduction
\[
T\text{-BZ-gamma-quotient-critical-window-reduction}
\]
and then the source problem
\[
T\text{-Bulboaca-Zayed-gamma-quotient-full-monotonicity}.
\]

The finite-envelope candidate is not separately promoted: the proof above is a
clean analytic ratio-kernel argument rather than the finite rational cover
specified in that candidate.  The diagnostic obstruction-map candidate is also
not needed after the single-crossing proof.
\end{proof}

\begin{counterexample}[APP-0035: Keady self-bijection inverse-CM question is negative]
\phantomsection\addcontentsline{toc}{subsubsection}{Counterexample~\thetheorem: APP-0035: Keady self-bijection inverse-CM question is negative}
\label{res:keady-q3-self-bijection-inverse-cm-negative-answer}
Keady Question 3 has a negative answer for the self-range case: there is a completely monotone decreasing bijection \((0,\infty)\to(0,\infty)\) whose inverse is not completely monotone.

This resolves the source problem recorded as \textbf{APP-0035} \cite{app-app-0035-source}.
\end{counterexample}

\begin{proof}
There exists a completely monotone decreasing bijection
\[
f:(0,\infty)\to(0,\infty)
\]
whose inverse is not completely monotone. One explicit example is
\[
f(x)=\frac1x+100e^{-x}.
\]

For every \(n\ge0\),
\[
(-1)^n f^{(n)}(x)=\frac{n!}{x^{n+1}}+100e^{-x}>0.
\]
Thus \(f\) is strictly completely monotone. Also
\[
f'(x)=-x^{-2}-100e^{-x}<0,
\]
so \(f\) is strictly decreasing. Finally,
\[
\lim_{x\downarrow0}f(x)=\infty,
\qquad
\lim_{x\to\infty}f(x)=0.
\]
Therefore \(f\) is a decreasing bijection from \((0,\infty)\) onto \((0,\infty)\).

Let \(g=f^{-1}\) and write \(y=f(x)\). Inverse differentiation gives
\[
g'''(f(x))=\frac{3(f''(x))^2-f'''(x)f'(x)}{(f'(x))^5}.
\]
For \(f_a(x)=x^{-1}+ae^{-x}\), define
\[
N_a(x)=3(f_a''(x))^2-f_a'''(x)f_a'(x).
\]
A direct calculation gives
\[
N_a(x)=\frac6{x^6}+ae^{-x}\left(-\frac6{x^4}+\frac{12}{x^3}-\frac1{x^2}\right)+2a^2e^{-2x}.
\]
At \(a=100\) and \(x_0=1/8\),
\[
N_{100}(1/8)=1{,}572{,}864-1{,}849{,}600e^{-1/8}+20{,}000e^{-1/4}.
\]
Using \(e^{-1/8}>7/8\) and \(e^{-1/4}<1\),
\[
N_{100}(1/8)<1{,}572{,}864-1{,}849{,}600\cdot\frac78+20{,}000=-25{,}536<0.
\]
Since \(f'(1/8)<0\), the denominator \((f'(1/8))^5\) is negative. Hence
\[
g'''(f(1/8))=\frac{N_{100}(1/8)}{(f'(1/8))^5}>0.
\]
A completely monotone function must satisfy \((-1)^3g'''\ge0\), i.e. \(g'''\le0\). Therefore \(g=f^{-1}\) is not completely monotone.
\end{proof}

\begin{counterexample}[APP-0036: Baskakov even-line complete-monotonicity fails at \(r=8\)]
\phantomsection\addcontentsline{toc}{subsubsection}{Counterexample~\thetheorem: APP-0036: Baskakov even-line complete-monotonicity fails at r=8}
\label{res:not-baskakov-higher-power-even-conjecture}
Abel--Gawronski--Neuschel's even-power Baskakov complete-monotonicity conjecture is false; it fails at \(\alpha=1\), \(r=8\).

This resolves the source problem recorded as \textbf{APP-0036} \cite{abel-gawronski-neuschel}.
\end{counterexample}

\begin{proof}
Set
\[
D_m(x)=(1+x)^{2m}-x^{2m}.
\]
The roots are obtained from
\[
\frac{1+x}{x}=\omega_k,\qquad
\omega_k=e^{\pi i k/m},\qquad k=1,\ldots,2m-1.
\]
Thus
\[
\lambda_{k,m}
=\frac{1}{\omega_k-1}
=-\frac12-\frac{i}{2}\cot\frac{\pi k}{2m}.
\]
At this root,
\[
D_m'(\lambda_{k,m})
=2m\lambda_{k,m}^{2m-1}(\omega_k^{-1}-1),
\]
and direct simplification gives the real residue
\[
R_{k,m}
=\frac{1}{D_m'(\lambda_{k,m})}
=\frac{(-1)^{m+k}}{2m}
\left(2\sin\frac{\pi k}{2m}\right)^{2m-2}.
\]
Consequently the inverse-Laplace density is
\[
\rho_m(t)
=\sum_{k=1}^{2m-1}R_{k,m}e^{\lambda_{k,m}t}.
\]
Pairing conjugate roots gives
\[
\rho_m(t)
=\frac{e^{-t/2}}{2m}
\left[
2^{2m-2}
+2\sum_{k=1}^{m-1}
(-1)^{m+k}
\left(2\sin\frac{\pi k}{2m}\right)^{2m-2}
\cos\!\left(
\frac{t}{2}\cot\frac{\pi k}{2m}
\right)
\right].
\]

This proves the finite trigonometric density formula used by the diagnostic
candidate.

For \(m=2\), this gives
\[
\rho_2(t)=e^{-t/2}\left(1-\cos\frac t2\right)
=2e^{-t/2}\sin^2\frac t4\ge0,
\]
recovering the previously admitted \(r=4,\alpha=1\) seed.

For \(m=3\),
\[
\rho_3(t)
=\frac{e^{-t/2}}6
\left[
16+2\cos\left(\frac{\sqrt3\,t}{2}\right)
-18\cos\left(\frac{t}{2\sqrt3}\right)
\right].
\]
Set \(u=t/(2\sqrt3)\).  Since
\(\cos(3u)=4\cos^3u-3\cos u\),
\[
\rho_3(t)
=\frac{4}{3}e^{-t/2}(1-\cos u)^2(2+\cos u)\ge0.
\]
Thus \(r=6,\alpha=1\) is completely monotone by an explicit positive density.

For \(m=4\), write \(h_4(t)=e^{t/2}\rho_4(t)\).  The density formula gives
\[
h_4(t)
=8+2\cos\frac t2
-\frac{(2-\sqrt2)^3}{4}
\cos\left(\frac{1+\sqrt2}{2}t\right)
-\frac{(2+\sqrt2)^3}{4}
\cos\left(\frac{\sqrt2-1}{2}t\right).
\]
At \(t=10\pi\),
\[
\cos(t/2)=\cos(5\pi)=-1,
\]
and both other cosine terms equal \(-\cos(5\pi\sqrt2)\).  Since
\[
(2-\sqrt2)^3+(2+\sqrt2)^3=40,
\]
we obtain
\[
h_4(10\pi)=6+10\cos(5\pi\sqrt2).
\]
Now
\[
0<5\sqrt2-7<\frac14
\]
because \(7/5<\sqrt2<29/20\).  Hence, with
\(\delta=5\sqrt2-7\),
\[
\cos(5\pi\sqrt2)=\cos(7\pi+\pi\delta)=-\cos(\pi\delta)
<-\cos\frac\pi4=-\frac{\sqrt2}{2}<-\frac35.
\]
Therefore
\[
h_4(10\pi)<6+10\left(-\frac35\right)=0,
\]
and
\[
\rho_4(10\pi)=e^{-5\pi}h_4(10\pi)<0.
\]
Since the inverse-Laplace transform is unique for measures on
\([0,\infty)\), \(f^{[8]}_1(x)=1/((1+x)^8-x^8)\) cannot be completely monotone.
\end{proof}

\begin{counterexample}[APP-0037: Du-Wang \(h\_3\) is increasing exactly on \(\frac12\le a\le1\)]
\phantomsection\addcontentsline{toc}{subsubsection}{Counterexample~\thetheorem: APP-0037: Du-Wang h\\_3 is increasing exactly on 12 a1}
\label{res:du-wang-h3-open-problem-2-classification}
For \(0<a<2\), Du--Wang's function \(h_3\) is increasing on \((0,\infty)\) exactly for \(1/2\le a\le1\), and is not monotone for \(0<a<1/2\) or \(1<a<2\).

This resolves the source problem recorded as \textbf{APP-0037} \cite{du-wang-h3}.
\end{counterexample}

\begin{proof}
For \(t>1/2\), define
\[
S_m(t)=\sum_{n=0}^{\infty}\frac1{(n+t)^m}.
\]
The polygamma identities give
\[
\psi''(t)=-2S_3(t),\qquad \psi'''(t)=6S_4(t).
\]
Thus the desired bound
\[
-\frac{2\psi''(t)}{\psi'''(t)}\ge t-\frac12
\]
is equivalent to
\[
2S_3(t)-3\left(t-\frac12\right)S_4(t)\ge0.
\]
Write \(y=t-\frac12>0\).  Using
\[
S_m(t)=\frac1{\Gamma(m)}
\int_0^\infty \frac{x^{m-1}e^{-tx}}{1-e^{-x}}\,dx
\]
and
\[
\frac{e^{-tx}}{1-e^{-x}}
=e^{-yx}\frac1{2\sinh(x/2)},
\]
set
\[
K(x)=\frac1{2\sinh(x/2)}.
\]
Then
\[
2S_3(t)-3yS_4(t)
=
\int_0^\infty e^{-yx}x^2K(x)\,dx
-\frac y2\int_0^\infty e^{-yx}x^3K(x)\,dx.
\]
Since \(x^3K(x)\sim x^2\) as \(x\downarrow0\) and decays exponentially at
infinity, integration by parts gives
\[
y\int_0^\infty e^{-yx}x^3K(x)\,dx
=
\int_0^\infty e^{-yx}(x^3K(x))'\,dx.
\]
Therefore
\[
2S_3(t)-3yS_4(t)
=
\int_0^\infty e^{-yx}
\left(x^2K(x)-\frac12(x^3K(x))'\right)dx.
\]
Now
\[
K'(x)=-\frac12K(x)\coth(x/2),
\]
so
\[
x^2K(x)-\frac12(x^3K(x))'
=
\frac12x^2K(x)\left(\frac x2\coth(x/2)-1\right).
\]
For \(z>0\), \(z\coth z>1\).  The integrand is therefore strictly positive,
and
\[
2S_3(t)-3\left(t-\frac12\right)S_4(t)>0.
\]
This proves
\[
-\frac{2\psi''(t)}{\psi'''(t)}>t-\frac12,
\qquad t>\frac12.
\]

This proves the claim.

Let \(1/2\le a\le1\), \(u>0\), and \(t=a+u\).  Since \(a\ge1/2\),
\[
u=t-a\le t-\frac12.
\]
The ratio lemma gives
\[
t-\frac12<-\frac{2\psi''(t)}{\psi'''(t)}.
\]
Because \(\psi'''(t)>0\),
\[
2\psi''(t)+u\psi'''(t)<0.
\]
Consequently
\[
H_a'(u)=-u^2\{2\psi''(t)+u\psi'''(t)\}>0.
\]
Together with \(H_a(0+)=2\log\Gamma(a)\ge0\), this gives
\[
H_a(u)\ge0
\qquad (u>0,\ 1/2\le a\le1).
\]

This proves the claim.

For \(x>0\), \(u=x\), and \(t=x+a\),
\[
h_3'(x)=\widetilde h_3'(x+a)=\frac{h_{31}(x+a)}{x^2}
=\frac{H_a(x)}{x^2}\ge0.
\]
The inequality is strict for \(x>0\) except possibly at the endpoint
\(a=1,u=0\), which is outside the open \(x\)-interval.  Thus \(h_3\) is
increasing on \((0,\infty)\) for every
\[
\frac12\le a\le1.
\]

This proves the claim.
the corresponding result.

The earlier the preceding theorem proves that
\(h_3\) is not monotone for
\[
0<a<\frac12
\qquad\text{or}\qquad
1<a<2.
\]
The middle-window theorem above proves increasing behavior for
\[
\frac12\le a\le1.
\]
Therefore Du--Wang Open Problem 2 is classified on \(0<a<2\):
\[
h_3 \text{ is increasing on }(0,\infty)
\quad\Longleftrightarrow\quad
\frac12\le a\le1,
\]
and it is not monotone on the two outer windows.

This proves the claim.
the corresponding result true by source-problem classification.
\end{proof}

\begin{proposition}[APP-0038: Ma-Weigert derivative-sign regions form a chain]
\phantomsection\addcontentsline{toc}{subsubsection}{Proposition~\thetheorem: APP-0038: Ma-Weigert derivative-sign regions form a chain}
\label{res:ma-weigert-log-function-dk-chain-conjecture}
In the Ma--Weigert log-function class \(\mathcal F_{1,n}\), the derivative-sign regions \(D_k=\{f:L^k f\ge0\}\), with \(L=-d/dx\), form a descending chain: \(D_{k+1}\subseteq D_k\) for every \(k,n\).

This resolves the source problem recorded as \textbf{APP-0038} \cite{ma-weigert-log-functions}.
\end{proposition}

\begin{proof}
We first prove that for each \(k\ge0\),
\[
L^k f(x)=x^{-k-1}p_k(\log x)
\]
for some polynomial \(p_k\).

For \(k=0\), this is the definition with \(p_0=p\).  Suppose
\[
L^k f(x)=x^{-k-1}p_k(\log x).
\]
Then
\[
\begin{aligned}
L^{k+1}f(x)
&=-\frac{d}{dx}\left(x^{-k-1}p_k(\log x)\right)\\
&=x^{-k-2}\left((k+1)p_k(\log x)-p_k'(\log x)\right).
\end{aligned}
\]
Thus \(p_{k+1}(y)=(k+1)p_k(y)-p_k'(y)\), which is again a polynomial.

Since every polynomial in \(\log x\) grows slower than every positive power of \(x\),
\[
x^{-k-1}p_k(\log x)\to0
\qquad (x\to\infty).
\]
Therefore
\[
\lim_{x\to\infty}L^k f(x)=0
\]
for every \(k\ge0\).

Assume \(f\in D_{k+1}\).  Then
\[
L^{k+1}f(x)\ge0
\]
on \((0,\infty)\).  Put
\[
g(x)=L^k f(x).
\]
By definition of \(L\),
\[
-g'(x)=L^{k+1}f(x)\ge0.
\]
Using \(g(x)\to0\) as \(x\to\infty\), integrate from \(x\) to \(R\):
\[
g(x)-g(R)=\int_x^R L^{k+1}f(t)\,dt.
\]
Letting \(R\to\infty\) gives
\[
L^k f(x)=g(x)=\int_x^\infty L^{k+1}f(t)\,dt\ge0.
\]
Hence \(f\in D_k\).  Therefore
\[
D_{k+1}\subseteq D_k
\]
for all \(k\ge0\) and all \(n\).
\end{proof}

\begin{proposition}[APP-0046: Tao--Sawin finite Weyl \(\ell^1\) minimum is exact]
\phantomsection\addcontentsline{toc}{subsubsection}{Proposition~\thetheorem: APP-0046: Tao--Sawin finite Weyl ^1 minimum is exact}
\label{res:tao-sawin-weyl-l1-exact-minimum}
In the Weyl algebra DX-XD=1, the least l1 coefficient norm of a balanced degree (n,n) word combination equal to 1 is 2^n/n!.
\[
L_n=min{sum_w |a_w|: sum_w a_w w=1, w has n D's and n X's}=2^n/n!.
\]

This resolves the source problem recorded as \textbf{APP-0046} \cite{app-app-0046-source}.
\end{proposition}

\begin{proof}
Let \(A_1\) be the Weyl algebra over a field of characteristic \(0\), generated by \(X,D\) with
\[
DX-XD=1.
\]
Let \(W_n\) be the set of words containing exactly \(n\) copies of \(D\) and \(n\) copies of \(X\). For
\[
P=\sum_{w\in W_n}a_w w
\]
define
\[
\|P\|_1=\sum_{w\in W_n}|a_w|.
\]
Let
\[
L_n=\min\left\{\|P\|_1:\sum_{w\in W_n}a_w w=1\text{ in }A_1\right\}.
\]
Then
\[
L_n=\frac{2^n}{n!}.
\]

Set
\[
P_n=\frac1{n!}\operatorname{ad}_D^n(X^n).
\]
Since \(\operatorname{ad}_D(X)=1\), repeated derivation gives
\[
\operatorname{ad}_D^n(X^n)=n!.
\]
Also
\[
\operatorname{ad}_D^n(X^n)
=
\sum_{a=0}^n(-1)^a\binom naD^{n-a}X^nD^a.
\]
Thus \(P_n=1\) in \(A_1\), and
\[
\|P_n\|_1
=
\frac1{n!}\sum_{a=0}^n\binom na
=
\frac{2^n}{n!}.
\]
Therefore \(L_n\le2^n/n!\).

For \(q\in\mathbb C\), define \(\Phi_q\) on normal-ordered monomials by
\[
\Phi_q(X^rD^s)=
\begin{cases}
q^r r!,&r=s,\\
0,&r\ne s.
\end{cases}
\]
The functional relevant to the lower bound is
\[
\Lambda_n=\Phi_{-1/2},
\]
so
\[
\Lambda_n(X^iD^i)=(-1)^i2^{-i}i!.
\]
In particular, \(\Lambda_n(1)=1\).

For letters \(a_i\in\{X,D\}\),
\[
\Phi_q(a_1\cdots a_m)
=
\sum_\pi
\prod_{\{i,j\}\in\pi,\ i<j} C_q(a_i,a_j),
\]
where \(\pi\) runs over pairings pairing each \(X\) with a \(D\), and
\[
C_q(X,D)=q,\qquad
C_q(D,X)=q+1,\qquad
C_q(X,X)=C_q(D,D)=0.
\]

To prove this, use formal exponentials. Since
\[
e^{uD}e^{vX}=e^{uv}e^{vX}e^{uD},
\]
normal ordering gives
\[
e^{t_1a_1}\cdots e^{t_ma_m}
=
\exp\left(\sum_{\substack{i<j\\a_i=D,\ a_j=X}}t_it_j\right)
e^{UX}e^{VD},
\]
where
\[
U=\sum_{i:a_i=X}t_i,\qquad
V=\sum_{i:a_i=D}t_i.
\]
Applying \(\Phi_q\) gives
\[
\Phi_q(e^{UX}e^{VD})=e^{qUV}.
\]
Hence
\[
\Phi_q(e^{t_1a_1}\cdots e^{t_ma_m})
=
\exp\left(\sum_{i<j}C_q(a_i,a_j)t_it_j\right).
\]
Extracting the coefficient of \(t_1\cdots t_m\) gives the Wick pairing formula.

Take \(q=-1/2\). Then
\[
C_{-1/2}(X,D)=-\frac12,\qquad
C_{-1/2}(D,X)=\frac12.
\]
Every nonzero contraction has absolute value \(1/2\). If \(w\in W_n\), then pairings of the \(n\) \(X\)'s with the \(n\) \(D\)'s are in bijection with permutations, so there are exactly \(n!\) possible pairings. Therefore
\[
|\Lambda_n(w)|
=
|\Phi_{-1/2}(w)|
\le
n!\left(\frac12\right)^n
=
\frac{n!}{2^n}.
\]

Now let
\[
P=\sum_{w\in W_n}a_w w
\]
and suppose \(P=1\) in \(A_1\). Applying \(\Lambda_n\),
\[
1=|\Lambda_n(1)|
=
|\Lambda_n(P)|
\le
\sum_{w\in W_n}|a_w|\,|\Lambda_n(w)|
\le
\frac{n!}{2^n}\sum_{w\in W_n}|a_w|.
\]
Thus
\[
\|P\|_1\ge\frac{2^n}{n!}.
\]
Combining this with the upper bound proves
\[
L_n=\frac{2^n}{n!}.
\]

This solves the finite homogeneous Weyl algebra \(\ell^1\) coefficient model problem from the Tao-Sawin comment thread. It does not solve the full operator-norm problem for commutators close to the identity.
\end{proof}

\begin{counterexample}[APP-0047: BPV noncentral chi-square HCM range has no positive noncentrality]
\phantomsection\addcontentsline{toc}{subsubsection}{Counterexample~\thetheorem: APP-0047: BPV noncentral chi-square HCM range has no positive noncentrality}
\label{res:bpv-noncentral-chi-square-hcm-no-positive-noncentrality}
For every \(\mu>0\) and every positive noncentrality \(\lambda>0\), the noncentral chi-square density \(\chi_{\mu,\lambda}\) is not hyperbolically completely monotone. Equivalently, the BPV optimal HCM range for positive noncentrality is empty; if the central case \(\lambda=0\) is adjoined, the remaining HCM range is exactly the central gamma line.

This resolves the source problem recorded as \textbf{APP-0047} \cite{app-app-0047-source}.
\end{counterexample}

\begin{proof}
Let \(a=\mu/2>0\). The complete Bell generating function is
\[
\sum_{n\ge0}\frac{B_nz^n}{n!}
=\exp\left(\sum_{j\ge1}\frac{c_jz^j}{j!}\right)
=e^{-z/2}S_\mu(\lambda z).
\]
Indeed \(b_1(\mu)=1/(2\mu)\), so the \(j=1\) term of \(\log S_\mu(\lambda z)\) contributes \(\lambda z/(2\mu)\), and the exponential factor \(e^{-z/2}\) supplies the remaining \(-z/2\).

Define
\[
A_n=\frac{(-1)^nB_n}{n!}.
\]
Then
\[
\sum_{n\ge0}A_nt^n
=e^{t/2}S_\mu(-\lambda t)
=e^{t/2}\,{}_0F_1\left(;a;-\frac{\lambda t}{4}\right).
\]
Equivalently,
\[
A_n=\frac{2^{-n}}{(a)_n}L_n^{a-1}\left(\frac{\lambda}{2}\right),
\]
so the leading HCM signs are exactly the fixed-argument Laguerre signs.

Use the classical identity
\[
{}_0F_1\left(;a;-\frac{x^2}{4}\right)
=\Gamma(a)\left(\frac{x}{2}\right)^{1-a}J_{a-1}(x).
\]
Since \(a>0\), the Bessel function \(J_{a-1}\) has a positive real zero \(j_{a-1,1}>0\). For \(\lambda>0\), set
\[
t_0=\frac{j_{a-1,1}^2}{\lambda}>0.
\]
Then
\[
e^{t_0/2}{}_0F_1\left(;a;-\frac{\lambda t_0}{4}\right)=0.
\]
If every leading HCM sign were nonnegative, then all coefficients \(A_n\) would be nonnegative and \(A_0=1\), forcing
\[
\sum_{n\ge0}A_nt_0^n>0,
\]
contradicting the displayed zero. Therefore, for every \(\mu>0\) and \(\lambda>0\), there is an integer \(n\ge1\) such that
\[
(-1)^nB_n<0.
\]

The leading expansion gives, for any fixed compact \(w\)-window,
\[
\frac{1}{H_u(w)}\frac{d^n}{dw^n}H_u(w)
=B_nu^n+O(u^{n+1})
\qquad(u\downarrow0).
\]
Choosing \(u>0\) sufficiently small yields
\[
(-1)^n\frac{d^n}{dw^n}H_u(w)<0,
\]
which violates complete monotonicity of \(w\mapsto H_u(w)\).

\emph{Conclusion.}
For every \(\mu>0\) and every positive noncentrality \(\lambda>0\),
\[
\chi_{\mu,\lambda}\notin HCM.
\]
Thus the BPV optimal HCM range for positive noncentrality is empty. If one adjoins the central case \(\lambda=0\), the remaining density is a gamma density and is HCM; hence the extended closed range is exactly the central line \(\lambda=0\).

Strict APP status: candidate for the next APP slot after APP-0046 because the source explicitly asks the optimal range, the local theorem answers it, and the result is not among public APP-0001--APP-0045.
\end{proof}

\begin{counterexample}[APP-0048: BMR tau-Gauss ordinary concavity is false]
\phantomsection\addcontentsline{toc}{subsubsection}{Counterexample~\thetheorem: APP-0048: BMR tau-Gauss ordinary concavity is false}
\label{res:bmr-tau-gauss-not-globally-concave}
The Bansal--Mehrez--Raina tau-Gauss source function \(a\mapsto {}_2\phi^\tau_1(a,c-a;c;z)\) is not concave in general on \((0,c)\). In the subcase \(c=2\), \(\tau=1\), and \(z=1-e^{-4}\), one has \(F''(a,z)>0\) for all sufficiently small \(a>0\).

This resolves the source problem recorded as \textbf{APP-0048} \cite{app-app-0048-source}.
\end{counterexample}

\begin{proof}
For \(0<z<1\), define probability weights
\[
w_k(a,z)=\frac{A_k(a)z^k/k!}{F(a,z)}.
\]
Let
\[
g_k(a)=\partial_a\log A_k(a),
\qquad
h_k(a)=\partial_a^2\log A_k(a).
\]
Then direct differentiation of the log-sum gives
\[
\partial_a^2\log F(a,z)
=
\sum_{k\ge0}w_k(a,z)h_k(a)
+\operatorname{Var}_{w(a,z)}(g_k(a)).
\]
Thus global log-concavity is equivalent to the variance-domination inequality
\[
\operatorname{Var}_{w(a,z)}(g_k(a))
\le
\sum_{k\ge0}w_k(a,z)(-h_k(a))
\]
for all \(0<a<c\), \(c,\tau>0\), and \(0<z<1\). This is the sharp remaining gate; coefficientwise log-concavity alone is insufficient.

Take the classical subcase
\[
c=2,\qquad \tau=1,\qquad z=1-e^{-4}.
\]
Then
\[
F(a,z)=
\sum_{k=0}^{\infty}
\frac{(a)_k(2-a)_k}{(2)_k}\frac{z^k}{k!}.
\]
For \(k\ge1\), as \(a\downarrow0\),
\[
\frac{(a)_k(2-a)_k}{(2)_k\,k!}
=
\frac{a}{k}
+a^2\left(
\frac1k-\frac1{k^2}-\frac1{k(k+1)}
\right)+O(a^3),
\]
with locally uniform summability for fixed \(0<z<1\). Hence
\[
F(a,z)=1+aL(z)+a^2M(z)+O(a^3),
\]
where
\[
L(z)=\sum_{k\ge1}\frac{z^k}{k}=-\log(1-z)
\]
and
\[
M(z)=
\sum_{k\ge1}z^k
\left(
\frac1k-\frac1{k^2}-\frac1{k(k+1)}
\right)
=
\frac{L(z)}{z}-\operatorname{Li}_2(z)-1.
\]
For \(z=1-e^{-4}\), \(L(z)=4\), \(L(z)/z>4\), and \(\operatorname{Li}_2(z)<\pi^2/6<2\). Therefore
\[
M(z)>4-2-1=1>0.
\]
It follows that
\[
F''(a,z)\to2M(z)>0
\qquad(a\downarrow0),
\]
so \(F''(a,z)>0\) for all sufficiently small \(a>0\). Therefore the source function is not concave in \(a\) in general.
\end{proof}

\begin{proposition}[APP-0049: Baricz arithmetic/arithmetic zero-balanced hypergeometric threshold]
\phantomsection\addcontentsline{toc}{subsubsection}{Proposition~\thetheorem: APP-0049: Baricz arithmetic/arithmetic zero-balanced hypergeometric threshold}
\label{res:baricz-aa-zero-balanced-concavity-threshold}
For \(c>0\), define \(F_a(r)={}_2F_1(a,c-a;c;r)\) on \(0<a<c\), \(0<r<1\). The uniform arithmetic/arithmetic inequality \((F_{a_1}(r)+F_{a_2}(r))/2\le F_{(a_1+a_2)/2}(r)\) for all \(a_1,a_2\in(0,c)\) and all \(r\in(0,1)\) holds if and only if \(0<c\le1\).

This resolves the source problem recorded as \textbf{APP-0049} \cite{app-app-0049-source}.
\end{proposition}

\begin{proof}
For \(c>0\), set
\[
F_a(r)={}_2F_1(a,c-a;c;r),\qquad 0<a<c,\quad 0<r<1.
\]
The Baricz arithmetic/arithmetic source slice asks when
\[
\frac{F_{a_1}(r)+F_{a_2}(r)}2\le F_{(a_1+a_2)/2}(r)
\]
holds for all \(a_1,a_2\in(0,c)\) and all \(r\in(0,1)\).

The answer is exactly \(0<c\le1\).

The existing local node the corresponding result imports Baricz's theorem:
\[
\sqrt{F_{a_1}(r)F_{a_2}(r)}
\le
\frac{F_{a_1}(r)+F_{a_2}(r)}2
\le
F_{(a_1+a_2)/2}(r)
\]
for \(0<c\le1\), \(a_1,a_2\in(0,c)\), and \(0<r<1\). This proves the A/A inequality on that range.

Fix \(c>1\). For \(0<r<1\),
\[
F_a(r)=\sum_{k=0}^{\infty}
\frac{(a)_k(c-a)_k}{(c)_k}\frac{r^k}{k!}.
\]
For \(k\ge1\), as \(a\to0\),
\[
(a)_k=a(k-1)!\left(1+aH_{k-1}+O(a^2)\right)
\]
and
\[
\frac{(c-a)_k}{(c)_k}
=
1-a\sum_{j=0}^{k-1}\frac1{c+j}+O(a^2).
\]
Hence
\[
\frac{(a)_k(c-a)_k}{(c)_k k!}
=
\frac ak+\frac{a^2}{k}
\left(
H_{k-1}-\sum_{j=0}^{k-1}\frac1{c+j}
\right)
 +O(a^3).
\]
For fixed \(r<1\), normal convergence of the differentiated hypergeometric series justifies termwise expansion, so
\[
F_a(r)=1+aL(r)+a^2M_c(r)+O(a^3),
\]
where
\[
L(r)=\sum_{k\ge1}\frac{r^k}{k}=-\log(1-r)
\]
and
\[
M_c(r)=
\sum_{k\ge1}
\left(
H_{k-1}-\sum_{j=0}^{k-1}\frac1{c+j}
\right)\frac{r^k}{k}.
\]
The bracket is
\[
B_{c,k}
=
\psi(k)+\gamma-\psi(c+k)+\psi(c),
\]
so
\[
B_{c,k}\to \gamma+\psi(c).
\]
Because \(c>1\), \(\psi(c)>\psi(1)=-\gamma\), and therefore \(\gamma+\psi(c)>0\). Thus \(B_{c,k}\ge\eta>0\) for all sufficiently large \(k\), and
\[
M_c(r)\to+\infty\qquad (r\uparrow1).
\]
Choose \(r_0\in(0,1)\) with \(M_c(r_0)>0\). Then
\[
\partial_a^2F_a(r_0)\to 2M_c(r_0)>0
\qquad (a\downarrow0).
\]
For some sufficiently small interior \(a_0\in(0,c)\), \(\partial_a^2F_{a_0}(r_0)>0\). By continuity, \(\partial_a^2F_a(r_0)>0\) on a small interval around \(a_0\). Taking \(h>0\) small enough that \(a_0\pm h\in(0,c)\), Taylor's theorem gives
\[
\frac{F_{a_0-h}(r_0)+F_{a_0+h}(r_0)}2>F_{a_0}(r_0).
\]
Since \(a_0=((a_0-h)+(a_0+h))/2\), the A/A inequality fails.

Therefore the source A/A inequality holds uniformly if and only if \(0<c\le1\).
\end{proof}

\begin{counterexample}[APP-0050: Stolarsky shifted power means with \(p>1\) are not Bernstein]
\phantomsection\addcontentsline{toc}{subsubsection}{Counterexample~\thetheorem: APP-0050: Stolarsky shifted power means with p>1 are not Bernstein}
\label{res:not-stolarsky-power-mean-pgt1-bernstein-all-shifts}
not(For some \(p>1\), the shifted power mean \(x\mapsto H_p(x+a,x+b)\), where \(H_p(a,b)=((a^p+b^p)/2)^{1/p}\), is a Bernstein function for every positive pair of shifts \(a,b\).)

This resolves the source problem recorded as \textbf{APP-0050} \cite{app-app-0050-source}.
\end{counterexample}

\begin{proof}
Let
\[
c=\frac{a+b}{2},\qquad d=\frac{a-b}{2}.
\]
Since \(a,b>0\) and \(a\ne b\), we have \(c>|d|>0\). Put \(y=x+c\). Then on the natural half-line \(y>|d|\),
\[
F(x)=G(y),
\qquad
G(y)=\left(\frac{(y+d)^p+(y-d)^p}{2}\right)^{1/p}.
\]
Translation does not affect the sign pattern of derivatives, so it is enough to show that \(G\) is not Bernstein.

For \(|t|<1\), define
\[
\phi_p(t)=H_p(1+t,1-t)
=2^{-1/p}\left((1+t)^p+(1-t)^p\right)^{1/p}.
\]
By homogeneity,
\[
G(y)=y\phi_p(d/y).
\]
Let
\[
A=1+t,\qquad B=1-t,\qquad N=A^p+B^p.
\]
Then \(\phi_p(t)=2^{-1/p}N^{1/p}\). Differentiating twice gives
\[
\phi_p''(t)
=2^{-1/p}(p-1)N^{1/p-2}
\left(N(A^{p-2}+B^{p-2})-(A^{p-1}-B^{p-1})^2\right).
\]
The bracket simplifies to
\[
A^{p-2}B^{p-2}(A+B)^2,
\]
because
\[
(A^p+B^p)(A^{p-2}+B^{p-2})-(A^{p-1}-B^{p-1})^2
=A^{p-2}B^{p-2}(A^2+2AB+B^2).
\]
For \(p>1\) and \(|t|<1\), all factors are positive, so
\[
\phi_p''(t)>0.
\]
Since \(t=d/y\),
\[
G'(y)=\phi_p(t)-t\phi_p'(t)
\]
and
\[
G''(y)=\frac{t^2}{y}\phi_p''(t)
=\frac{d^2}{y^3}\phi_p''(d/y)>0
\]
for every \(y>|d|\).

For \(|u|<1\), define
\[
A(u)=\frac{(1+u)^p+(1-u)^p}{2}.
\]
The odd terms cancel in the binomial expansion, giving
\[
A(u)=1+\frac{p(p-1)}{2}u^2+O(u^4)
\qquad (u\to0).
\]
Since \(z\mapsto z^{1/p}\) is analytic near \(z=1\),
\[
A(u)^{1/p}
=1+\frac1p\cdot\frac{p(p-1)}{2}u^2+O(u^4)
=1+\frac{p-1}{2}u^2+O(u^4).
\]
With \(u=d/y\), this yields
\[
G(y)
=yA(d/y)^{1/p}
=y+\frac{(p-1)d^2}{2y}+O(y^{-3})
\qquad (y\to\infty).
\]
Differentiating the same analytic expansion gives
\[
G'(y)
=1-\frac{(p-1)d^2}{2y^2}+O(y^{-4}),
\]
and hence
\[
G''(y)
=\frac{(p-1)d^2}{y^3}+O(y^{-5}).
\]
Because \(p>1\) and \(d\ne0\), the leading coefficient is positive. This recovers the tail obstruction from the exact curvature certificate above.

If \(G\) were a Bernstein function, then \(G'\) would be completely monotone. In particular,
\[
(G')'(y)=G''(y)\le0
\]
throughout the interval. This contradicts the positivity of \(G''(y)\) proved above.

Thus \(G\), and equivalently \(F\), is not a Bernstein function.
\end{proof}

\section{Conclusion}

The theory developed here is organized around a single principle: positivity
should be transported only through mechanisms that preserve the relevant
analytic structure.  Zeta-law compression turns Dirichlet sums into
probabilistic identities, special-function curvature is tested through
Laplace kernels, and application-level claims are resolved by moving between
Stieltjes, Bernstein, determinant, ratio, and endpoint-obstruction forms.

This perspective gives the paper its hierarchy.  Definitions fix the common
objects, lemmas and propositions isolate reusable transformations, theorems
state the main transport mechanisms, corollaries and counterexamples record
their consequences, and the application ledger points each source problem to
the result that resolves it.  The result is a connected theory of positivity
transport rather than a list of unrelated solved examples.
\begin{thebibliography}{24}

\bibitem{nantomah}
K. Nantomah,
\emph{Open Problem on Riemann Zeta Function},
ResearchGate problem note, October 2024.
\url{https://www.researchgate.net/publication/384676538_Open_Problem_on_Riemann_Zeta_Function}.

\bibitem{sroysang}
B. Sroysang,
\emph{Two Inequalities for the Riemann Zeta Functions},
Mathematica Aeterna 3, no. 1 (2013), 21--24.
\url{https://arastirmax.com/en/system/files/dergiler/135290/makaleler/3/1/arastirmax-two-inequalities-riemann-zeta-functions.pdf}.

\bibitem{alzer-kwong}
H. Alzer and M. K. Kwong,
\emph{On the concavity and convexity of \(1/\zeta\)},
International Journal of Number Theory 21, no. 8 (2025), 1825--1835.
\url{https://doi.org/10.1142/S1793042125500897}.

\bibitem{qi-lim-nantomah-polygamma}
F. Qi, D. Lim, and K. Nantomah,
\emph{Monotonicity and positivity of several functions involving ratios and products of polygamma functions},
Journal of Inequalities and Applications 2025, article 5.
\url{https://doi.org/10.1186/s13660-024-03245-8}.

\bibitem{bulboaca-zayed-gamma}
T. Bulboaca and H. M. Zayed,
\emph{Monotonic nature of the Gamma function},
Journal of Inequalities and Applications 2026, article 27.
\url{https://doi.org/10.1186/s13660-025-03425-0}.

\bibitem{kim-song-tail}
D. Kim and K. Song,
\emph{The inverses of tails of the Riemann zeta function},
Journal of Inequalities and Applications 2018, article 157.
\url{https://doi.org/10.1186/s13660-018-1743-6}.

\bibitem{pan-wu-tail}
Z. Pan and Z. Wu,
\emph{The inverse of tails of Riemann zeta function, Hurwitz zeta function and Dirichlet L-function},
AIMS Mathematics 9, no. 6 (2024), 16564--16585.
\url{https://doi.org/10.3934/math.2024803}.


\bibitem{yin-zhang-k-beta}
L. Yin and J. Zhang,
\emph{On some properties of special functions involving \(k\)-gamma and \(k\)-digamma functions},
arXiv:2502.15852, 2025.
\url{https://arxiv.org/abs/2502.15852}.


\bibitem{baricz-turan-thesis}
A. Baricz,
\emph{Turan type inequalities for some special functions},
Ph.D. thesis, University of Debrecen, 2008.
\url{https://dea.lib.unideb.hu/bitstreams/bdd4469f-1b1e-4996-9b90-73f5e1a6b0e8/download}.

\bibitem{yin-pk-digamma}
L. Yin,
\emph{Complete monotonicity of a function involving the \((p,k)\)-digamma function},
International Journal of Open Problems in Computer Mathematics 11(2) (2018), 103--108.
\url{https://www.ijopcm.org/Vol/2018/2.9.pdf}.

\bibitem{ma-weigert-log-functions}
R. Ma and J. Weigert,
\emph{Complete monotonicity of log-functions},
arXiv:2505.04225.
\url{https://arxiv.org/abs/2505.04225}.

\bibitem{qi-agarwal-gamma-ratios}
F. Qi and R. P. Agarwal,
\emph{On complete monotonicity for several classes of functions related to
ratios of gamma functions},
Journal of Inequalities and Applications 2019, article 36.
\url{https://doi.org/10.1186/s13660-019-1976-z}.


\bibitem{mishra-swaminathan-ramanujan}
D. Mishra and A. Swaminathan,
\newblock Inequalities involving a Ramanujan integral,
\newblock arXiv:2511.07443.
\url{https://arxiv.org/abs/2511.07443}.

\bibitem{sibisi-prabhakar}
N. K. Sibisi,
\newblock A probabilistic perspective on Feller, Pollard and the complete
  monotonicity of the Mittag-Leffler function,
\newblock arXiv:2301.01466.
\url{https://arxiv.org/abs/2301.01466}.

\bibitem{from-mills}
S. G. From,
\newblock Some new upper and lower bounds for the Mills ratio,
\newblock J. Math. Anal. Appl. 486 (2020), 123872.
\url{https://doi.org/10.1016/j.jmaa.2020.123872}.

\bibitem{baricz-turan}
A. Baricz,
\newblock Turan type inequalities for hypergeometric functions,
\newblock Proc. Amer. Math. Soc. 136 (2008), 3223--3229.
\url{https://doi.org/10.1090/S0002-9939-08-09353-2}.

\bibitem{qi-guo-2004}
F. Qi and B.-N. Guo,
\newblock Complete monotonicities of functions involving gamma and digamma
  functions,
\newblock RGMIA Res. Rep. Coll. 7 (2004), Article 8.
\url{https://vuir.vu.edu.au/18037/}.

\bibitem{wakrim-w}
M. Wakrim,
\newblock The W-operator: a Volterra fractional time operator with
  non-Bernstein symbol,
\newblock arXiv:2601.02876.
\url{https://arxiv.org/abs/2601.02876}.

\bibitem{qi-hlambda}
F. Qi,
\newblock Completely monotonic degree of a function involving trigamma and
  tetragamma functions,
\newblock AIMS Math. 5 (2020), 3391--3407.
\url{https://doi.org/10.3934/math.2020219}.

\bibitem{szabo-cm}
V. E. S. Szabo,
\newblock Completely monotone functions in general and some applications,
\newblock arXiv:2411.17670.
\url{https://arxiv.org/abs/2411.17670}.

\bibitem{chen-choi-gurland}
C.-P. Chen and J. Choi,
\newblock Completely monotonic functions related to Gurland's ratio for the
  gamma function,
\newblock Math. Inequal. Appl. 20 (2017), 651--659.
\url{https://doi.org/10.7153/mia-20-43}.

\bibitem{sokal-gs}
A. D. Sokal,
\newblock Real-variables characterization of generalized Stieltjes functions,
\newblock Expo. Math. 28 (2010), 179--185.
\url{https://doi.org/10.1016/j.exmath.2009.06.004}.

\bibitem{simon-moment}
T. Simon,
\newblock Moment problems related to Bernstein functions,
\newblock Ann. Fac. Sci. Toulouse Math. 29 (2020), 577--594.
\url{https://doi.org/10.5802/afst.1640}.

\bibitem{du-wang-h3}
P. Du and G. Wang,
\newblock Monotonicity, convexity, and inequalities for functions involving
  gamma function,
\newblock arXiv:2205.12530.
\url{https://arxiv.org/abs/2205.12530}.

\bibitem{abel-gawronski-neuschel}
U. Abel, W. Gawronski, and T. Neuschel,
\newblock Complete monotonicity and zeros of sums of squared Baskakov
  functions,
\newblock Appl. Math. Comput. 258 (2015), 130--137.
\url{https://arxiv.org/abs/1411.7945}.

\bibitem{apostol}
T. M. Apostol,
\emph{Introduction to Analytic Number Theory},
Undergraduate Texts in Mathematics, Springer, 1976.
\url{https://doi.org/10.1007/978-1-4757-5579-4}.

\bibitem{davenport}
H. Davenport,
\emph{Multiplicative Number Theory},
3rd ed., revised by H. L. Montgomery, Graduate Texts in Mathematics, vol. 74, Springer, 2000.
\url{https://doi.org/10.1007/978-1-4757-5927-3}.

\bibitem{titchmarsh}
E. C. Titchmarsh,
\emph{The Theory of the Riemann Zeta-function},
2nd ed., revised by D. R. Heath-Brown, Oxford University Press, 1986.
\url{https://global.oup.com/academic/product/the-theory-of-the-riemann-zeta-function-9780198533696}.

\bibitem{cover-thomas}
T. M. Cover and J. A. Thomas,
\emph{Elements of Information Theory},
2nd ed., Wiley-Interscience, 2006.
\url{https://doi.org/10.1002/047174882X}.

\bibitem{app-app-0001-source}
Horst Alzer and Man Kam Kwong, "On the concavity and convexity of \(1/\zeta\)", International Journal of Number Theory, Vol. 21, No. 8 (2025), 1825-1835. DOI: [DOI](https://doi.org/10.1142/S1793042125500897).
\url{https://doi.org/10.1142/S1793042125500897}.

\bibitem{app-app-0004-source}
Feng Qi, Dongkyu Lim, and Kwara Nantomah, "Monotonicity and positivity of several functions involving ratios and products of polygamma functions", Journal of Inequalities and Applications 2025, article 5. DOI: [DOI](https://doi.org/10.1186/s13660-024-03245-8).
\url{https://doi.org/10.1186/s13660-024-03245-8}.

\bibitem{app-app-0009-source}
Teodor Bulboaca and Hanaa M. Zayed, "Monotonic nature of the Gamma function", Journal of Inequalities and Applications 2026, article 27. DOI: [DOI](https://doi.org/10.1186/s13660-025-03425-0).
\url{https://doi.org/10.1186/s13660-025-03425-0}.

\bibitem{app-app-0035-source}
G. Keady, N. Khajohnsaksumeth, and B. Wiwatanapataphee, On functions and inverses, both positive, decreasing and convex: and Stieltjes functions, Cogent Mathematics \& Statistics 5:1 (2018), 1477543..
\url{https://doi.org/10.1080/25742558.2018.1477543}.

\bibitem{app-app-0040-source}
Zhen-Hang Yang and Jing-Feng Tian, Convexity of ratios of the modified Bessel functions of the first kind with applications, Revista Matematica Complutense, 2022..
\url{https://doi.org/10.1007/s13163-022-00439-w}.

\bibitem{app-app-0041-source}
Feng Qi, "Completely monotonic degree of a function involving trigamma and tetragamma functions", AIMS Mathematics 5 (2020), 3391-3407. DOI: [DOI](https://doi.org/10.3934/math.2020219).
\url{https://doi.org/10.3934/math.2020219}.

\bibitem{app-app-0042-source}
Mihaly Baricz, Pietro Ponnusamy, and Matti Vuorinen, Functional inequalities for modified Bessel functions, arXiv:1009.4814, DOI: [DOI](https://doi.org/10.1016/j.exmath.2011.07.001).
\url{https://doi.org/10.1016/j.exmath.2011.07.001}.

\bibitem{app-app-0046-source}
Terence Tao, "Commutators close to the identity," blog post and comment thread, with Will Sawin's \(\ell^1\)/LP framing. https://terrytao.wordpress.com/2018/04/11/commutators-close-to-the-identity/.
\url{https://terrytao.wordpress.com/2018/04/11/commutators-close-to-the-identity/}.

\bibitem{app-app-0047-source}
Baricz--Prabhu--Singh--Vijesh, Infinitely divisible modified Bessel distributions, Pacific Journal of Mathematics 343 (2026), 261--313. https://doi.org/10.2140/pjm.2026.343.261.
\url{https://doi.org/10.2140/pjm.2026.343.261}.

\bibitem{app-app-0048-source}
Deepak Bansal, Khaled Mehrez, and Ravinder Krishna Raina, Certain functional inequalities for the tau-hypergeometric functions, Journal of Inequalities and Special Functions 12(3), 33--43, 2021. https://www.ilirias.com/jiasf/repository/docs/JIASF12-3-4.pdf.
\url{https://www.ilirias.com/jiasf/repository/docs/JIASF12-3-4.pdf}.

\bibitem{app-app-0049-source}
Anderson--Vuorinen--Zhang, Topics in special functions III, arXiv:1209.1696, quoting Baricz's bivariate-means open problem..
\url{https://arxiv.org/abs/1209.1696}.

\bibitem{app-app-0050-source}
Adam Bessenyei, On complete monotonicity of some functions related to means, Mathematical Inequalities and Applications 16 (2013), 233--239. https://files.ele-math.com/articles/mia-16-17.pdf.
\url{https://files.ele-math.com/articles/mia-16-17.pdf}.

\end{thebibliography}

\end{document}
